% Auto-generated from clustering/cards/methodology_cards.json by paper/render_cards.py
% Do not edit by hand — re-run render_cards.py.

\begin{tcolorbox}[colback=bgskill, colframe=cblue!70, title={\textbf{Audit and Relax Implicit Assumption} \textcolor{gray}{\small\texttt{assumption\_audit\_and\_relax}} \hfill confidence: \textbf{high}}, breakable, before skip=4pt, after skip=4pt]
\textcolor{gray}{\small Oral $n=122$ (meta 73/122) | HC $n=20$ (meta 16/20) | Reject $n=108$ (meta 108/108)}

\textbf{Sample note.} \textit{All sample-size thresholds satisfied (n\_oral=122, n\_reject=108, HC=20 with 16/20 meta coverage); Oral meta coverage at 60\% is the lowest stratum and the main caveat for absolute strength of Oral-side claims, but reject coverage at 100\% gives high confidence in failure-mode patterns.}

\textbf{Step-by-Step:}
\begin{enumerate}[leftmargin=14pt,topsep=2pt,itemsep=1pt]
  \item Locate the implicit assumption: trace the existing method M and identify a condition A (often unstated) that its results depend on.
  \item Prove A is binding in the user's regime via a separation result (impossibility, lower bound, counterexample) --- not just empirical underperformance.
  \item Substitute A with a weaker, structurally richer, or reformulated A'; name A' as a precise mathematical condition.
  \item Re-derive the guarantee under A' (identifiability theorem, regret bound, sample complexity, privacy budget) --- the relaxation IS the theorem statement.
  \item Build a minimal instantiation (one architectural change, one regularizer, one substitution) so gains attribute unambiguously to the audited assumption.
  \item Evaluate in the regime where A breaks sharply (small-sample, distribution shift, adversarial input), not benchmarks where A approximately holds.
\end{enumerate}
\textbf{Success conditions} (Oral):
\begin{itemize}[leftmargin=12pt,topsep=2pt,itemsep=2pt]
  \item \textbf{Name the implicit assumption explicitly and prove it is the binding constraint, not a peripheral factor --- typically via a separation result (impossibility, lower bound, or counterexample) that shows what is gained or lost when A is held vs. relaxed to A'.} \textcolor{gray}{\footnotesize evidence: \texttt{ICLR\_2022\_0094}, \texttt{ICML\_2023\_0551}, \texttt{ICLR\_2025\_1631}, \texttt{ICML\_2024\_1086}} \\ LP-FT proves OOD error is controlled by head initialization in a 2-layer linear model before claiming the remedy works; the DP continual-observation paper proves a polynomial T\textasciicircum{}(1/3) gap (vs. the assumed log T) by constructing adversarial instances; the CoT-on-parity paper proves an exponential separation between with/without teacher forcing. In each case the audit is anchored by a formal separation rather than empirical contrast alone.
  \item \textbf{Demonstrate the relaxation by re-deriving guarantees (identifiability, regret, sample complexity, privacy budget) under the weaker A', not by patching the existing algorithm and benchmarking it.} \textcolor{gray}{\footnotesize evidence: \texttt{NeurIPS\_2023\_0587}, \texttt{NeurIPS\_2023\_0640}, \texttt{ICML\_2023\_0417}, \texttt{ICLR\_2025\_1648}} \\ Additive Decoders replace distributional/temporal auxiliary supervision with a single architectural constraint and then prove block-wise identifiability under additivity alone. Generalizing nonlinear ICA replaces structural sparsity with partial independence and re-proves identifiability via geometric separability. Bayesian-frequentist paper relaxes the 'priors require true Bayesian commitment' assumption and re-derives best-of-both-worlds regret bounds. The relaxation is the theorem statement, not just the method.
  \item \textbf{Show that A' creates a new design variable or capability that A foreclosed --- the relaxation pays off because the system can now exploit something previously invisible (a signal, a degree of freedom, a regime).} \textcolor{gray}{\footnotesize evidence: \texttt{ICLR\_2025\_1525}, \texttt{ICML\_2023\_0476}, \texttt{NeurIPS\_2023\_0708}, \texttt{ICML\_2024\_1133}} \\ Modeling instantaneous dependencies (previously eliminated as a nuisance) turns them into an identification source; analyzing intervention-induced support changes (rather than only conditional independence) gives identifiability invisible to covariance methods; treating DP's parallel composition as exploitable enables one-run privacy auditing; recharacterizing classifier-free guidance as signal amplification reveals it scales adversarial perturbations. The relaxed assumption converts a hidden constraint into a usable lever.
  \item \textbf{Pair the audit with an instantiation that is minimal --- one architectural change, one regularizer, one substitution --- so the gain is unambiguously attributable to the relaxed assumption, not to compounded engineering.} \textcolor{gray}{\footnotesize evidence: \texttt{ICLR\_2023\_0333}, \texttt{ICML\_2023\_0445}, \texttt{ICLR\_2022\_0094}, \texttt{ICML\_2023\_0549}} \\ Replay-ratio resets are a single periodic reinitialization. PAG paper directly trains for input-gradient alignment as the only objective change. LP-FT is two phases of standard optimization. Dormant neuron paper adds one reinitialization rule. Oral papers consistently isolate the relaxed assumption with a minimal intervention so reviewers cannot attribute gains to confounds.
  \item \textbf{Specify the operating regime where A fails sharply --- small sample, distribution shift, adversarial input, multi-pass, sparse supervision --- and evaluate there rather than on benchmarks where A is approximately satisfied.} \textcolor{gray}{\footnotesize evidence: \texttt{ICML\_2024\_1048}, \texttt{ICML\_2024\_1136}, \texttt{NeurIPS\_2024\_1289}, \texttt{ICML\_2024\_1144}} \\ Long-term causal effects are measured precisely in the regime where short-term experiments are the only available signal. Random-forest privacy attack works in the worst-case categorical regime where rules are invertible. ViP shows DP-SGD's parameter-count dependence dominates only when DP fine-tuning is the bottleneck. Each paper makes the regime where A breaks the focus of evaluation, not an afterthought.
\end{itemize}
\textbf{Failure modes} (Reject):
\begin{itemize}[leftmargin=12pt,topsep=2pt,itemsep=2pt]
  \item \textbf{Asserting an assumption is being relaxed without demonstrating it was actually binding --- the proposed change works empirically but the paper never shows that the identified assumption (rather than capacity, data quality, or regularization) drove the prior failure.} \textcolor{gray}{\footnotesize evidence: \texttt{ICLR\_2024\_0767}, \texttt{ICLR\_2025\_1644}, \texttt{ICLR\_2025\_1566}, \texttt{ICLR\_2025\_1516}} \\ AdaFlood reframes a uniform constant as per-instance adaptive but reviewers note ablations don't isolate this from capacity/data confounds. VCReg borrows VICReg's variance-covariance constraint to supervised learning without showing the supposedly suppressed structural property is the bottleneck. Reviewers explicitly call out the 'failure mode is invisible at benchmark level' problem in these abstracts and meta-reviews.
  \item \textbf{Identifying the right hidden assumption but offering only an empirical fix --- no theorem characterizes when the relaxation succeeds, so the contribution reduces to a heuristic with one new dataset.} \textcolor{gray}{\footnotesize evidence: \texttt{ICLR\_2023\_0388}, \texttt{ICLR\_2024\_0788}, \texttt{ICLR\_2024\_0917}, \texttt{ICLR\_2025\_1500}} \\ The fairness-under-covariate-shift paper relaxes MAR-style assumptions but reviewers found theorems unclear and baselines missing. Air-quality causal estimation has the right premise (continuous-treatment shift) but is methodologically incremental on top of existing semiparametric work. Edge-dependency hierarchy paper proposes the right framing but lacks comparison to canonical random-graph models. Empirical patches don't substitute for re-derived guarantees.
  \item \textbf{Combining the relaxation with multiple other novelties at once, so reviewers cannot attribute gains to the audited assumption --- papers stack components and skip ablations isolating the targeted assumption.} \textcolor{gray}{\footnotesize evidence: \texttt{ICLR\_2024\_0776}, \texttt{ICLR\_2025\_1400}, \texttt{ICLR\_2025\_1597}, \texttt{ICML\_2025\_1827}} \\ Sparse link prediction reframes as ranking AND adds attribute features AND adds hierarchical partitioning --- reviewers can't tell which addresses the imbalance assumption. Sheaf-theoretic recommender stacks topological machinery on standard pipelines. Multimodal peptide design unifies multiple targets but lacks the targeted ablations to attribute the gain to the assumption-relaxation. The card is a tool for surgical change, not bundled novelty.
  \item \textbf{Treating the relaxation as a 'natural extension' that inherits the prior framework's validity automatically, without verifying that the inherited theory actually carries over --- most visible when porting techniques across communities.} \textcolor{gray}{\footnotesize evidence: \texttt{ICLR\_2024\_0919}, \texttt{ICLR\_2025\_1481}, \texttt{ICML\_2025\_1745}, \texttt{NeurIPS\_2024\_1300}} \\ Online fractional knapsack paper assumes prediction-augmented analysis carries over without sufficient motivation for its prediction model. High-dim instrument paper extends partial-identification to representations but reviewers wanted concrete baselines to verify validity transfers. Light-as-deception attack ports adversarial framing to relighting without rigorously establishing the new attack surface. Ported assumptions need re-verification, not inheritance.
  \item \textbf{Auditing an assumption that is genuinely violated, but in a regime where the violation is small enough that the proposed fix yields only marginal improvements --- the field rejects on significance, not on logic.} \textcolor{gray}{\footnotesize evidence: \texttt{ICLR\_2023\_0381}, \texttt{ICML\_2025\_1657}, \texttt{ICLR\_2025\_1343}, \texttt{NeurIPS\_2024\_1302}} \\ Bank-loan paper combines unbiased decisions with optimistic exploration but gains over baseline are statistically marginal. DPOT-L0 backdoor attack design is principled but does not show practically meaningful risk over existing methods on broad evaluations. Reviewers consistently say 'right idea, gains too small to matter' --- proving the assumption mattered in theory does not show it matters in practice.
  \item \textbf{Mistaking notational or framing differences for assumption violations --- papers claim a 'new perspective' but the underlying assumption is unchanged or the new model is mathematically equivalent to existing ones in disguise.} \textcolor{gray}{\footnotesize evidence: \texttt{ICLR\_2025\_1385}, \texttt{ICLR\_2025\_1548}, \texttt{ICLR\_2024\_0923}, \texttt{ICML\_2025\_1718}} \\ Counterfactual learning under rank preservation introduces ordinal invariance but reviewers note the 'weaker' assumption is in fact strong and the formalism overlaps confusingly with existing SCM machinery. RFM-vs-linear paper relaxes one assumption from prior equivalence results but reviewers find contributions incremental relative to Ba et al. and Mousavi-Hosseini et al. PNL maximal correlation paper reframes optimization but the underlying identifiability is from existing PNL theory.
\end{itemize}
\textbf{Oral vs Reject gap.} Oral papers state the audited assumption as a precise mathematical condition (a measurability property, an independence structure, a regularity assumption, an equilibrium constraint) and then prove a separation --- an impossibility result, a tight bound, or an identifiability theorem --- that holds under the original assumption A but breaks under A'. The relaxation is the theorem statement; the algorithm is one minimal instantiation chosen so that gains are attributable. Reject papers describe the audited assumption in prose, replace it with a 'natural extension' or 'principled relaxation', and validate empirically --- but skip the step of proving that the assumption (rather than capacity, regularization, or data quality) was the binding constraint. Reviewers consistently demand ablations isolating the assumption-relaxation from confounds, and demand a regime where the violation is sharp rather than incidental. The visible signal in meta-reviews: Oral reviews mention 'derives identifiability conditions', 'proves separation', 'characterizes when'; Reject reviews mention 'incremental over prior work', 'gains marginal', 'cannot tell which component drove improvement'.

\textbf{Oral vs HC gap.} HC papers (n=20, coverage=16/20) execute the same audit-and-relax move but typically deliver a broadly-usable artifact (SAM, Tent, VICReg, SDEdit) whose appeal is the empirical generalizability of the technique enabled by the relaxation, often with theory in a follow-up role. Oral papers tend to deliver the relaxation itself as the primary contribution: the new theorem (e.g., identifiability under partial independence, T\textasciicircum{}(1/3) gap under continual observation, support-change identification) is the result, and the algorithm is a demonstration. Put differently, HC papers say 'because A' is enough, here is a method that works'; Oral papers say 'A' is exactly what is necessary and sufficient --- here is the theorem, and here is the method that realizes it'. HC papers also more often relax an assumption to enable a technique; Oral papers more often relax an assumption to expose a previously-invisible signal or capability (instantaneous dependencies as identification source, support changes as identifiability, parameter dormancy as capacity loss).

\textbf{Reviewer expectations:}
\begin{itemize}[leftmargin=12pt,topsep=2pt,itemsep=1pt]
  \item \textit{[both]} Want a precise statement of the assumption being relaxed and the assumption replacing it --- vague 'principled relaxation' framings are flagged as insufficient.
  \item \textit{[reject\_reviews]} Demand ablations or counterexamples isolating the relaxed assumption from other architectural or training changes --- reject reviewers consistently fault papers that bundle multiple novelties without disentanglement.
  \item \textit{[oral\_reviews]} Praise papers that derive necessary AND sufficient conditions or prove sharp separations (impossibility, tight bound, exact equivalence) under the relaxed assumption.
  \item \textit{[both]} Require evaluation in the regime where the original assumption fails sharply --- benchmark performance under regimes where the assumption approximately holds is treated as not informative about the contribution.
  \item \textit{[reject\_reviews]} Require explicit positioning against prior work that may have implicitly used the same relaxation --- when the auditor is unaware of close predecessors, reviewers downgrade the contribution to incremental.
  \item \textit{[oral\_reviews]} Reward minimal interventions (one regularizer, one architectural change, one substitution) over comprehensive frameworks, because minimality makes attribution to the audited assumption credible.
\end{itemize}
\textbf{Cognitive barriers:}
\begin{itemize}[leftmargin=12pt,topsep=2pt,itemsep=1pt]
  \item Field-wide convergence on a paradigm makes its embedded assumptions invisible --- practitioners stop seeing them as choices and treat them as facts (e.g., 'more fine-tuning is always better', 'rare elements are class-imbalance', 'global attention is the fix for local aggregation', 'continual release adds only log-factor overhead').
  \item Cross-disciplinary categorization blind spots: a tool from one community (causal inference's IPW, optimization's local error bounds, signal processing's harmonic analysis, cryptography's ZKPs) is filed under one problem class, so researchers in adjacent classes never consult it --- the audit requires recognizing that an assumption in one's own field has been resolved as a methodological choice in another.
  \item Treating empirically-observed correlates as diagnostic side-effects rather than as objectives that can be directly imposed (PAG correlated with robustness $\to$ train PAG directly; flat minima correlate with generalization $\to$ train flatness directly; dormant neurons correlate with stalled RL $\to$ reset directly).
  \item Confidence asymmetry: the original assumption was reinforced by repeated empirical successes or by negative results from earlier attempts to relax it, making the cost of questioning it appear higher than the expected gain --- until the gain is shown to be a phase transition, not an increment.
\end{itemize}
\textbf{Representative examples:}
\begin{itemize}[leftmargin=12pt,topsep=2pt,itemsep=1pt]
  \item \textbf{[Oral]} \texttt{ICLR\_2022\_0094} \textemdash\ Cleanly inverts the implicit 'full fine-tuning dominates linear probing' assumption with a two-layer linear theory that pinpoints head initialization as the controlling variable, then derives LP-FT as the minimal remedy.
  \\ \textcolor{gray}{\footnotesize \textit{Lesson:} Pin the controlling variable with a minimal theory (here: a 2-layer linear model) before proposing the remedy --- this lets reviewers verify head initialization, not capacity, drove prior failure.}
  \item \textbf{[Oral]} \texttt{ICML\_2023\_0551} \textemdash\ Proves the field's assumed log T overhead for differential privacy under continual observation is wrong by a polynomial T\textasciicircum{}(1/3) factor via combinatorial reductions --- the audit is the theorem itself.
  \\ \textcolor{gray}{\footnotesize \textit{Lesson:} When the audit IS the theorem, the proof construction (combinatorial reductions) replaces the algorithmic novelty --- papers that overturn a field's assumed bound by constructive separation are read as foundational.}
  \item \textbf{[Oral]} \texttt{ICLR\_2025\_1525} \textemdash\ Reframes instantaneous dependencies --- eliminated as a nuisance in prior temporal causal representation work --- as an additional identifiability resource, deriving conditions under which more structure aids rather than hinders identification.
  \\ \textcolor{gray}{\footnotesize \textit{Lesson:} Look for nuisances the field has eliminated by convention --- they often hide identifiability resources. Reframing a constraint as a degree of freedom is the highest-yield move.}
  \item \textbf{[Oral]} \texttt{ICML\_2025\_1810} \textemdash\ Audits the unexamined assumption that better prediction always improves equity outcomes, then proves analytically that capacity expansion can dominate prediction improvement under specific distributional conditions.
  \\ \textcolor{gray}{\footnotesize \textit{Lesson:} Audit even the unexamined assumptions in adjacent welfare / equity claims --- analytic dominance results expose tradeoffs benchmarks cannot.}
  \item \textbf{[Reject]} \texttt{ICLR\_2023\_0388} \textemdash\ Identifies the right hidden assumption (fairness tradeoffs are distribution-dependent under shift) but reviewers found theorems unclear, baselines missing, and Pareto-frontier results unconvincing --- the relaxation lacked rigorous re-derived guarantees.
  \\ \textcolor{gray}{\footnotesize \textit{Lesson:} Identifying the right hidden assumption is necessary but not sufficient --- without a re-derived guarantee that pins the relaxation to a sharp condition, the contribution reads as heuristic.}
  \item \textbf{[Reject]} \texttt{ICLR\_2024\_0767} \textemdash\ Relaxes the uniform-flooding-constant assumption to per-instance adaptive flooding but cannot demonstrate via ablations that uniformity (rather than capacity or noise) was the binding constraint.
  \\ \textcolor{gray}{\footnotesize \textit{Lesson:} The adaptive replacement must be ablated against capacity / noise alternatives --- otherwise reviewers cannot confirm uniformity (rather than a confounder) was the binding constraint.}
  \item \textbf{[Reject]} \texttt{ICLR\_2025\_1644} \textemdash\ Borrows VICReg's variance-covariance regularization to supervised learning under the premise that standard objectives implicitly suppress feature diversity, but reviewers found novelty insufficient and the suppression-is-the-bottleneck claim under-supported.
  \\ \textcolor{gray}{\footnotesize \textit{Lesson:} Borrowing a regularizer from another setting requires showing the suppressed property is the bottleneck in the new setting; without that, the borrow looks unmotivated.}
\end{itemize}
\end{tcolorbox}

\begin{tcolorbox}[colback=bgskill, colframe=cblue!70, title={\textbf{Encode Invariance as Hard Constraint} \textcolor{gray}{\small\texttt{invariance\_as\_architectural\_constraint}} \hfill confidence: \textbf{medium}}, breakable, before skip=4pt, after skip=4pt]
\textcolor{gray}{\small Oral $n=62$ (meta 48/62) | HC $n=3$ (meta 2/3) | Reject $n=36$ (meta 36/36)}

\textbf{Sample note.} \textit{Oral (n=62) and Reject (n=36) samples are large with strong meta-review coverage on the Reject side (100\%) and adequate coverage on the Oral side (77\%); the HC sample (n=3) is too small to support an Oral/HC contrast, so that field is intentionally limited.}

\textbf{Step-by-Step:}
\begin{enumerate}[leftmargin=14pt,topsep=2pt,itemsep=1pt]
  \item Characterize the symmetry / invariance / equivariance the data carries --- group structure, geometry, conservation law, algebraic property --- name it precisely.
  \item Identify what current architectures handle softly (data augmentation, soft penalty, ignored). Cite the prior baseline that pays for this softness.
  \item Design operations that commute with or preserve the symmetry by construction --- every layer / step provably equivariant or invariant.
  \item Prove (or constructively verify) that the entire computation inherits the symmetry --- local equivariance composes to global equivariance.
  \item Quantify the structural advantage: parameter count, sample complexity, generalization, OR a universal-approximation result that pre-empts expressivity objections.
  \item Evaluate against augmentation-trained and soft-penalty baselines on tasks where the symmetry holds --- zero advantage means the symmetry was not load-bearing.
\end{enumerate}
\textbf{Success conditions} (Oral):
\begin{itemize}[leftmargin=12pt,topsep=2pt,itemsep=2pt]
  \item \textbf{Prove (not just assert) that the local constraint propagates to the global property required by the task --- e.g., that an SE(3)-equivariant kernel induces an SE(3)-invariant stationary distribution, or that ensemble averaging marginalizes the group orbit.} \textcolor{gray}{\footnotesize evidence: \texttt{ICLR\_2022\_0097}, \texttt{ICML\_2024\_1028}, \texttt{ICLR\_2022\_0096}} \\ Oral papers consistently include a theorem/lemma that the chosen architectural device (kernel, averaging operator, frame) actually delivers the claimed invariance through the rest of the pipeline. Reviewers explicitly cite this 'theoretical argument that local property implies global property' as the load-bearing contribution rather than the engineering.
  \item \textbf{Enumerate the FULL symmetry group of the target space before designing the architecture --- including non-obvious classes (gating translations, scaling, orthogonal/invertible) --- and resolve every class.} \textcolor{gray}{\footnotesize evidence: \texttt{NeurIPS\_2024\_1303}, \texttt{NeurIPS\_2025\_1878}, \texttt{NeurIPS\_2025\_1908}} \\ Oral metanetwork/LMC papers all begin with a complete algebraic characterization and explicitly call out previously-missed symmetry classes; partial enumeration is treated as a defect, not a simplification. The framing is 'we missed X, here is the complete list, here is the algorithm that handles each.'
  \item \textbf{Borrow an existing algebraic or physical structure (Brauer algebra, Clifford algebra, Schrödinger bridge, Hamiltonian mechanics, finite group Fourier theory) so that invariance is inherited rather than re-engineered.} \textcolor{gray}{\footnotesize evidence: \texttt{ICML\_2023\_0424}, \texttt{NeurIPS\_2023\_0605}, \texttt{ICLR\_2025\_1600}, \texttt{NeurIPS\_2025\_1903}} \\ Oral papers route the invariance through a mathematically canonical object (Clifford group, transposition noise on S\_n, time-reversal symmetry of conservative dynamics). This buys both rigor and a natural class of operations, and reviewers reward the bridging of abstract math into network design.
  \item \textbf{Achieve invariance WITHOUT collapsing expressivity --- show universal approximation, recover prior special cases, or pair the invariance with a complementary property (locality, multi-scale, hierarchy) that sharpens what the model can still discriminate.} \textcolor{gray}{\footnotesize evidence: \texttt{ICLR\_2022\_0096}, \texttt{ICML\_2024\_1030}, \texttt{ICML\_2023\_0544}, \texttt{ICML\_2024\_1033}} \\ Oral submissions defuse the 'equivariance hurts capacity' objection up front: Frame Averaging proves universality, EquiPocket couples E(3) with locality, Subequivariant GNN hierarchically combines global+local symmetry, fragment-bias work shows expressivity AND generalization improve together.
  \item \textbf{Replace a costly canonicalization, augmentation, or alignment pipeline with a constraint that is provably equivalent but cheaper, and quantify the saved cost.} \textcolor{gray}{\footnotesize evidence: \texttt{ICLR\_2022\_0096}, \texttt{ICLR\_2023\_0323}, \texttt{ICLR\_2023\_0243}} \\ Oral papers frequently sell the constraint as removing a previously-accepted overhead: Frame Averaging avoids averaging the full group, Relative Representations eliminates Procrustes/CCA, Git Re-Basin replaces the assumption of distinct basins entirely. The contrast against the prior costly workflow is part of the pitch.
  \item \textbf{Apply at the right level of the pipeline: bake the symmetry into the noise/transition/message-passing operator rather than only the final layer or as a post-hoc projection.} \textcolor{gray}{\footnotesize evidence: \texttt{ICLR\_2022\_0097}, \texttt{NeurIPS\_2024\_1225}, \texttt{ICLR\_2025\_1600}, \texttt{ICLR\_2024\_0792}} \\ Oral generative-model papers all install the constraint in the per-step kernel (graph-dependent noise, SE(3)-equivariant denoiser, transposition noise, conservation-respecting drift). Reviewers treat this as the difference between a method that 'respects' a constraint and one that merely approximates it.
\end{itemize}
\textbf{Failure modes} (Reject):
\begin{itemize}[leftmargin=12pt,topsep=2pt,itemsep=2pt]
  \item \textbf{Constraint added as a soft loss/regularizer rather than enforced by construction, then evaluated against generic baselines rather than against the architectural-enforcement alternative.} \textcolor{gray}{\footnotesize evidence: \texttt{ICLR\_2025\_1534}, \texttt{ICLR\_2024\_0789}, \texttt{NeurIPS\_2023\_0686}} \\ Reviewers of these papers note that 'physics-informed' or 'geometry-aware' losses do not deliver the hard guarantees the framing implies, and ablations do not isolate whether the structural term actually matters versus an unconstrained model with similar capacity.
  \item \textbf{Architectural choice that constrains expressivity (e.g., bilinear-only, linear-only, or a single basis class) without showing it can match the more general equivariant family it restricts.} \textcolor{gray}{\footnotesize evidence: \texttt{NeurIPS\_2023\_0721}, \texttt{ICLR\_2025\_1519}} \\ Meta-reviews explicitly call out that the simplified equivariant construction is a special case of more general known architectures, and the experiments do not demonstrate parity with those richer alternatives --- making the 'simpler invariant design' look like a capacity downgrade.
  \item \textbf{Domain port of a known equivariant/invariant idea where the new contribution is an engineering rewiring rather than a new structural insight, and head-to-head comparison against the unconstrained baseline shows only marginal lift.} \textcolor{gray}{\footnotesize evidence: \texttt{NeurIPS\_2024\_1324}, \texttt{ICML\_2025\_1692}, \texttt{NeurIPS\_2024\_1332}} \\ Reviewers acknowledge the construction is sound but reject because the empirical advantage of the invariance is not demonstrated to be the cause of any measured improvement, and missing analyses of symmetry/invariance properties leave the central claim unsupported.
  \item \textbf{Claim of theoretical completeness (a 'complete invariant' or 'unified framework') without the empirical experiments needed to ground the claim in a downstream task.} \textcolor{gray}{\footnotesize evidence: \texttt{ICML\_2025\_1692}, \texttt{ICLR\_2025\_1519}, \texttt{NeurIPS\_2024\_1228}} \\ Across rejects, AC notes the math is correct and elegant but cannot be assessed as useful: no scaling study, no comparison against architectures that satisfy the constraint differently, no real downstream evaluation --- so the contribution remains a definition rather than a method.
  \item \textbf{Constraint motivated by domain physics or biology but the writing leaves the symmetry class informal, lacks a precise statement of what is preserved, and conflates multiple notions (equivariance vs. invariance vs. consistency).} \textcolor{gray}{\footnotesize evidence: \texttt{ICLR\_2024\_0848}, \texttt{ICLR\_2024\_0811}, \texttt{ICLR\_2024\_0789}} \\ Reject meta-reviews flag mathematical statements that are not properly stated as theorems, no proof for the central proposition, and unclear positioning relative to existing equivariant baselines --- reviewers cannot verify whether the constraint actually holds.
\end{itemize}
\textbf{Oral vs Reject gap.} Oral papers couple the architectural device to a proof (or near-proof) that the desired invariance survives composition through every subsequent operation --- kernels, denoising chains, message passing, ensemble averaging --- and they enumerate the full symmetry group rather than picking one obvious class. Rejects either install the constraint as a soft penalty (PIAE for CO2, centroid/orientation losses) or pick a single restricted construction (bilinear-only, linear superposition only) without proving completeness or empirically beating the unconstrained baseline at the constraint's job. Oral papers also routinely pair invariance with a second axis --- universal-approximation guarantee, locality, hierarchy, or replacement of a known costly pipeline (Procrustes, group averaging, BPTT) --- so the contribution is 'constraint plus capability', whereas rejects offer 'constraint alone' and leave reviewers unable to attribute any measured improvement to the invariance itself.

\textbf{Oral vs HC gap.} HC sample (n=3) too small for reliable contrast. The single most directly aligned HC paper (Caduceus, ICML\_2024\_1007) shares the Oral pattern of enforcing symmetry through weight-tying and input decomposition rather than augmentation, suggesting no obvious behavioral gap, but three papers cannot support a confident claim.

\textbf{Reviewer expectations:}
\begin{itemize}[leftmargin=12pt,topsep=2pt,itemsep=1pt]
  \item \textit{[both]} A clean theorem or proof that the local architectural property (equivariant kernel, invariant operator, conservation-respecting drift) implies the global property the task requires --- not just a verbal assurance.
  \item \textit{[oral\_reviews]} Complete characterization of the symmetry group of the target space, with all classes resolved; reviewers reward extensions that uncover previously-missed symmetries (gating translation, scaling, invertible).
  \item \textit{[reject\_reviews]} Comparison against architectures that solve the same problem with a DIFFERENT enforcement mechanism (data augmentation, soft penalty, post-hoc canonicalization) so that the contribution of architectural enforcement is isolated.
  \item \textit{[both]} Evidence that the imposed constraint does not collapse expressivity --- universal-approximation result, parity with a less-constrained baseline, or recovery of known special cases.
  \item \textit{[reject\_reviews]} Precise mathematical statements (formal definitions, properly numbered theorems with proofs); reviewers reject when 'invariance' is invoked informally or when a central proposition appears without a proof.
  \item \textit{[reject\_reviews]} Empirical demonstration that the constraint pays off on a downstream task that domain practitioners care about, not only on synthetic benchmarks designed to highlight the symmetry.
\end{itemize}
\textbf{Cognitive barriers:}
\begin{itemize}[leftmargin=12pt,topsep=2pt,itemsep=1pt]
  \item The relevant symmetry is invisible at the function-class level --- it only appears in the optimization algorithm, the representation's parameter space, or the gating structure --- so practitioners trained on input-space symmetry never look there (ICLR\_2021\_0076, NeurIPS\_2024\_1303, NeurIPS\_2025\_1908).
  \item Equivariance and ensemble/randomized-averaging come from different research communities, so seeing that random initialization marginalizes the group, or that algebraic constructions from pure math can be a parameter basis, requires bridging traditions that rarely speak (ICML\_2024\_1028, ICML\_2023\_0424, NeurIPS\_2023\_0605).
  \item The dominant practice --- induce symmetry through data augmentation or soft penalties --- actively suppresses the question 'why not bake it into the weights?' because the augmentation answer feels sufficient and architecturally cheaper (ICML\_2024\_1007, ICLR\_2025\_1900).
  \item It feels intuitively necessary to average over the full group or canonicalize globally to guarantee invariance; the move to a small representative subset, a per-step kernel, or a relative representation requires trusting a non-trivial propagation argument over an obvious-looking workaround (ICLR\_2022\_0096, ICLR\_2022\_0097, ICLR\_2023\_0323).
\end{itemize}
\textbf{Representative examples:}
\begin{itemize}[leftmargin=12pt,topsep=2pt,itemsep=1pt]
  \item \textbf{[Oral]} \texttt{ICLR\_2022\_0097} \textemdash\ GeoDiff is the canonical 'local-implies-global' invariance argument: prove that SE(3)-equivariant transition kernels make the entire diffusion chain induce an SE(3)-invariant distribution, then build kernels that satisfy the condition.
  \\ \textcolor{gray}{\footnotesize \textit{Lesson:} The cleanest invariance argument proves local property $\to$ global property --- equivariant transition kernels make the entire chain inherit SE(3) invariance, no side-loss needed.}
  \item \textbf{[Oral]} \texttt{ICLR\_2022\_0096} \textemdash\ Frame Averaging defuses the 'you need the whole group' intuition by proving exact invariance from averaging over a tiny representative subset, plus a universal-approximation result that pre-empts the expressivity objection.
  \\ \textcolor{gray}{\footnotesize \textit{Lesson:} Defuse the 'you need the whole group' intuition with averaging over a representative subset, plus a universal-approximation theorem --- this pre-empts the standard expressivity attack.}
  \item \textbf{[Oral]} \texttt{NeurIPS\_2024\_1303} \textemdash\ Scale Equivariant GMNs exemplifies the 'enumerate-all-symmetries-then-fix-the-missing-one' pattern: catalogs symmetries of weight space, identifies that scaling has not been enforced, designs equivariant message passing for it.
  \\ \textcolor{gray}{\footnotesize \textit{Lesson:} Catalog all symmetries of the input space first, then pick the un-enforced one --- completeness is a generative move that exposes architectural gaps.}
  \item \textbf{[Oral]} \texttt{ICML\_2023\_0424} \textemdash\ Brauer's Group Equivariant Networks shows the 'borrow an existing algebraic structure' move at its purest: lift Brauer algebra diagrams from pure math directly into a concrete weight basis for O(n)/SO(n)/Sp(n)-equivariant layers.
  \\ \textcolor{gray}{\footnotesize \textit{Lesson:} Lift existing algebraic structure (Brauer diagrams) directly into a weight basis --- borrowing from pure math at the right level of abstraction is faster than re-deriving the equivariance machinery.}
  \item \textbf{[Reject]} \texttt{ICLR\_2025\_1534} \textemdash\ Physics-Informed Autoencoder for CO2 illustrates the soft-penalty failure: the SDE governing the dynamics is encoded as a loss term rather than an architectural constraint, and the paper does not ablate against an unconstrained baseline to show the physics term matters.
  \\ \textcolor{gray}{\footnotesize \textit{Lesson:} Encoding a constraint as a soft loss term loses the structural guarantee an architectural constraint provides --- and skipping the unconstrained-baseline ablation makes the soft-penalty contribution invisible.}
  \item \textbf{[Reject]} \texttt{NeurIPS\_2023\_0721} \textemdash\ Bilinear Tensor Networks restricts the equivariant family to a special case (bilinear-only) without demonstrating parity with the more general equivariant architectures it implicitly competes against --- reviewers ask for the head-to-head that would justify the constraint.
  \\ \textcolor{gray}{\footnotesize \textit{Lesson:} Restricting to a special case of an equivariant family demands head-to-head parity with the more general family --- without that comparison, the constraint looks unmotivated.}
  \item \textbf{[Reject]} \texttt{ICML\_2025\_1692} \textemdash\ The complete SE(n)-invariant for Euclidean graphs reaches mathematical completeness but never demonstrates that the invariant pays off on a real downstream task --- illustrating the 'definition not method' rejection pattern.
  \\ \textcolor{gray}{\footnotesize \textit{Lesson:} Mathematical completeness is not a downstream contribution --- invariants must be paired with an evaluation showing they enable real tasks.}
\end{itemize}
\end{tcolorbox}

\begin{tcolorbox}[colback=bgskill, colframe=cblue!70, title={\textbf{Reframe via Alternative Formalism} \textcolor{gray}{\small\texttt{formal\_reframing\_via\_equivalence}} \hfill confidence: \textbf{medium}}, breakable, before skip=4pt, after skip=4pt]
\textcolor{gray}{\small Oral $n=162$ (meta 103/162) | HC $n=12$ (meta 10/12) | Reject $n=85$ (meta 84/85)}

\textbf{Sample note.} \textit{n\_oral=162 and n\_reject=85 are well above thresholds, but the methodology is broad and many tagged Oral papers are only partial instances of formal reframing (some are closer to architecture transfer or empirical reformulation), reducing the signal density; HC sample is at n=12 with two papers (RFA, FreeMatch) only loosely fitting the methodology, limiting the Oral-vs-HC contrast.}

\textbf{Step-by-Step:}
\begin{enumerate}[leftmargin=14pt,topsep=2pt,itemsep=1pt]
  \item Identify problem P that is currently intractable, theoretically opaque, or stuck in its native formalism.
  \item Search for a neighboring formalism F (algebraic, geometric, probabilistic, information-theoretic, category-theoretic) whose native objects map cleanly to P's objects.
  \item Construct an exact correspondence P $\leftrightarrow$ P' in F --- equivalence theorem with a constructive direction (not just existence).
  \item Deploy F's native algorithms or theorems on P'; F's machinery should yield an algorithm or guarantee invisible under P's original framing.
  \item Map the result back to P: the new algorithm runs on P's inputs and produces P's outputs, but its analysis lives in F.
  \item Show the reframe pays a derivability dividend --- closed form, parallelizable algorithm, sharper bound, identifiability --- not just notational shuffling.
\end{enumerate}
\textbf{Success conditions} (Oral):
\begin{itemize}[leftmargin=12pt,topsep=2pt,itemsep=2pt]
  \item \textbf{Prove a bidirectional, constructive equivalence (not just analogy) between P and P', so that solutions in F can be mapped back to P with no loss of fidelity.} \textcolor{gray}{\footnotesize evidence: \texttt{ICLR\_2021\_0018}, \texttt{ICLR\_2022\_0141}, \texttt{ICLR\_2023\_0331}, \texttt{ICLR\_2023\_0380}, \texttt{NeurIPS\_2025\_1902}} \\ Oral papers do not stop at 'P looks like P'' --- they construct an explicit reduction with a back-mapping. EigenGame proves the Nash equilibrium of the constructed game equals the PCA solution; Hidden Convex Landscape gives a constructive convex cone whose solutions reproduce all global ReLU optima; SAM is shown to coincide exactly with the Fenchel biconjugate of a Bayes objective; Transformers$\to$automata maps depth bounds to a specific algebraic invariant of the transformation semigroup.
  \item \textbf{Show that the imported framework's native guarantees (algorithms, rates, optimality, identifiability) survive the round-trip and produce a concrete computational or theoretical payoff in the original problem.} \textcolor{gray}{\footnotesize evidence: \texttt{ICLR\_2023\_0223}, \texttt{ICML\_2023\_0491}, \texttt{NeurIPS\_2025\_1884}, \texttt{ICML\_2025\_1706}, \texttt{ICML\_2023\_0531}} \\ Adversarial Team Markov Games imports LP duality + non-standard KKT to obtain polynomial scaling; Multicalibration-as-Boosting transfers the swap-regret oracle to give a fairness algorithm without realizability; High-Dimensional Calibration reduces to OLO swap regret to inherit explicit complexity; EVIs use the distributional relaxation that makes correlated-equilibrium tractable to solve nonmonotone VIs in poly time; Caratheodory decomposition turns exponential action sets into polynomial-time updates.
  \item \textbf{Recover existing results, methods, or empirical heuristics as special cases of the new framing, providing retroactive justification and a unified taxonomy.} \textcolor{gray}{\footnotesize evidence: \texttt{ICLR\_2022\_0083}, \texttt{ICLR\_2023\_0331}, \texttt{NeurIPS\_2023\_0753}, \texttt{ICML\_2025\_1662}, \texttt{ICLR\_2025\_1437}} \\ Reviewers consistently praise reframings that subsume prior art: the JSD/MMD discrepancy paper recovers both as decision-theoretic special cases; SAM$\leftrightarrow$Bayes recovers SAM exactly as the variational lower bound; the diffusion-ELBO paper unifies a family of empirically tuned weightings as a single ELBO with augmentation; VES recovers EI as the exponential-distribution case; Generator Matching recovers flow matching as the vector-field special case of arbitrary Markov generators.
  \item \textbf{Close a previously open or ill-posed problem in P by exploiting machinery that exists only in F (e.g., topology, algebraic decomposition, convex duality, optimal transport, complexity theory).} \textcolor{gray}{\footnotesize evidence: \texttt{ICLR\_2021\_0027}, \texttt{ICLR\_2023\_0356}, \texttt{NeurIPS\_2023\_0749}, \texttt{ICML\_2025\_1726}, \texttt{ICML\_2024\_1120}} \\ Three-layer mean-field convergence uses algebraic topology to prove a time-uniform expressivity lemma that closes the global-convergence question; Symbolic Physics Learner uses MCTS exploration-exploitation theory to obtain principled symbolic regression; CoT theory uses circuit complexity to prove formal separations; Tubal tensors reduce to matrix factorization via the frequency-domain block-diagonal product to import implicit-bias proofs; Doubly-Efficient Debate gives formal complexity-theoretic safety guarantees.
  \item \textbf{Bridge two communities with disjoint vocabularies --- the contribution is the translation layer itself, often validated by reviewers as 'unifying' or 'opening new directions.'} \textcolor{gray}{\footnotesize evidence: \texttt{NeurIPS\_2023\_0631}, \texttt{ICML\_2023\_0535}, \texttt{ICLR\_2023\_0286}, \texttt{ICLR\_2025\_1623}, \texttt{ICML\_2025\_1804}} \\ PGF-based exact inference reframes Bayesian inference as automatic differentiation; SimRFs imports random-matrix-theory simplex coupling into kernel approximation; Action-Impact Embedding reframes labels as observed effects (perception $\leftrightarrow$ planning); Logical Framework for GNN Expressiveness builds a complete bidirectional translation between GNN architectures and FO logic; TD Flows lifts the Bellman equation from scalars to probability paths to inherit flow-matching.
\end{itemize}
\textbf{Failure modes} (Reject):
\begin{itemize}[leftmargin=12pt,topsep=2pt,itemsep=2pt]
  \item \textbf{Reframing-by-analogy without a proven equivalence: paper claims P resembles P' but provides no formal mapping showing F's guarantees transfer back.} \textcolor{gray}{\footnotesize evidence: \texttt{ICLR\_2023\_0203}, \texttt{ICLR\_2025\_1554}, \texttt{ICLR\_2025\_1351}, \texttt{ICLR\_2024\_0828}, \texttt{ICLR\_2025\_1411}} \\ Reviewers reject these as 'reformulation does not constitute methodological novelty.' Differentiable Logic Programming dropped discrete operations for differentiable proxies but lost soundness guarantees; Reformer cast kernel selection as classification with no theoretical link; AlphaIntegrator wrapped symbolic integration in RL without a verified action-validity equivalence; ExID and Quantum Resource Scheduling recast classical problems as MDPs without justifying why MDP machinery is well-suited.
  \item \textbf{The new formulation loses or weakens the original problem's guarantees, so the reframed solution is not actually a solution to P.} \textcolor{gray}{\footnotesize evidence: \texttt{NeurIPS\_2023\_0591}, \texttt{ICLR\_2024\_0941}, \texttt{ICML\_2025\_1675}, \texttt{ICLR\_2025\_1595}} \\ Unbiased NE estimation produced an objective whose stationary points need not be Nash equilibria; consistency-model convergence relied on a Lipschitz constant the authors could not bound; meta-learned regret minimization preserved CCE convergence but reviewers found the Nash gap claim unsupported; constraint inference paper had only approximate constraint satisfaction. Reviewers explicitly cite 'no finite-time guarantees of reaching a Nash equilibrium' or 'theoretical results crucially depend on notions\ldots{} not justified.'
  \item \textbf{Generalization-by-extension that adds an axis (more dimensions, more agents, more modalities) without a non-trivial structural insight beyond the prior reframing.} \textcolor{gray}{\footnotesize evidence: \texttt{ICLR\_2023\_0173}, \texttt{ICLR\_2024\_0975}, \texttt{ICLR\_2024\_0954}, \texttt{ICML\_2025\_1754}} \\ Many-to-one matching markets just generalized one-to-one; extensive-form linear function approximation extended fully observed MDP analysis; data-dependent couplings for stochastic interpolants overlapped substantially with Pooladian et al. and Liu et al.; meta-learning for self-play regret minimization extended one-sided meta-learning. The reframing exists but reviewers conclude the technical novelty over the prior reframing is 'slim.'
  \item \textbf{Empirical equivalence claim unsupported by the formal bridge required by the methodology --- the authors assert structural similarity, but experiments do not isolate the reframing from confounding architectural changes.} \textcolor{gray}{\footnotesize evidence: \texttt{ICLR\_2023\_0190}, \texttt{ICLR\_2023\_0214}, \texttt{ICLR\_2023\_0220}, \texttt{ICLR\_2025\_1483}} \\ Convolutions-vs-Transformers for proteins reframed sequence modeling but reviewers found the empirical comparison did not isolate the architectural claim; FGSO joined GSOs and Fourier with marginal gains and no theoretical motivation for the new operator; model-based Quality-Diversity combined two paradigms without showing the combination is structurally principled; G-SFormer claimed graph-structured attention as the key but reviewers found ablations insufficient.
  \item \textbf{Cross-domain transplant that ignores why the source-domain machinery worked: the paper imports a tool but the assumptions that made it powerful in the source do not hold in the target.} \textcolor{gray}{\footnotesize evidence: \texttt{ICLR\_2024\_0916}, \texttt{ICLR\_2025\_1370}, \texttt{NeurIPS\_2025\_1889}, \texttt{ICLR\_2024\_0976}} \\ Quaternion DFT/convolution proved spectral identities but reviewers found the architectural niche too narrow to justify the algebraic machinery; Cayley Maze proposed a beautiful group-theoretic environment family but lacked empirical evaluation; UDV factorization made implicit bias 'explicit' but provided no theoretical guarantee that the constraint actually realizes the desired bias; TinT simulation was a sophisticated reduction but reviewers questioned soundness of the chained approximations.
  \item \textbf{Reformulation produces a tractable algorithm but the resulting solution concept is uninformative or behaviorally equivalent to a much simpler baseline.} \textcolor{gray}{\footnotesize evidence: \texttt{NeurIPS\_2024\_1240}, \texttt{ICML\_2025\_1671}, \texttt{ICML\_2024\_1085}} \\ K-recall poker abstraction added bounded history but reviewers found the framework redundant with existing solvers; APPD implementation paper provided an O(n$^2$) algorithm reviewers said already existed; Personhood/accountability reframing replaced one framework with another without showing the new one yields different decisions in concrete cases.
\end{itemize}
\textbf{Oral vs Reject gap.} Oral papers ground the reframing in a proof obligation: they construct an explicit, often constructive, mapping P$\leftrightarrow$P', verify that the imported framework's solution concepts collapse to (or exactly equal) the originals (EigenGame's Nash = PCA solution; SAM = Fenchel biconjugate of Bayes; CoT separation = circuit lower bound), and demonstrate that algorithms or rates from F yield a quantifiable payoff back in P (polynomial scaling, exact constraint satisfaction, closed-form parameterization). Reject papers stop at structural analogy --- 'X is like Y' --- and either (i) leave the equivalence as motivation rather than a theorem, (ii) prove it only under assumptions that void the original problem (e.g., interior NE, bounded data, infinite-width limits the experiments don't satisfy), or (iii) substitute the original solution concept with a strictly weaker one (CCE for Nash, approximate for exact, marginal for joint) without an argument that the substitution preserves what the user actually wanted. Oral papers also typically recover at least one prior method as a special case of the new framing; reject papers rarely do, leaving reviewers unable to see what was unified.

\textbf{Oral vs HC gap.} HC sample (n=12) too small for reliable contrast --- and inspection shows several HC papers (Random Feature Attention, FreeMatch, Selective Frequency Network) are only loosely instances of formal reframing rather than canonical examples. The few HC papers that do reframe (Mamba/SSM-attention duality, LLM-as-optimizer) tend to deliver strong empirical results from a clean reformulation but lack the closed-form equivalence or recovery-as-special-case proofs that distinguish Oral papers; the Oral bar appears to be 'the reframing produces a theorem,' while HC papers can succeed by 'the reframing produces a competitive algorithm.'

\textbf{Reviewer expectations:}
\begin{itemize}[leftmargin=12pt,topsep=2pt,itemsep=1pt]
  \item \textit{[oral\_reviews]} Demand a formal, often bidirectional, equivalence --- not just an analogy or motivational mapping. Reviewers explicitly praise 'beautiful connection,' 'non-trivial KKT argument,' 'exact reformulation,' and 'characterization theorem.'
  \item \textit{[oral\_reviews]} Expect the reframing to recover at least one well-known prior method as a special case, used as evidence that the new framing is genuinely unifying rather than merely renaming.
  \item \textit{[reject\_reviews]} Reject papers when the new formulation appears to lose guarantees that motivated the original problem ('no finite-time convergence to NE,' 'reduces to weaker solution concept,' 'assumption is too strong / artificial / unjustified').
  \item \textit{[both]} Probe whether the import of a specialist framework (topology, group theory, optimal transport, complexity theory) actually pays off in the target domain --- reviewers reject if the algebra is impressive but the architectural niche or empirical lift is small.
  \item \textit{[reject\_reviews]} Question incremental extensions of an existing reframing (one-to-one $\to$ many-to-one, scalar $\to$ distributional, single-agent $\to$ multi-agent) when the structural insight beyond the prior reframing is unclear.
  \item \textit{[oral\_reviews]} Look for cross-disciplinary bridges that 'open a new direction' or 'inspire future research,' especially when the bridge is between communities (e.g., game theory $\leftrightarrow$ generative models, algebra $\leftrightarrow$ deep learning) that rarely cite each other.
\end{itemize}
\textbf{Cognitive barriers:}
\begin{itemize}[leftmargin=12pt,topsep=2pt,itemsep=1pt]
  \item Disciplinary vocabulary lock-in: the source community describes the problem in terms (calibration, attention, score-matching, regression) that obscure its structural identity with an object the target community has fully developed (swap regret, kernel, IRL, classification).
  \item Working backward from a known empirical winner to the right mathematical bridge feels post-hoc and unprincipled, so practitioners avoid it even though that is exactly what successful reframings (SAM$\leftrightarrow$Bayes, MES$\leftrightarrow$EI, diffusion-as-ELBO) do.
  \item Apparent loss of structure: the original problem's natural framing (non-convex, discrete, exponential, intractable) is held as essential, masking the existence of a lifted or dualized form (convex, distributional, polynomial, tractable) that preserves all the information.
  \item Standard ranking metrics (WL hierarchy, expressivity, regret, depth) acquire foundational status, blocking the realization that an alternative algebraic or combinatorial measure (homomorphism counts, polynomial degree, swap regret, lattice height) gives finer control with less effort.
\end{itemize}
\textbf{Representative examples:}
\begin{itemize}[leftmargin=12pt,topsep=2pt,itemsep=1pt]
  \item \textbf{[Oral]} \texttt{ICLR\_2021\_0018} \textemdash\ Canonical example: utility-function design transforms PCA (a classical centralized eigenproblem) into a game whose Nash equilibrium is provably the PCA solution, and the meta-review specifically highlights that this exact algebraic structure enables a parallelizable, orthonormalization-free algorithm.
  \\ \textcolor{gray}{\footnotesize \textit{Lesson:} The reframe must enable an algorithm the original framing forbids --- here, parallelizability via game-theoretic structure that centralized eigendecomposition cannot expose.}
  \item \textbf{[Oral]} \texttt{ICLR\_2022\_0141} \textemdash\ Provides a constructive --- not merely existential --- convex reformulation of regularized two-layer ReLU optimization, showing how to reconstruct any optimal network from a solution to the convex program; the equivalence is exact, not an approximation.
  \\ \textcolor{gray}{\footnotesize \textit{Lesson:} The equivalence must be exact and constructive --- given a solution to the convex program, you can reconstruct the optimal network. Approximate equivalences invite 'loose surrogate' attacks.}
  \item \textbf{[Oral]} \texttt{ICLR\_2023\_0331} \textemdash\ Works backward from an empirically successful heuristic (SAM) to identify it as the Fenchel biconjugate of an expected log-loss in a Bayesian objective --- the meta-review calls this 'a beautiful connection' that retroactively grounds an opaque method.
  \\ \textcolor{gray}{\footnotesize \textit{Lesson:} Working backward from an empirically successful method to its formal characterization is a legitimate move --- retroactive grounding is as valuable as forward derivation when the method was previously opaque.}
  \item \textbf{[Oral]} \texttt{ICLR\_2023\_0380} \textemdash\ Connects learned transformer depth bounds to the algebraic group-theoretic classification of automata via transformation semigroups, predicting depth O(log T) generally and O(1) for solvable groups --- a cross-disciplinary bridge that classical CS theory enables.
  \\ \textcolor{gray}{\footnotesize \textit{Lesson:} Borrow classification machinery from classical theoretical CS at the right structural level --- depth bounds for transformers fall out of automata theory once the bridge is established.}
  \item \textbf{[Reject]} \texttt{ICLR\_2023\_0203} \textemdash\ Reframed discrete logic programming as continuous optimization via differentiable proxies, but the paper substitutes operations rather than proving an equivalence, and reviewers found the experimental section too thin to demonstrate the substitution preserves the original solution concept.
  \\ \textcolor{gray}{\footnotesize \textit{Lesson:} Substituting operations is not equivalence --- without a theorem proving the substitution preserves the original solution concept, the reframe is unjustified analogy.}
  \item \textbf{[Reject]} \texttt{NeurIPS\_2023\_0591} \textemdash\ Recasts Nash equilibrium computation as an unbiased stochastic optimization problem, but the reformulation only guarantees stationary points of the surrogate objective rather than Nash equilibria, and the assumption of an interior NE is itself solvable in polynomial time --- undermining the motivation.
  \\ \textcolor{gray}{\footnotesize \textit{Lesson:} A reformulation that only guarantees stationary points of a surrogate fails when the original problem demanded a stronger property --- the equivalence's direction must reach the actual goal, not a weaker one.}
  \item \textbf{[Reject]} \texttt{ICML\_2025\_1675} \textemdash\ Meta-learns regret-minimization updates to push toward Nash in general-sum games, but reviewers questioned whether the parameterization preserves the CCE convergence guarantee while actually closing the Nash gap, with no clean theorem to anchor the empirical claim.
  \\ \textcolor{gray}{\footnotesize \textit{Lesson:} Empirically pushing toward a target without a theorem anchoring the parameterization to the convergence guarantee leaves reviewers unable to assess whether the reframe preserves the original property.}
\end{itemize}
\end{tcolorbox}

\begin{tcolorbox}[colback=bgskill, colframe=cblue!70, title={\textbf{Decompose and Recombine} \textcolor{gray}{\small\texttt{decomposition\_into\_tractable\_stages}} \hfill confidence: \textbf{high}}, breakable, before skip=4pt, after skip=4pt]
\textcolor{gray}{\small Oral $n=31$ (meta 27/31) | HC $n=5$ (meta 5/5) | Reject $n=32$ (meta 32/32)}

\textbf{Sample note.} \textit{n\_oral=31, n\_reject=32 with 87\% and 100\% meta-coverage respectively support strong claims; HC sample is exactly n=5, so the oral\_hc\_gap should be read as suggestive rather than statistically firm.}

\textbf{Step-by-Step:}
\begin{enumerate}[leftmargin=14pt,topsep=2pt,itemsep=1pt]
  \item Pick the decomposition axis: timescale, locality, privacy boundary, sub-population, computational stage --- the axis must be intrinsic to the problem, not arbitrary.
  \item Partition P into sub-problems \{P\_i\} where each P\_i has a known or constructible solution method targeted at its specific structure.
  \item Specify the interface contract between stages: what each exposes, what invariants each preserves, how composition errors propagate.
  \item Solve each P\_i; for each, name the method and why it's appropriate at this granularity (not a generic solver).
  \item Aggregate the sub-solutions back to a solution for P with stated guarantees --- preservation theorem, ablation chain, or known invariant.
  \item Demonstrate via ablation that the decomposition (not bundled tricks) is the binding factor in the observed gain.
\end{enumerate}
\textbf{Success conditions} (Oral):
\begin{itemize}[leftmargin=12pt,topsep=2pt,itemsep=2pt]
  \item \textbf{The decomposition axis is justified by an explicit independence or separation claim that is then proved or measured (timescale separation, factor independence, multiplicative decomposability), not merely asserted.} \textcolor{gray}{\footnotesize evidence: \texttt{ICML\_2025\_1738}, \texttt{NeurIPS\_2025\_1864}, \texttt{NeurIPS\_2025\_1873}, \texttt{NeurIPS\_2025\_1945}, \texttt{ICML\_2024\_1138}} \\ Oral papers turn 'decompose' from a design intuition into a load-bearing technical claim: 1738 proves loss factors multiplicatively into LR-anneal and distribution-shift components and fits each separately; 1864 proves overfitting and generalization are separated by an analytically computed timescale gap; 1873 and 1138 use timescale separation to make a coupled nonconvex system tractable; 1945 derives $\tau$\_gen=O(1) vs $\tau$\_mem=Ω(n). The independence is not assumed --- it is the central theorem.
  \item \textbf{Each sub-problem is solvable with strictly simpler / pre-existing machinery, and the recombination operator is explicit and lossless on the property being claimed (accuracy, identity, convergence rate).} \textcolor{gray}{\footnotesize evidence: \texttt{ICLR\_2023\_0337}, \texttt{ICLR\_2024\_0826}, \texttt{NeurIPS\_2025\_1854}, \texttt{ICML\_2024\_1072}, \texttt{ICLR\_2025\_1603}} \\ 0337 splits multi-step reasoning into Selection and Inference, each within a 7B LM's single-step capability, and the loop is the recombination; 1854 splits long-motion generation into SPDM (segments) + TPDM (transitions) with a phase-state interface; 1072 factorizes speech into independent codec streams that recombine coherently at synthesis; 0826 uses execution states (not syntax) as the interface between sub-programs; 1603 retrieves clips then interpolates. The interface between stages is concrete and the gain over end-to-end is directly attributable to it.
  \item \textbf{The chosen axis exposes a *non-obvious* compositional structure --- semantic execution state instead of syntax, factorized latent instead of holistic signal, hierarchical/recursive instead of flat --- and the paper articulates why the prevailing axis was wrong.} \textcolor{gray}{\footnotesize evidence: \texttt{ICLR\_2024\_0826}, \texttt{ICLR\_2024\_0881}, \texttt{ICML\_2023\_0466}, \texttt{NeurIPS\_2024\_1328}, \texttt{NeurIPS\_2025\_1915}} \\ 0826 explicitly argues that execution states (semantic) decompose better than program syntax; 0881 argues theorem libraries should be mutable, not static; 0466 extends a flat automaton with a call-stack subroutine semantics; 1328 and 1915 decompose a generative signal into part-level streams (identity/motion or head/body/hand) whose independence is geometrically motivated. Reviewers reward the reframe of *what* to decompose along, not just the act of decomposing.
  \item \textbf{Decomposition is paired with a recombination guarantee --- convergence rate matched, global coherence preserved, or empirical equivalence-to-noise validated --- rather than relying on the parts being 'good enough' separately.} \textcolor{gray}{\footnotesize evidence: \texttt{ICML\_2023\_0413}, \texttt{ICLR\_2025\_1541}, \texttt{ICML\_2025\_1791}, \texttt{NeurIPS\_2025\_1830}} \\ 0413 proves Error Feedback in non-convex FL recovers full-precision rates with linear speedup despite being decomposed across heterogeneous clients; 1541 proves convergence with parameters depending only on system-level (not per-component) quantities; 1791 measures inter-model curve differences against intra-run seed noise to show the decomposition is exact, not just approximate; 1830 derives a multiscale equation that matches three observed phases. The 'recombine' step is held to a quantitative bar.
  \item \textbf{When the decomposition surfaces a meta-decision (when/where to spend compute), it is implemented as a separate dedicated module with its own training signal, not absorbed into the primary model.} \textcolor{gray}{\footnotesize evidence: \texttt{ICLR\_2023\_0272}, \texttt{ICLR\_2025\_1505}, \texttt{ICLR\_2024\_0922}} \\ 0272 adds a dedicated mask-predictor that tells the symbolic reasoner where to act; 1505 adds an introspection network separate from the prediction network for early-exit decisions; 0922 separates hypothesis generation from hypothesis testing/refinement. Decomposition becomes architectural --- the meta-component has a distinct training objective --- rather than a prompt-level trick.
\end{itemize}
\textbf{Failure modes} (Reject):
\begin{itemize}[leftmargin=12pt,topsep=2pt,itemsep=2pt]
  \item \textbf{Decomposition is proposed but ablations do not isolate which component drives the gain, leaving open whether a simpler baseline matches it.} \textcolor{gray}{\footnotesize evidence: \texttt{ICLR\_2023\_0234}, \texttt{ICLR\_2024\_0814}, \texttt{ICLR\_2025\_1605}, \texttt{ICLR\_2025\_1514}} \\ Meta-reviews flag this directly: 0234 --- random selection performs within one std of the proposed method on final accuracy despite the elaborate decomposition; 0814 --- selection procedure is not ablated independently from the single-objective oracle; 1605 --- clustering and prompt-init ablations missing; 1514 --- modular planner/executor not ablated, accuracy gains do not justify 6$\times$ compute. The paper says 'we decompose'; reviewers ask 'and what part of the decomposition is doing the work?'
  \item \textbf{Two known techniques are stitched together along a plausible decomposition axis without showing the axis addresses the binding constraint.} \textcolor{gray}{\footnotesize evidence: \texttt{NeurIPS\_2024\_1160}, \texttt{NeurIPS\_2024\_1174}, \texttt{ICLR\_2025\_1429}} \\ 1160 layers blockchain + FL + anomaly detection without showing which layer solves the cited problem; 1174 conditions a GAN on a structural embedding for inertial signals but reviewers cite limited novelty over standard augmentation; without convincing motivation; 1429 reformulates centralized ML-IRL into federated form but reviewers find limited differentiation from prior work. The decomposition is engineering composition, not a structural insight that resolves a specific obstruction.
  \item \textbf{Distributed/federated reformulation that assumes away the hard case of the decomposition rather than confronting it.} \textcolor{gray}{\footnotesize evidence: \texttt{ICLR\_2024\_0960}, \texttt{ICLR\_2023\_0310}, \texttt{NeurIPS\_2023\_0634}} \\ 0960's Byzantine-resilience theorem holds only when Byzantine clients form a small fraction over all iterations --- the AC says 'the really interesting case is the one where Byzantine clients can form a majority in some iterations [...] this difficult case is not studied'; 0310's bilevel formulation is criticized for shaky basic assumptions; 0634 sits in a crowded space without isolating its contribution. The decomposition is technically correct but pre-empts the hardness it claims to address.
  \item \textbf{Decomposition is recast amortization or a known wrapper, presented as a new method rather than as the application of a known pattern.} \textcolor{gray}{\footnotesize evidence: \texttt{ICLR\_2023\_0271}, \texttt{ICLR\_2023\_0280}, \texttt{NeurIPS\_2024\_1311}} \\ 0271 (predict params from datasets), 0280 (Meta OT --- predict warm-start), and 1311 (predict quantum-circuit params) are each reviewed as plausible amortization frameworks but rejected for lacking either compelling applications, scope, or differentiation from concurrent work. When decomposition reduces to 'train an offline predictor for the inner loop,' reviewers want either a new theoretical guarantee or a use case where the pattern is uniquely enabling.
  \item \textbf{Decomposition adds complexity without a clear net win on the metric that motivated it.} \textcolor{gray}{\footnotesize evidence: \texttt{NeurIPS\_2023\_0635}, \texttt{NeurIPS\_2023\_0666}, \texttt{ICML\_2025\_1812}, \texttt{ICLR\_2024\_0902}} \\ 0635 --- 'idea of local and global knowledge distillation significantly complicates the federated learning setting' with too many tunable variables; 0666 --- disentangling lighting/texture is conceptually clean but quantitative gains are marginal and qualitative results inconsistent; 1812 --- manual transition-token selection makes the decomposition fragile and narrow; 0902 --- multi-objective evolutionary SR cannot demonstrate advantage over baselines. The decomposition is real but its overhead is not amortized by the result.
  \item \textbf{Multi-objective / Pareto reformulations that replace a scalar weighting with a richer structure but cannot show the new structure unlocks a regime the old one missed.} \textcolor{gray}{\footnotesize evidence: \texttt{ICLR\_2025\_1424}, \texttt{ICLR\_2025\_1356}, \texttt{ICLR\_2024\_0814}} \\ 1424 reframes IB trade-offs as multi-objective optimization, but reviewers cite small datasets and unproven stability --- the non-convexity-of-frontier claim is not isolated as the limiting factor; 1356 combines MTO statistics with linear scalarization but reviewers find the contribution limited; 0814 provably traces the Pareto front but is compared to only one baseline. Replacing scalar with vector objectives is structurally appealing, but reviewers demand evidence the scalar form was demonstrably the bottleneck.
\end{itemize}
\textbf{Oral vs Reject gap.} Oral papers earn the decomposition by promoting it from design intuition to a load-bearing technical claim: they either prove the parts are independent (multiplicative loss in 1738; timescale gap in 1864/1945; mean-field PDE separation in 1138/1830) or measure inter-component coupling against a noise floor (1791). Reject papers assert the decomposition and ablate the *whole pipeline*, leaving reviewers unable to tell whether the chosen axis (per-client, per-objective, planner/executor, blockchain+FL+anomaly) is the binding constraint or just one of many possible carvings --- 0234, 0814, 1514, 1605, and 1424 are all flagged for this exact gap. Second, Oral papers specify the recombination operator with a quantitative property it must preserve (rate, identity, coherence-within-noise), while Reject papers stitch components with heuristic glue and report end-task numbers --- when those numbers are marginal (0666, 0902) or the cost balloons (1514: 6$\times$ slower), the decomposition has nothing to fall back on. Third, Oral papers articulate *why the prevailing axis was wrong* (semantic vs syntactic in 0826; mutable vs static library in 0881; recursive vs flat formalism in 0466), whereas Reject papers decompose along the most obvious axis without engaging the cognitive barrier they had to break.

\textbf{Oral vs HC gap.} HC sample (n=5) is at the threshold; reading them as a contrast: HC papers (0624 EmbodiedGPT, 0656 HuggingGPT, 0677 LLM+PDDL, 0756 VideoComposer, 0942 SEINE) use decomposition as an *engineering integration* pattern --- compose pre-existing systems, route subtasks to specialists, surface intermediate latents as conditioning --- and meta-reviews praise utility and breadth while flagging "limited technical novelty" or "concept previously explored." Oral papers, by contrast, use decomposition to expose *new analytical structure* (timescale separation, factorizable training dynamics, recursive automata semantics) that yields a theorem or a measurement at the noise floor. The HC bar is "this composition works and is useful"; the Oral bar is "this decomposition reveals a previously invisible structural property of the system."

\textbf{Reviewer expectations:}
\begin{itemize}[leftmargin=12pt,topsep=2pt,itemsep=1pt]
  \item \textit{[reject\_reviews]} Show, with an ablation that holds everything else fixed, which component of the decomposition is responsible for the gain --- random or simpler baselines must be ruled out, not just outperformed in aggregate.
  \item \textit{[reject\_reviews]} When the decomposition is a federated/distributed reformulation, do not assume away the hard case (e.g., majority-Byzantine rounds, full participation) the original problem made hard; address it head-on.
  \item \textit{[oral\_reviews]} Provide a quantitative recombination guarantee --- convergence rate matched, generalization gap derived, curves collapsed within seed noise --- not just qualitative coherence.
  \item \textit{[both]} When introducing modular pipelines (planner+executor, mask+reasoner, codec+diffusion), demonstrate compute/latency cost is justified by the accuracy or controllability gain; reviewers explicitly weigh this trade-off.
  \item \textit{[oral\_reviews]} Articulate why the natural axis of decomposition is the *right* one --- what hidden structure (timescale, factor, semantic state) the chosen axis exposes that the obvious axis would obscure.
  \item \textit{[reject\_reviews]} If the decomposition reduces to amortization or to a wrapper around an existing solver, provide either a new theoretical property or a setting where the wrapper is uniquely enabling --- otherwise reviewers call it incremental.
\end{itemize}
\textbf{Cognitive barriers:}
\begin{itemize}[leftmargin=12pt,topsep=2pt,itemsep=1pt]
  \item Treating a process as monolithic by default --- reasoning as a single capability (0337), symbolic constraints as globally enforced (0272), libraries as static expert artifacts (0881), face dynamics as a single high-dimensional signal (1328) --- blocks the move to view the process as composed of separable sub-operations.
  \item Assuming coupled nonconvex dynamics are intractable end-of-story rather than searching for a different analysis frame (mean-field, timescale separation, multiscale) that makes them decomposable --- the cognitive cost of switching mathematical machinery (1138, 1864, 1830) is what stops the decomposition from being attempted.
  \item Confusing the natural unit with the right unit: program syntax instead of execution state (0826), holistic motion instead of part-level streams (1915), absolute outputs instead of ordinal rankings --- practitioners default to the most visible decomposition axis even when a less obvious axis carries the actual independence structure.
  \item Believing that joint modeling is needed for global coherence (factorized synthesis 'risks inconsistencies' --- 1072; selection/inference 'must be end-to-end' --- 0337), so the burden of proof feels asymmetrically high on the decomposer to show recombination is lossless, even when the joint formulation has no such evidence in its favor.
\end{itemize}
\textbf{Representative examples:}
\begin{itemize}[leftmargin=12pt,topsep=2pt,itemsep=1pt]
  \item \textbf{[Oral]} \texttt{ICML\_2025\_1738} \textemdash\ Proves the continual-pretraining loss decomposes *multiplicatively* into a distribution-shift factor and an LR-annealing factor, fits each independently, and recombines them into a predictive law that holds across hyperparameter settings --- the cleanest example of 'independence is a theorem, not a hope'.
  \\ \textcolor{gray}{\footnotesize \textit{Lesson:} Independence is a theorem, not a hope --- proving the loss decomposes multiplicatively into independent factors lets each be fit separately and recombined into a predictive law.}
  \item \textbf{[Oral]} \texttt{ICLR\_2024\_0826} \textemdash\ ExeDec argues the right unit of decomposition is *execution state* (what each sub-program produces), not program syntax --- a deliberate reframe of the decomposition axis that exposes the compositional generalization gap others were missing.
  \\ \textcolor{gray}{\footnotesize \textit{Lesson:} The decomposition AXIS can itself be the contribution --- choosing execution state over program syntax exposes a compositional gap others missed entirely.}
  \item \textbf{[Oral]} \texttt{NeurIPS\_2025\_1854} \textemdash\ Splits long-motion generation into SPDM (semantic clips) + TPDM (transitions) with a phase-space interface borrowed from physics-based animation --- each sub-problem becomes independently tractable and the recombination is geometrically meaningful, not heuristic stitching.
  \\ \textcolor{gray}{\footnotesize \textit{Lesson:} Borrow interface contracts from neighboring fields (here: phase-space animation) --- geometric interfaces between stages produce meaningful recombination, not heuristic stitching.}
  \item \textbf{[Oral]} \texttt{NeurIPS\_2025\_1945} \textemdash\ Derives two distinct training timescales ($\tau$\_gen=O(1), $\tau$\_mem=Ω(n)) for diffusion models and proves the gap grows with dataset size --- turning early stopping from a heuristic into an emergent property of a decomposable training dynamic.
  \\ \textcolor{gray}{\footnotesize \textit{Lesson:} When the decomposition has emergent structure ($\tau$\_gen vs $\tau$\_mem with provable gap), heuristics like early stopping become first-principles consequences.}
  \item \textbf{[Reject]} \texttt{ICLR\_2025\_1514} \textemdash\ MSR-ViR decomposes VideoQA into planner + grounded executors but the meta-review notes the gain does not justify a 6$\times$ compute increase and the decomposition's per-module contribution is not isolated --- a textbook 'modularity without ablations' failure.
  \\ \textcolor{gray}{\footnotesize \textit{Lesson:} A decomposition that demands 6$\times$ compute must demonstrate per-module contribution via ablation --- without isolating which module pays for the cost, modular gain looks like a brute-force win.}
  \item \textbf{[Reject]} \texttt{ICLR\_2024\_0960} \textemdash\ Extends Byzantine-resilient FL to partial participation but, per the AC, 'assumes away' the hard case where Byzantine clients can form a majority in some rounds --- the decomposition is technically correct but evades the regime that made the problem interesting.
  \\ \textcolor{gray}{\footnotesize \textit{Lesson:} Decompositions that 'assume away' the hardest regime (Byzantine majority) are technically correct but evade the case that made the problem interesting.}
  \item \textbf{[Reject]} \texttt{ICLR\_2023\_0234} \textemdash\ Proposes a sophisticated ranking-preserving condensation surrogate, but the meta-review observes random selection performs within one std on final accuracy --- exemplifying decomposition that cannot demonstrate it is the binding factor in observed gains.
  \\ \textcolor{gray}{\footnotesize \textit{Lesson:} If random selection performs within one std of the proposed decomposition, the decomposition is not the binding factor in the reported gains --- ablation against the trivial baseline is mandatory.}
\end{itemize}
\end{tcolorbox}

\begin{tcolorbox}[colback=bgskill, colframe=cblue!70, title={\textbf{Diagnose Then Surgically Fix} \textcolor{gray}{\small\texttt{root\_cause\_surgical\_remediation}} \hfill confidence: \textbf{high}}, breakable, before skip=4pt, after skip=4pt]
\textcolor{gray}{\small Oral $n=28$ (meta 19/28) | HC $n=40$ (meta 30/40) | Reject $n=27$ (meta 26/27)}

\textbf{Sample note.} \textit{Sample sizes are healthy (n\_oral=28, n\_reject=27, n\_hc=40) and meta-review coverage is strong on the Reject side (96\%) and adequate on the Oral side (68\%); a few Oral papers (e.g., several ICML 2023 Orals) lack meta-reviews so reviewer-voice claims lean on the ICLR/NeurIPS subset.}

\textbf{Step-by-Step:}
\begin{enumerate}[leftmargin=14pt,topsep=2pt,itemsep=1pt]
  \item Observe a specific, reproducible failure F: characterize precisely which input regime, metric, or behavior collapses.
  \item Isolate the causal component C via analysis or controlled ablation; formal analysis (a theorem about which step introduces the pathology) beats empirical localization.
  \item Propose a minimal edit E targeting only C --- one line of the loss, one layer, one training-phase rule, one estimator step.
  \item Prove or empirically verify a preservation guarantee: everything that worked before continues to work; the edit does not propagate side effects.
  \item Validate that F is resolved by E specifically, not by capacity / data / optimization confounds --- ablate the edit against close alternatives.
  \item Position the contribution as the formal localization, not the patch --- the edit demonstrates the localization, but the localization is the result.
\end{enumerate}
\textbf{Success conditions} (Oral):
\begin{itemize}[leftmargin=12pt,topsep=2pt,itemsep=2pt]
  \item \textbf{The failure mechanism is formalized as a measurable quantity at a specific internal site (a loss term, a numerical constant, an estimator's bias) before any fix is proposed --- diagnosis precedes and constrains the intervention.} \textcolor{gray}{\footnotesize evidence: \texttt{ICLR\_2024\_0884}, \texttt{ICLR\_2024\_0855}, \texttt{ICML\_2024\_1011}, \texttt{ICLR\_2025\_1583}, \texttt{NeurIPS\_2025\_1905}} \\ These Orals each produce a concrete object that pinpoints the pathology --- Lipschitz singularity, MSE sensitivity in high-perturbation regions, Hessian/gradient-conflict patterns, gradient-statistic imbalance, an analytic decomposition of label-smoothing --- and only then design the edit. Reviewer voice praises this as 'theoretical analysis of the limitations in previous approaches.'
  \item \textbf{The intervention is paired with a formal preservation guarantee (unbiasedness, equivalence under a degenerate parameter, convergence proof) that bounds collateral damage from the edit.} \textcolor{gray}{\footnotesize evidence: \texttt{ICLR\_2023\_0282}, \texttt{ICLR\_2024\_0860}, \texttt{ICML\_2025\_1693}} \\ Each pairs the surgical fix with a theorem: gradient-bias correction provably yields unbiased combined direction; InfoBatch's rescaling makes the gradient estimator unbiased in expectation; the non-parametric cost memory is bias-free by construction. The proof is what distinguishes 'surgical' from 'plausible patch.'
  \item \textbf{The fix is targeted at a single named component (one loss term, one normalization layer, one schedule, one numerical operation) rather than a broad architectural redesign --- and its scope is explicitly bounded.} \textcolor{gray}{\footnotesize evidence: \texttt{ICLR\_2021\_0046}, \texttt{ICLR\_2023\_0376}, \texttt{NeurIPS\_2025\_1905}, \texttt{ICLR\_2024\_0855}} \\ Binaural replaces only the loss; SAR fixes only BN behavior plus a flat-minima regularizer; MaxSup edits only the penalized logit; iCT swaps only the metric and the schedule. The minimality of the edit is what makes 'unrelated behavior preserved' verifiable.
  \item \textbf{The fix simultaneously eliminates multiple symptoms or accumulated patches that the field had stacked around the same root cause, demonstrating that the diagnosis is structurally upstream of those symptoms.} \textcolor{gray}{\footnotesize evidence: \texttt{ICML\_2023\_0465}, \texttt{ICLR\_2024\_0832}, \texttt{ICLR\_2024\_0860}} \\ Hiera shows that fixing the learning objective lets a stack of architectural 'bells and whistles' be deleted; FSQ shows that one reformulation eliminates the entire heuristic stack around codebook collapse; InfoBatch's bias fix also removes the need for elaborate pruning-criterion machinery. This 'patch-collapse' signal is a strong reviewer marker for genuine root-cause.
  \item \textbf{The diagnosis is validated by directly measuring the targeted property after the fix (e.g., empirical Lipschitz constant, gradient-statistic distributions, covariance spectra) and not only by aggregate task metrics.} \textcolor{gray}{\footnotesize evidence: \texttt{ICLR\_2024\_0884}, \texttt{ICLR\_2025\_1583}, \texttt{ICML\_2024\_1011}} \\ Each paper presents a 'before/after' on the specific internal quantity it claimed was the failure source --- this closes the diagnostic loop and is precisely what Reject papers typically omit.
\end{itemize}
\textbf{Failure modes} (Reject):
\begin{itemize}[leftmargin=12pt,topsep=2pt,itemsep=2pt]
  \item \textbf{Diagnosis is plausible but no controlled ablation isolates the proposed fix as the dominant lever --- reviewers cannot rule out that data quality, capacity, or a peripheral factor is the real driver.} \textcolor{gray}{\footnotesize evidence: \texttt{ICLR\_2025\_1345}, \texttt{ICLR\_2025\_1369}, \texttt{ICLR\_2025\_1584}, \texttt{ICLR\_2025\_1604}} \\ Multiple Reject meta-reviews explicitly say the structural-mismatch claim 'appears to be a theoretical impurity rather than a practical bottleneck' and that without isolating ablations the targeted correction's causal role is unverified.
  \item \textbf{Targeted fix is presented but its theoretical claims are internally inconsistent or formally wrong, undercutting the 'surgical correctness' that this methodology depends on.} \textcolor{gray}{\footnotesize evidence: \texttt{ICLR\_2025\_1584}, \texttt{ICLR\_2025\_1604}, \texttt{ICML\_2025\_1794}} \\ MoE Smoothness assumes Lipschitz continuity for known-discontinuous Top-K routers; the DP-GNN privacy analysis is judged simply incorrect; SemiDICE diagnosis is right but reviewers note prior papers already established the mechanism --- a surgical fix only lands if the proof holds and is novel.
  \item \textbf{Fix is validated only on toy or narrow regimes, so the claim that the localized mechanism is 'the' root cause cannot be defended at scale or across distributions.} \textcolor{gray}{\footnotesize evidence: \texttt{ICLR\_2024\_0765}, \texttt{ICLR\_2024\_0864}, \texttt{ICLR\_2025\_1504}, \texttt{ICLR\_2025\_1415}} \\ Reject meta-reviews flag restricted dimensionality ($\le$4D for MetaGFN), narrow data-generating processes (Brownian/OU only for signature), or marginal gains on some datasets --- the surgical fix may be real but its scope of applicability is left unestablished.
  \item \textbf{Failure to position the diagnosis against concurrent or prior work that already identified the same mechanism, so the 'we localized the root cause' contribution collapses.} \textcolor{gray}{\footnotesize evidence: \texttt{NeurIPS\_2023\_0684}, \texttt{ICLR\_2024\_0894}, \texttt{ICML\_2025\_1794}, \texttt{ICLR\_2025\_1573}} \\ Reviewers explicitly cite uncited adjacent papers (Spectral Feature Augmentation; PlanE; Dual RL/ODICE; DeepGCN). Surgical-remediation papers are penalized hard when the localized mechanism is presented as new but is independently known.
  \item \textbf{Claimed mechanism is asserted but the structural property it targets is never directly measured --- reviewers see only aggregate accuracy and cannot confirm the diagnosis.} \textcolor{gray}{\footnotesize evidence: \texttt{ICLR\_2024\_0829}, \texttt{ICLR\_2025\_1564}, \texttt{ICLR\_2025\_1594}} \\ Feature-accentuation visualizations are not psychophysically validated; VQ restoration's discrete-bottleneck claim isn't probed with diagnostics that would expose hard-assignment failure; StoryAgent's cross-component consistency is not measured. The methodology demands 'F is resolved' be observable on the targeted variable, not just on aggregate metrics.
\end{itemize}
\textbf{Oral vs Reject gap.} Oral papers couple three moves that Reject papers almost always miss at least one of: (1) the failure mode is operationalized as a *measurable quantity* on the suspected internal site (gradient norm statistics, covariance log-determinant, Lipschitz constant near t=0, energy/phase decomposition of the loss) so the diagnosis is falsifiable; (2) the fix is paired with a *preservation guarantee* --- a theorem of unbiasedness, equivalence under degenerate parameters, or proof that unrelated behavior is held constant --- not just an empirical 'works as well or better'; (3) ablations explicitly isolate the surgical edit from confounding factors (capacity, data, schedule) so reviewers can see the targeted mechanism is the dominant lever. Reject papers typically supply (1) at the level of intuition only, skip (2) entirely, and offer aggregate-metric improvements without ablations that rule out alternative explanations --- which is exactly the language Reject meta-reviews use ('without controlled ablations\ldots{} it remains unclear whether the targeted metric is the dominant bottleneck').

\textbf{Oral vs HC gap.} Both Oral and HC papers execute the diagnose$\to$localize$\to$minimal-fix loop, but Oral papers more consistently anchor the localization in a *formal* object (a closed-form bias term, a proven Lipschitz singularity, a loss-landscape Hessian/conflict analysis, an analytic decomposition of a loss into named summands) and then prove that the fix preserves a precise property --- unbiasedness in expectation, equivalence to full-gradient updates, convergence guarantees, or mathematical equivalence to a degenerate special case. HC papers more often perform the diagnosis empirically (ablation tables, controlled component swaps, measurement of internal statistics) and validate the fix by downstream metric recovery rather than by a theorem. The other Oral-distinctive move is breadth-of-recovery: Oral fixes are shown to simultaneously eliminate multiple downstream symptoms or accumulated patches (Hiera removes 'bells and whistles'; FSQ removes a stack of VQ heuristics; MaxSup fixes representation collapse and miscalibration with one swap), whereas HC papers typically resolve the single targeted symptom.

\textbf{Reviewer expectations:}
\begin{itemize}[leftmargin=12pt,topsep=2pt,itemsep=1pt]
  \item \textit{[oral\_reviews]} A formal characterization of the failure mechanism --- a theorem, a closed-form bias expression, a proven singularity, or a decomposition of the loss into named terms --- not just an intuitive narrative.
  \item \textit{[oral\_reviews]} A preservation guarantee that the surgical fix leaves untouched behavior actually untouched (unbiasedness, equivalence to a known good case, convergence, no quality drop on adjacent tasks).
  \item \textit{[reject\_reviews]} Ablations that isolate the proposed edit from capacity, data-quality, and hyperparameter confounds --- Reject meta-reviews repeatedly demand this and reject when it is missing.
  \item \textit{[reject\_reviews]} Explicit comparison against concurrent or prior work that may have already identified the same mechanism; failure to cite is treated as evidence the diagnosis is not novel.
  \item \textit{[both]} Direct measurement of the targeted internal property (the 'F' in the operational signature) before and after the fix, not just downstream task metrics --- both Orals supply this and Rejects are penalized for omitting it.
  \item \textit{[reject\_reviews]} Demonstration that the localized mechanism is the *dominant* limiting factor, not a peripheral one --- reviewers consistently challenge whether the surgically targeted component is actually load-bearing in the regime of interest.
\end{itemize}
\textbf{Cognitive barriers:}
\begin{itemize}[leftmargin=12pt,topsep=2pt,itemsep=1pt]
  \item Misattribution of root cause: practitioners default to blaming optimizer hyperparameters, capacity, or data quality, when the actual culprit is a specific term in a loss, a numerical singularity, or a normalization layer treated as 'neutral preprocessing' (e.g., BN in TTA, L2 in binaural, semi-gradient in DICE).
  \item Hidden failure modes that look like mild suboptimality rather than bugs: methods 'still converge, just suboptimally,' so the bias never crosses the visibility threshold that would prompt diagnosis (gradient bias in multi-objective, info-loss vs bias confusion in pruning).
  \item Component is treated as a fixed canonical primitive whose internals are off-limits for surgery; field-wide convention to patch around the component rather than open it up (label smoothing, MSE in CM training, fixed timestep schedule near singularities).
  \item Output-level vs mechanism-level framing: practitioners measure the failure at the output and propose architectural or data-level remedies, missing that a single internal site fully explains the symptom when probed with the right mathematical lens (loss-landscape geometry, Lipschitz analysis, covariance of hidden states).
\end{itemize}
\textbf{Representative examples:}
\begin{itemize}[leftmargin=12pt,topsep=2pt,itemsep=1pt]
  \item \textbf{[Oral]} \texttt{ICLR\_2024\_0884} \textemdash\ Proves a Lipschitz singularity at t$\to$0 as a theorem, measures it across architectures, then applies timestep truncation that removes only the singular region --- a textbook diagnose$\to$formalize$\to$minimally-edit chain.
  \\ \textcolor{gray}{\footnotesize \textit{Lesson:} A textbook diagnose $\to$ formalize as theorem $\to$ minimally edit chain. The truncation removes only the singular region; everything else is untouched.}
  \item \textbf{[Oral]} \texttt{ICLR\_2024\_0860} \textemdash\ Identifies pruning-induced gradient bias (not information loss) as the failure mechanism and corrects it with a per-batch rescaling whose unbiasedness is proven in expectation, leaving the rest of training untouched.
  \\ \textcolor{gray}{\footnotesize \textit{Lesson:} Reframing the failure as gradient bias (not information loss) localizes the cause to a correctable expectation --- per-batch rescaling fixes it with a proof, not a heuristic.}
  \item \textbf{[Oral]} \texttt{ICLR\_2023\_0282} \textemdash\ Reorders nonlinear-aggregation and stochastic-estimation steps so the resulting combined gradient is provably unbiased, with no extra sample-complexity cost --- surgical move at exactly the algebraic site that introduced the bias.
  \\ \textcolor{gray}{\footnotesize \textit{Lesson:} The surgical move is at the algebraic site that introduced the bias --- reordering nonlinear-aggregation and stochastic-estimation produces an unbiased gradient with no extra sample-complexity cost.}
  \item \textbf{[Oral]} \texttt{NeurIPS\_2025\_1905} \textemdash\ Analytically decomposes label smoothing to isolate the ground-truth-logit penalty as the term causing representation collapse, then surgically replaces only that term with a max-logit penalty.
  \\ \textcolor{gray}{\footnotesize \textit{Lesson:} Analytically decompose a multi-term loss to isolate the offending term; then surgically replace only that term --- preservation by construction.}
  \item \textbf{[Reject]} \texttt{ICLR\_2025\_1584} \textemdash\ Targets a real mechanism (sparsity vs stability in MoE routing) but the central theorem assumes Lipschitz continuity for routers known to be discontinuous, so the formal localization step that this methodology requires collapses.
  \\ \textcolor{gray}{\footnotesize \textit{Lesson:} If the formal localization step assumes a property the system doesn't have (Lipschitz on a discontinuous router), the central theorem collapses regardless of the empirical fix.}
  \item \textbf{[Reject]} \texttt{ICML\_2025\_1794} \textemdash\ Correctly diagnoses that semi-gradient DICE estimates a different fixed point than full-gradient DICE, but fails to position the diagnosis against concurrent ICLR 2024 papers that established the same mechanism, undercutting the 'we localized it' contribution.
  \\ \textcolor{gray}{\footnotesize \textit{Lesson:} The right diagnosis fails when concurrent papers established the same mechanism --- surgical remediation without explicit positioning against close prior work erases the localization claim.}
  \item \textbf{[Reject]} \texttt{ICLR\_2025\_1369} \textemdash\ Proposes a structure-aware substitution justified as a relaxed-assumption fix, but provides no ablations isolating the substituted component, leaving reviewers unable to confirm the structural mismatch is the dominant bottleneck.
  \\ \textcolor{gray}{\footnotesize \textit{Lesson:} Without ablations isolating the substituted component, reviewers cannot confirm the structural mismatch was the dominant bottleneck --- surgical attribution requires per-component evidence.}
\end{itemize}
\end{tcolorbox}

\begin{tcolorbox}[colback=bgskill, colframe=cblue!70, title={\textbf{Soft Probabilistic Relaxation} \textcolor{gray}{\small\texttt{soft\_probabilistic\_relaxation}} \hfill confidence: \textbf{low}}, breakable, before skip=4pt, after skip=4pt]
\textcolor{gray}{\small Oral $n=3$ (meta 2/3) | HC $n=0$ (meta 0/0) | Reject $n=12$ (meta 12/12)}

\textbf{Sample note.} \textit{n\_oral=3 with only 2/3 meta-reviews and n\_hc=0, so the oral/HC contrast is unavailable and success-condition claims rest on a small Oral sample; the Reject side is well-covered (12/12 meta-reviews).}

\textbf{Step-by-Step:}
\begin{enumerate}[leftmargin=14pt,topsep=2pt,itemsep=1pt]
  \item Identify the hard rule, discrete decision, or non-differentiable step H in the existing pipeline that creates instability, information loss, or brittleness.
  \item Construct a continuous relaxation H\_$\lambda$ --- temperature-scaled softmax, distributional surrogate, penalized continuous decision --- with $\lambda$ controlling relaxation strength.
  \item Show the relaxation recovers a structural property H discards: calibration, optimization landscape smoothness, identifiability, or derivable distributional structure.
  \item Optimize the relaxed objective; ensure it reduces to H exactly when $\lambda$ $\to$ limit (strict relaxation, not approximation).
  \item Recover a discrete solution or dual guarantee from the relaxed solution --- rounding scheme, MAP, or theorem bridging the relaxed objective back to the original.
  \item Quantify the relaxation's gain with significance testing OR derive a structural advantage --- stability claims must be empirically grounded.
\end{enumerate}
\textbf{Success conditions} (Oral):
\begin{itemize}[leftmargin=12pt,topsep=2pt,itemsep=2pt]
  \item \textbf{Modify the objective/likelihood at the mathematical level rather than introducing the softening as a preprocessing step or an add-on regularizer. The relaxation is derived inside the loss, not stapled on top.} \textcolor{gray}{\footnotesize evidence: \texttt{ICML\_2025\_1792}, \texttt{ICML\_2024\_1008}, \texttt{ICLR\_2025\_1540}} \\ ICML\_2025\_1792 explicitly rejects imputation and derives a new score-matching objective over observed-only coordinates; ICML\_2024\_1008 replaces argmax pseudolabeling with the partial-label training objective itself; ICLR\_2025\_1540 encodes workload control inside a mixture prior rather than as a post-hoc constraint. The softening is the objective, not a penalty applied to an unchanged one.
  \item \textbf{Show that the relaxation reduces to the standard hard/complete-data method when the slack parameter degenerates (no missing data, no uncertainty, single expert).} \textcolor{gray}{\footnotesize evidence: \texttt{ICML\_2025\_1792}, \texttt{ICLR\_2025\_1540}} \\ ICML\_2025\_1792's meta-review highlights that 'when no values are missing, this reduces to the standard expected score-matching loss' as a strength; ICLR\_2025\_1540's mixture formulation recovers standard L2D when all annotators cover all examples. This reduction property is what lets reviewers trust that the relaxation is a strict generalization rather than a parallel method.
  \item \textbf{Import a mature theoretical framework from an adjacent subfield by identifying a structural isomorphism, rather than inventing a bespoke relaxation.} \textcolor{gray}{\footnotesize evidence: \texttt{ICML\_2024\_1008}, \texttt{ICML\_2025\_1792}, \texttt{ICLR\_2025\_1540}} \\ ICML\_2024\_1008 identifies that a CLIP top-k distribution is structurally identical to a partial-label set and transplants that literature wholesale; ICML\_2025\_1792 carries score-matching tricks (IS, VI variants) into the missing-data regime; ICLR\_2025\_1540 casts L2D as a latent-variable mixture so EM applies directly. Oral papers inherit proofs; they do not rederive them.
  \item \textbf{Carry concrete theoretical guarantees (finite-sample bounds, consistency, identifiability) through the relaxation, not only empirical ablations.} \textcolor{gray}{\footnotesize evidence: \texttt{ICML\_2025\_1792}, \texttt{ICML\_2024\_1008}} \\ ICML\_2025\_1792 ships an importance-weighted version with finite-sample guarantees alongside the VI scalable version; ICML\_2024\_1008 argues the partial-label theoretical tools apply directly to the pseudolabel case. Oral meta-reviews still probe for more (asymptotic normality in 1792), signaling reviewers expect the relaxation to be *analyzable*, not only trainable.
  \item \textbf{Deliver a second, orthogonal capability that the hard version structurally cannot offer (e.g., dual regime, workload/constraint control, scalable variant).} \textcolor{gray}{\footnotesize evidence: \texttt{ICLR\_2025\_1540}, \texttt{ICML\_2025\_1792}} \\ ICLR\_2025\_1540 wins reviewer enthusiasm specifically for workload distribution control, which falls out of the mixture priors --- not the main goal but a structural bonus. ICML\_2025\_1792 offers both a finite-sample IS variant and a VI variant for high dimensions. The softening buys something beyond 'handles the edge case.'
\end{itemize}
\textbf{Failure modes} (Reject):
\begin{itemize}[leftmargin=12pt,topsep=2pt,itemsep=2pt]
  \item \textbf{Derivations contain undefined variables, ungrounded probability factorizations, or missing intermediate steps, so reviewers cannot verify the relaxation is valid.} \textcolor{gray}{\footnotesize evidence: \texttt{ICLR\_2023\_0398}, \texttt{ICLR\_2024\_0853}} \\ ICLR\_2023\_0398's meta-review lists undefined latent variables u\_i/z\_i, missing variational distributions, and unclear calibration metrics; ICLR\_2024\_0853's lists broken EM derivations where log P(X;$\theta$) is dropped without justification and Y is silently conflated with a single label. Because the paper's selling point *is* the principled relaxation, opaque math is fatal --- not cosmetic.
  \item \textbf{Positions the contribution as a unification of several imperfect-supervision settings but never works out per-setting specifics, leaving reviewers with a general template and no concrete algorithm.} \textcolor{gray}{\footnotesize evidence: \texttt{ICLR\_2024\_0853}, \texttt{ICLR\_2024\_0911}} \\ ICLR\_2024\_0853's meta-review notes 'only a very general description of the approach without giving details of the final methods'; ICLR\_2024\_0911 gets flagged for not distinguishing itself from prior probabilistic NMF (Tan \& Févotte). A unification claim raises the bar: each subsumed setting must yield a worked-out instantiation, otherwise the framework reads as relabeling.
  \item \textbf{Combines two validated components (e.g., learned prior + existing optimization, score prior + VAE, EBM + diffusion) where the interaction is not demonstrably non-trivial --- reviewers read it as additive and incremental.} \textcolor{gray}{\footnotesize evidence: \texttt{ICLR\_2025\_1518}, \texttt{ICLR\_2025\_1451}, \texttt{NeurIPS\_2023\_0737}} \\ ICLR\_2025\_1518's own abstract concedes constituent components 'had independent prior validation, making their combination appear straightforwardly additive'; ICLR\_2025\_1451 lacks justification for why score-based priors help beyond Gaussian; NeurIPS\_2023\_0737 is rejected for 'incremental innovation' despite strong empirical results. Reviewers demand either a theoretical coupling argument or ablations that isolate interaction terms.
  \item \textbf{Claims stability/benefit from the soft formulation but the empirical edge over hard baselines is within noise and no statistical-significance test is reported.} \textcolor{gray}{\footnotesize evidence: \texttt{NeurIPS\_2024\_1280}, \texttt{ICLR\_2025\_1451}} \\ NeurIPS\_2024\_1280's meta-review cites 'marginal results and the lack of statistical significance testing' as the primary rejection reason, and the authors explicitly deferred broader evaluation. When the relaxation's value proposition is calibration/stability, reviewers need the magnitude of the effect quantified, not asserted.
  \item \textbf{Overlaps heavily with known cooperative/joint-training or probabilistic-factorization literature that the paper does not engage, so the softening reads as a rename.} \textcolor{gray}{\footnotesize evidence: \texttt{ICLR\_2023\_0340}, \texttt{ICLR\_2024\_0911}, \texttt{NeurIPS\_2023\_0643}} \\ ICLR\_2023\_0340 is rejected with four citations to earlier Cooperative GAN work; ICLR\_2024\_0911 for missing Tan \& Févotte; NeurIPS\_2023\_0643 for insufficient motivation relative to existing generative noisy-label methods. Soft probabilistic relaxations live in a crowded neighborhood, and a paper that does not map itself onto that neighborhood is read as re-discovering it.
  \item \textbf{Restricts evaluation to a single modality or a narrow benchmark when the relaxation is pitched as general, inviting 'does it transfer?' skepticism.} \textcolor{gray}{\footnotesize evidence: \texttt{ICLR\_2023\_0340}, \texttt{NeurIPS\_2023\_0643}, \texttt{ICLR\_2025\_1531}} \\ ICLR\_2023\_0340 shows NLP only and is asked for Vision; NeurIPS\_2023\_0643 is told Clothing1M alone is insufficient and asked for ImageNet/Food-101N; ICLR\_2025\_1531 gets hit for no comparison against self-supervised/primal-dual baselines. A principled relaxation implicitly promises generality; reviewers test that promise with cross-domain evidence.
\end{itemize}
\textbf{Oral vs Reject gap.} Oral papers place the soft relaxation *inside* the objective --- they rewrite the loss, derive a new estimator, or introduce a latent variable whose marginalization yields the method --- and they prove the new objective reduces to the standard hard/complete-data form in the degenerate limit (ICML\_2025\_1792, ICLR\_2025\_1540). Reject papers tend to bolt softness on as an additive regularizer, a prior swap, or a joint-training loop over unchanged components (ICLR\_2025\_1518, ICLR\_2025\_1451, ICLR\_2023\_0340, ICLR\_2023\_0398), which reviewers consistently label 'straightforwardly additive' or 'incremental.' Oral papers also import a mature framework with its theorems intact (partial-label theory, EM, score-matching identities) and inherit guarantees; Reject papers frequently unify or combine without carrying per-setting derivations through (ICLR\_2024\_0853, ICLR\_2024\_0911). Finally, Oral papers ship a *second* structural capability that the hard version cannot produce --- workload distribution control, dual IS/VI estimators --- rather than only matching the hard baseline more robustly; Reject papers pitched on 'stability' or 'calibration' alone struggle when the empirical delta is marginal and untested for significance (NeurIPS\_2024\_1280).

\textbf{Oral vs HC gap.} HC sample (n=0) too small for reliable contrast.

\textbf{Reviewer expectations:}
\begin{itemize}[leftmargin=12pt,topsep=2pt,itemsep=1pt]
  \item \textit{[oral\_reviews]} The relaxation must reduce cleanly to the standard hard-assignment/complete-data method when the slack parameter is zero, and that reduction should be stated explicitly in the paper.
  \item \textit{[oral\_reviews]} Theoretical properties of the estimator (consistency, finite-sample behavior, asymptotic normality) should be characterized, or explicitly flagged as future work with a principled reason.
  \item \textit{[reject\_reviews]} Every latent variable, variational distribution, and integral in the derivation must be defined, typed, and consistent across equations --- gaps block assessment regardless of empirical results.
  \item \textit{[reject\_reviews]} When combining two validated components (learned prior + optimization, score prior + VAE, EBM + diffusion), the paper must argue non-trivial interaction --- either theoretically or via ablations that isolate the coupling term.
  \item \textit{[reject\_reviews]} Close prior work on probabilistic/cooperative/unified variants of the same relaxation must be engaged directly, not just cited --- including quantitative comparison where baselines exist.
  \item \textit{[both]} When the relaxation is justified by stability or calibration, empirical claims require statistical-significance testing and multi-domain evaluation, not single-benchmark margins.
\end{itemize}
\textbf{Cognitive barriers:}
\begin{itemize}[leftmargin=12pt,topsep=2pt,itemsep=1pt]
  \item Hard commitments feel more informative than soft ones --- the field's intuition is that argmax concentrates useful signal, making it counterintuitive that structured uncertainty is itself the signal rather than noise to be suppressed (ICML\_2024\_1008).
  \item Incompleteness, ambiguity, and missing data are culturally treated as data-quality problems to fix via imputation or filtering, not as latent structure to model; reframing them as first-class objects requires resisting the preprocessing reflex (ICML\_2025\_1792, ICLR\_2025\_1540).
  \item The structural isomorphism between settings (e.g., top-k confidence $\leftrightarrow$ partial label set; missing-coordinate score matching $\leftrightarrow$ latent-variable score matching; workload distribution $\leftrightarrow$ mixture priors) sits across subfield boundaries and is invisible from within either community alone (ICML\_2024\_1008, ICLR\_2025\_1540).
  \item Competing-paradigm framings (discriminative vs generative, EBM vs diffusion, hard-assignment vs marginalization) discourage reinterpreting one paradigm's artifacts as another's --- e.g., a diffusion trajectory as an MCMC step --- even when the math permits it (NeurIPS\_2023\_0737, NeurIPS\_2023\_0643).
\end{itemize}
\textbf{Representative examples:}
\begin{itemize}[leftmargin=12pt,topsep=2pt,itemsep=1pt]
  \item \textbf{[Oral]} \texttt{ICML\_2024\_1008} \textemdash\ Most distinctive cross-subfield import: recognizes CLIP's top-k confidence distribution as structurally identical to a partial-label set and transplants the partial-label training framework wholesale rather than designing a new pseudolabel heuristic.
  \\ \textcolor{gray}{\footnotesize \textit{Lesson:} Cross-subfield import --- recognize a structural identity (CLIP confidence $\cong$ partial-label set), then transplant the partial-label training framework wholesale rather than designing a new pseudolabel heuristic.}
  \item \textbf{[Oral]} \texttt{ICML\_2025\_1792} \textemdash\ Cleanest objective-level relaxation: derives a new score-matching loss over observed-only coordinates that provably reduces to the standard loss when no data is missing, and ships both a finite-sample IS variant and a scalable VI variant.
  \\ \textcolor{gray}{\footnotesize \textit{Lesson:} The relaxation must reduce to the standard loss when no relaxation is needed (no missing data) --- provable degeneration to baseline is a strong indicator the relaxation is principled.}
  \item \textbf{[Oral]} \texttt{ICLR\_2025\_1540} \textemdash\ Best example of the relaxation yielding a structural bonus capability: casting expert assignment as a latent variable in a mixture produces workload distribution control as a free consequence of the mixture priors rather than a separately engineered constraint.
  \\ \textcolor{gray}{\footnotesize \textit{Lesson:} The cleanest relaxations yield free structural bonuses --- workload distribution control falls out of the mixture priors rather than requiring a separately engineered constraint.}
  \item \textbf{[Reject]} \texttt{NeurIPS\_2024\_1280} \textemdash\ Textbook attempt at the methodology --- a log-quadratic Potts relaxation with soft self-labeling --- undone by marginal empirical gains and no statistical-significance testing, showing that soft-relaxation pitched on stability must quantify the stability.
  \\ \textcolor{gray}{\footnotesize \textit{Lesson:} Soft relaxation pitched on stability must quantify the stability --- marginal empirical gains without significance testing convert the contribution to 'a heuristic among many'.}
  \item \textbf{[Reject]} \texttt{ICLR\_2024\_0853} \textemdash\ Unifies noisy/partial/multi-candidate labels as MLE over a latent label under EM but stops at the general template; broken per-setting derivations and unresolved symbol types illustrate how unification framings collapse without careful math.
  \\ \textcolor{gray}{\footnotesize \textit{Lesson:} Unification framings need careful per-setting math; broken derivations and unresolved symbol types make the framework collapse despite the appealing template.}
  \item \textbf{[Reject]} \texttt{ICLR\_2025\_1451} \textemdash\ Swaps a Gaussian prior for a score-based one inside a VAE --- classic two-component relaxation --- but does not demonstrate non-trivial interaction between the score prior and the structural regularizer, exemplifying the 'additive combination' failure.
  \\ \textcolor{gray}{\footnotesize \textit{Lesson:} Two-component relaxation requires demonstrating non-trivial interaction between the components --- additive combination without interaction tests reads as a wrapper.}
\end{itemize}
\end{tcolorbox}

\begin{tcolorbox}[colback=bgskill, colframe=cblue!70, title={\textbf{Exploit Redundancy via Compression} \textcolor{gray}{\small\texttt{redundancy\_compression\_and\_distillation}} \hfill confidence: \textbf{high}}, breakable, before skip=4pt, after skip=4pt]
\textcolor{gray}{\small Oral $n=31$ (meta 25/31) | HC $n=9$ (meta 7/9) | Reject $n=29$ (meta 29/29)}

\textbf{Sample note.} \textit{Sample is adequate (n\_oral=31, n\_reject=29, meta\_coverage 25/31 oral and 29/29 reject); HC sample (n=9) supports a tendency-level Oral-vs-HC contrast but not strong claims.}

\textbf{Step-by-Step:}
\begin{enumerate}[leftmargin=14pt,topsep=2pt,itemsep=1pt]
  \item Detect redundancy R in the system: where in the pipeline is compute / memory / parameter budget being spent on structurally removable computation? Localize to a layer / module / phase.
  \item Classify R as structural (theoretically removable without loss --- low-rank, sparse, separable) vs. empirical (typically noisy, removable with quality cost).
  \item Design a compact proxy P that captures R's essential outputs --- orthogonal-decomposition (base + adapter), per-layer reconstruction, null-space projection, entanglement-aware factoring.
  \item Anchor P with a structural certificate: closed-form geometric guarantee, dropped independence assumption, or per-layer reconstruction objective --- not heuristic ratios.
  \item Train / calibrate P to match S on the task distribution; preserve the abstraction-level shift (global $\to$ per-layer, full $\to$ adapter).
  \item Validate the compute-quality frontier across the deployment regime --- practical throughput must match paper-only efficiency; theory-practice disconnects are flagged.
\end{enumerate}
\textbf{Success conditions} (Oral):
\begin{itemize}[leftmargin=12pt,topsep=2pt,itemsep=2pt]
  \item \textbf{Anchor compression in a previously-unexploited structural property of the system, then design the compact proxy as a direct algebraic consequence of that property --- not as a heuristic shortcut.} \textcolor{gray}{\footnotesize evidence: \texttt{ICML\_2024\_1032}, \texttt{ICLR\_2025\_1350}, \texttt{ICML\_2024\_1016}, \texttt{NeurIPS\_2024\_1203}, \texttt{ICML\_2024\_0996}} \\ Oral papers consistently locate a specific structural fact the field had overlooked (weight--momentum entanglement, null-space of preserved keys, training-induced parameter manifold, outlier distribution geometry, mutual-information bottleneck) and derive the compression scheme from it. The structural framing is what makes the resulting method principled rather than another point on the compression Pareto front.
  \item \textbf{Decompose the system into orthogonal concerns whose interface is explicit and provably separable, then compress only one side aggressively while keeping the coupled side at full fidelity.} \textcolor{gray}{\footnotesize evidence: \texttt{NeurIPS\_2023\_0714}, \texttt{ICML\_2023\_0538}, \texttt{ICLR\_2024\_0886}, \texttt{NeurIPS\_2023\_0685}} \\ QLoRA (frozen 4-bit base + full-precision LoRA buffer), SparseGPT (global pruning $\to$ per-layer reconstruction), LoftQ (initialize LoRA to absorb quantization residual), and Monarch Mixer (single primitive across two axes) all succeed by exposing an interface that lets compression error be quarantined. Without this clean split, compression-induced distortion contaminates downstream learning.
  \item \textbf{Provide a closed-form, geometric, or information-theoretic derivation that the compressed proxy preserves the relevant signal --- before the empirical sweep.} \textcolor{gray}{\footnotesize evidence: \texttt{ICLR\_2025\_1350}, \texttt{ICLR\_2025\_1559}, \texttt{ICLR\_2025\_1544}, \texttt{ICML\_2025\_1731}} \\ AlphaEdit's null-space projection, the Rotation Trick's exact Jacobian, Progressive Distillation's sample-complexity proof, and IMM's moment-matching guarantee all turn what could have been an empirical ablation into a structural certificate. Reviewers treat the derivation as evidence the gain is not coincidence.
  \item \textbf{Reformulate a discrete, non-differentiable, or multi-step bottleneck as a continuous (or single-step) operation in a transformed domain so end-to-end gradients flow through what was previously opaque.} \textcolor{gray}{\footnotesize evidence: \texttt{ICLR\_2022\_0105}, \texttt{ICLR\_2025\_1559}, \texttt{ICLR\_2021\_0054}, \texttt{ICML\_2025\_1793}} \\ Fourier-domain stride learning, Householder/Cayley rotations replacing nearest-neighbor lookup, randomized AD over the DAG, and score-of-mixture single-step training all show that a domain change converts a hard combinatorial choice into a differentiable parameter --- enabling joint optimization that uniform-precision baselines cannot reach.
  \item \textbf{Use the system itself (or a fixed teacher) as the supervision oracle for compression --- reconstruction loss, activation matching, or self-filtered synthetic data --- eliminating dependence on the original training corpus.} \textcolor{gray}{\footnotesize evidence: \texttt{NeurIPS\_2025\_1898}, \texttt{ICLR\_2025\_1527}, \texttt{NeurIPS\_2025\_1869}, \texttt{ICML\_2024\_1016}} \\ KVzip (model reconstructs its own context), data-free CLIP distillation (teacher filters its own synthetic data), grafting (activation matching across operator substitutions), and Riemannian data-free compression all sidestep the data-access bottleneck by treating the trained system's outputs/geometry as the importance signal. This converts compression from a data problem into a model-internal problem.
  \item \textbf{Validate across multiple model scales, architectural families, and downstream tasks --- not just a single benchmark --- and characterize the failure regime explicitly.} \textcolor{gray}{\footnotesize evidence: \texttt{ICML\_2023\_0538}, \texttt{NeurIPS\_2023\_0714}, \texttt{ICLR\_2023\_0181}, \texttt{NeurIPS\_2024\_1203}} \\ Oral compression papers routinely cover several LLM sizes (7B--65B), multiple model families, and ablations of the binding constraint (e.g., DuQuant's outlier analysis, SparseGPT's sparsity sweep, climate compression's rare-event failure). Reviewers treat absence of this scope as evidence the structural claim is parochial.
\end{itemize}
\textbf{Failure modes} (Reject):
\begin{itemize}[leftmargin=12pt,topsep=2pt,itemsep=2pt]
  \item \textbf{Stitching two existing compression primitives into a 'wrapper' or 'pipeline' without identifying a joint structural property that makes their combination more than additive.} \textcolor{gray}{\footnotesize evidence: \texttt{ICLR\_2023\_0283}, \texttt{NeurIPS\_2023\_0761}, \texttt{NeurIPS\_2023\_0690}, \texttt{ICLR\_2024\_0887}} \\ MixQuant (bit-width search over arbitrary quantizers), XYZ Data Efficiency (sampling + routing), Architectural Compression (prune + distill for SD), and binary+distill for long-tail all draw 'limited novelty' verdicts because each module was independently validated and the combination is presented as additive rather than as a new structural object.
  \item \textbf{Theory-practice disconnect: theoretical analysis assumes one structure (random vectors, smooth functions, balanced data) but the implementation silently swaps in a different one, leaving the headline guarantee unsupported.} \textcolor{gray}{\footnotesize evidence: \texttt{ICLR\_2023\_0221}, \texttt{ICLR\_2024\_0876}, \texttt{ICML\_2025\_1782}, \texttt{ICLR\_2025\_1421}} \\ Efficient HDC theorizes over random hypervectors but experiments use a trained BNN encoder; Laplacian Sparsification proves a sparsifier bound without ablating that the sparsified step is the dominant cost; Regression-tree gradients require smoothness atypical of the regimes trees are used in; CCWT's algebraic equivalence isn't backed by wall-clock evidence. Reviewers explicitly flag the gap.
  \item \textbf{Claimed efficiency gains do not translate to the resource the deployment pipeline actually cares about (wall-clock latency, real throughput, supported hardware), leaving the compression a paper-only win.} \textcolor{gray}{\footnotesize evidence: \texttt{ICLR\_2023\_0275}, \texttt{ICLR\_2024\_0868}, \texttt{ICLR\_2025\_1374}, \texttt{NeurIPS\_2024\_1197}} \\ LVQ-VAE's runtime is too slow to be practical despite SOTA RD curves; KVTQ's ternary scheme fits limited hardware configurations; ClusComp's 1-bit numbers don't match throughput at higher bit-widths; Diff-PCC's stochastic decoder reaches good perceptual quality but reviewers flag undefined fidelity-diversity trade. The proxy P beats S on a metric nobody deploys against.
  \item \textbf{Empirical scope confined to small or stale benchmarks (MNIST/CIFAR, ResNet-only, single model size) so the structural claim cannot be distinguished from overfitting to that regime.} \textcolor{gray}{\footnotesize evidence: \texttt{ICLR\_2023\_0247}, \texttt{NeurIPS\_2023\_0637}, \texttt{ICML\_2025\_1766}, \texttt{NeurIPS\_2024\_1265}} \\ HRBP evaluated only on CIFAR/ImageNet ResNets; FITS lacks multiple-seed runs; OD$^3$ benchmarks only on RetinaNet/Faster R-CNN with ResNet backbone; NeCGS shows significant compression-time overhead vs baselines. Reviewers downgrade because the binding constraint at deployment scale is invisible.
  \item \textbf{Treating compression and an adjacent stage (adaptation, rebalancing, decoder) as independent modules when they are coupled, so quantization/pruning error propagates into the downstream signal without compensation.} \textcolor{gray}{\footnotesize evidence: \texttt{NeurIPS\_2024\_1190}, \texttt{NeurIPS\_2023\_0623}, \texttt{ICLR\_2024\_0806}, \texttt{NeurIPS\_2024\_1264}} \\ decoupleQ adds a floating-point correction but reviewers note its similarity to layer-norm without a coupled-stage analysis; EmbedDistill aligns geometries without showing the layer choice is sufficient; orthogonal-gradient NAS treats parameter sharing as independent updates; Adaptive Bit Switching's two-stage rounding lacks coupled analysis. The Oral analog (LoftQ, IR-QLoRA, AlphaEdit) explicitly co-designs the stages.
  \item \textbf{Asserting a universal compression rule over a population that is structurally heterogeneous (outliers, long-tail importance, hetero-modular components), without partitioning the population first.} \textcolor{gray}{\footnotesize evidence: \texttt{NeurIPS\_2023\_0750}, \texttt{NeurIPS\_2024\_1266}, \texttt{ICLR\_2025\_1652}} \\ Subspace Projection imposes one global subspace on all transformer layers; Neural Expressiveness applies an overlap criterion uniformly; XoRA imposes a fixed expander mask. Compare to Oral counterparts (BiLLM/SqueezeLLM partition by sensitivity strata, MoDeGPT decomposes per sub-module, DuQuant routes outliers separately) --- uniform schemes reliably underperform when importance is long-tailed.
\end{itemize}
\textbf{Oral vs Reject gap.} Oral compression papers lead with a *structural* discovery --- an unrecognized coupling, geometry, or distribution (weight-momentum entanglement in ExCP, null space of preserved keys in AlphaEdit, long-tail sensitivity strata in BiLLM, outlier-rotation duality in DuQuant) --- and derive the compact proxy as the unique consequence of that discovery, often with closed-form initialization or a provable bound. Reject papers, by contrast, present the proxy first and motivate it post-hoc: the abstract describes the method as 'combining', 'wrapping', 'extending', or 'searching over' established primitives (MixQuant, XYZ, Architectural Compression, OD$^3$). A second consistent gap is interface design: Oral papers expose a clean separability between the compressed component and an adjacent fidelity-preserving component (QLoRA's full-precision adapter, SparseGPT's per-layer reconstruction, LoftQ's SVD compensation), whereas Reject papers either compress everything uniformly (Subspace Projection) or stack independent stages without analyzing their interaction (decoupleQ, Adaptive Bit Switching). Finally, Oral papers report wall-clock or downstream-metric evidence that the binding constraint actually moves at deployment scale and characterize the failure regime; Reject papers cite FLOPs, parameter counts, or single-benchmark accuracy and leave the latency/throughput trade-off unexplored (KVTQ, ClusComp, LVQ-VAE).

\textbf{Oral vs HC gap.} HC papers (Fourier Neural Operator, Quant-Noise, Progressive Distillation, LLM-Pruner, BiLLM, SqueezeLLM, ProlificDreamer, Language Modeling Is Compression, Long-Term Memory) clear the same structural-discovery bar as Orals --- each identifies a non-obvious property (Shannon-ML duality, dense-and-sparse heterogeneity, residual-correction recursion, parameter-drift cache invalidation). The Oral-vs-HC delta is sharper on two axes: (1) Orals more often produce a *theorem-grade* artifact --- closed-form null-space projection, provable curriculum complexity, exact rotation Jacobian --- whereas HC papers tend to land their structural claim through targeted empirical analysis (e.g., SqueezeLLM's sensitivity profile, BiLLM's importance distribution); and (2) Orals more frequently introduce a primitive that *generalizes beyond compression* (SDS for any 2D-to-3D bridge, Monarch matrix as a substrate, score-distillation as a pretraining objective), where HC papers more often refine compression for a specific target. HC sample is small (n=9) so this is a tendency, not a rule.

\textbf{Reviewer expectations:}
\begin{itemize}[leftmargin=12pt,topsep=2pt,itemsep=1pt]
  \item \textit{[oral\_reviews]} A theoretical or geometric derivation showing that the compact proxy preserves the property the original system relied on, not just empirical parity at a chosen budget.
  \item \textit{[both]} Experiments spanning multiple model scales, architectures, and downstream tasks; reviewers explicitly question the structural claim when validation is confined to a single family or small dataset.
  \item \textit{[reject\_reviews]} Wall-clock or hardware-realistic measurement of the claimed efficiency gain; FLOPs and parameter counts alone are repeatedly called insufficient.
  \item \textit{[reject\_reviews]} Comparison against the most recent competing compression methods (not just classical baselines) on the same metric the paper claims to advance.
  \item \textit{[both]} Explicit characterization of the failure regime --- when does the proxy degrade, on which inputs, at which compression ratio --- rather than headline numbers at a fixed operating point.
  \item \textit{[oral\_reviews]} Evidence that the structural claim is non-trivial: a clean ablation isolating the proposed insight from generic engineering improvements (better hyperparameters, more data, larger teacher).
\end{itemize}
\textbf{Cognitive barriers:}
\begin{itemize}[leftmargin=12pt,topsep=2pt,itemsep=1pt]
  \item Compression and adaptation are assumed mutually exclusive --- practitioners expect quantization to corrupt gradient flow, blocking the insight (QLoRA, IR-QLoRA, LoftQ) that a small full-precision buffer can absorb compression error and make the two concerns orthogonal.
  \item Standard practice equates protecting capability with protecting parameters; reaching gradient-space, null-space, or manifold-space framings (Gradient Projection Memory, AlphaEdit, Riemannian compression) requires shifting the object of protection up an abstraction level.
  \item Co-produced artifacts (weights and optimizer state, quantization and adaptation, encoder and decoder) are stored or trained as if statistically independent; the entanglement insight (ExCP, IR-QLoRA, EmbedDistill) requires connecting two literatures (optimization theory and compression theory) that operate at different abstractions.
  \item Discrete or non-differentiable operations (quantization, stride, sorting) are treated as inherently un-trainable; recognizing them as projections, rotations, or frequency-domain crops with exact Jacobians (Rotation Trick, Learning Strides, Randomized AD) requires a domain-change move that practitioners trained inside the original domain don't naturally make.
\end{itemize}
\textbf{Representative examples:}
\begin{itemize}[leftmargin=12pt,topsep=2pt,itemsep=1pt]
  \item \textbf{[Oral]} \texttt{ICLR\_2025\_1350} \textemdash\ AlphaEdit converts a 'do not interfere' statistical objective into an exact null-space projection --- a closed-form geometric guarantee that exemplifies how Oral compression/edit work derives the compact proxy from a structural certificate rather than an empirical heuristic.
  \\ \textcolor{gray}{\footnotesize \textit{Lesson:} Convert a statistical objective ('don't interfere') into a closed-form geometric guarantee (null-space projection) --- Oral compression derives the proxy from a structural certificate, not a heuristic.}
  \item \textbf{[Oral]} \texttt{ICML\_2024\_1032} \textemdash\ ExCP exposes a previously-unrecognized statistical entanglement between adaptive-optimizer momentum and weights --- the compression ratio is a direct consequence of dropping an unjustified independence assumption.
  \\ \textcolor{gray}{\footnotesize \textit{Lesson:} The compression ratio is a direct consequence of dropping an unjustified independence assumption between optimizer momentum and weights --- diagnosis of a hidden statistical structure beats engineering tuning.}
  \item \textbf{[Oral]} \texttt{NeurIPS\_2023\_0714} \textemdash\ QLoRA exemplifies the orthogonal-decomposition pattern: aggressive 4-bit base storage plus a full-precision adapter that acts as a buffer absorbing compression error --- the clean interface is what makes 65B finetuning on one GPU work.
  \\ \textcolor{gray}{\footnotesize \textit{Lesson:} Orthogonal decomposition with a clean interface --- aggressive base storage + full-precision adapter buffer absorbs compression error. The interface clarity is what makes 65B-on-one-GPU work.}
  \item \textbf{[Oral]} \texttt{ICML\_2023\_0538} \textemdash\ SparseGPT decomposes global one-shot pruning into independent per-layer reconstruction problems solved with local Hessians, showing the abstraction-level shift (global $\to$ layer-wise reconstruction) that distinguishes Oral compression from heuristic mass pruning.
  \\ \textcolor{gray}{\footnotesize \textit{Lesson:} Abstraction-level shift (global one-shot pruning $\to$ per-layer reconstruction) distinguishes Oral compression from heuristic mass pruning.}
  \item \textbf{[Reject]} \texttt{ICLR\_2023\_0283} \textemdash\ MixQuant proposes bit-width search as a generic wrapper over any quantizer --- a pure 'combine existing primitives' move with no joint structural property, drawing reject votes for limited novelty and weak baselines.
  \\ \textcolor{gray}{\footnotesize \textit{Lesson:} Combining existing primitives without a joint structural property reads as a wrapper --- bit-width search over a quantizer is not a methodological contribution.}
  \item \textbf{[Reject]} \texttt{ICLR\_2023\_0221} \textemdash\ Efficient HDC's theory assumes random hypervectors but the experiments use a trained BNN encoder, exemplifying the theory-practice disconnect that reviewers reliably flag in compression rejects.
  \\ \textcolor{gray}{\footnotesize \textit{Lesson:} Theory-practice disconnects (random hypervectors in theory, trained encoder in experiments) are reliably flagged in compression rejects.}
  \item \textbf{[Reject]} \texttt{ICLR\_2025\_1374} \textemdash\ ClusComp claims 1-bit clustering compression but throughput collapses at higher bit-widths and the practical impact remains unproven --- a textbook 'paper-only efficiency win' rejection.
  \\ \textcolor{gray}{\footnotesize \textit{Lesson:} Throughput collapse at higher bit-widths converts the contribution to 'paper-only efficiency win' --- practical throughput is the binding metric.}
\end{itemize}
\end{tcolorbox}

\begin{tcolorbox}[colback=bgskill, colframe=cblue!70, title={\textbf{Engineer Auxiliary Supervision Signal} \textcolor{gray}{\small\texttt{auxiliary\_signal\_engineering}} \hfill confidence: \textbf{high}}, breakable, before skip=4pt, after skip=4pt]
\textcolor{gray}{\small Oral $n=23$ (meta 18/23) | HC $n=8$ (meta 7/8) | Reject $n=23$ (meta 23/23)}

\textbf{Sample note.} \textit{Sample is adequate (n\_oral=23, n\_reject=23, HC=8 with 7/8 meta coverage); meta-review coverage on the Oral side is 18/23 (78\%), so a handful of Oral claims rest on abstract\_* fields rather than reviewer voice.}

\textbf{Step-by-Step:}
\begin{enumerate}[leftmargin=14pt,topsep=2pt,itemsep=1pt]
  \item Identify a capability C the standard end-to-end loss underspecifies (geometric consistency, calibration, retrieval, control flow, modality alignment) with a clean failure mode.
  \item Diagnose the confounder shortcut: what easier signal is the model riding on instead of inducing C? Often a structural mismatch the standard loss is silent about.
  \item Construct an auxiliary signal A that disambiguates C from confounders --- A must be cheap to compute (no new annotation) and structurally aligned with C, not just correlated.
  \item Add A to the training objective with the right blending --- fixed weight, scheduled, gated by perplexity, or self-corrective via simulation gradients.
  \item Demonstrate via controlled ablation that C emerges only when A is present --- without A the confounder dominates; with A the shortcut closes.
  \item Show A does not introduce a new shortcut --- probe C with a held-out, A-independent test.
\end{enumerate}
\textbf{Success conditions} (Oral):
\begin{itemize}[leftmargin=12pt,topsep=2pt,itemsep=2pt]
  \item \textbf{Name a single specific capability the primary loss cannot teach, and design the auxiliary signal so its information content maps directly onto that capability --- not a generic 'representation quality' improvement.} \textcolor{gray}{\footnotesize evidence: \texttt{NeurIPS\_2024\_1298}, \texttt{NeurIPS\_2023\_0744}, \texttt{ICML\_2025\_1702}} \\ RG-SAN isolates spatial-grounding by extracting position-only supervision from referring expressions; Toolformer isolates tool-use by filtering on perplexity reduction; DistiLLM-2 isolates teacher-vs-student differential via asymmetric contrastive --- in each case reviewers cite the precise capability-to-signal mapping as the contribution.
  \item \textbf{Source the auxiliary signal from something already present in the system or data (model's own loss, renderer gradients, search-tree rejected branches, hierarchical-clustering levels, taxonomy structure) rather than collecting new annotation --- and make the 'free' nature of the signal central to the claim.} \textcolor{gray}{\footnotesize evidence: \texttt{NeurIPS\_2023\_0744}, \texttt{ICML\_2024\_1037}, \texttt{ICLR\_2021\_0017}} \\ Toolformer mines API examples filtered by its own perplexity, hierarchical-clustering papers convert clustering output into multi-scale pseudo-labels, GAN-as-3D-supervisor uses a 2D-only model as the 3D oracle --- Oral meta-reviews emphasize that no new labels were needed, which is what makes the signal a methodological contribution rather than an annotation effort.
  \item \textbf{Either prove or empirically demonstrate the auxiliary signal carries information disjoint from the primary loss --- via a derivation (concentration bound, Bayes' theorem, orthogonality construction) or via a measurable probe of the targeted capability.} \textcolor{gray}{\footnotesize evidence: \texttt{ICML\_2023\_0569}, \texttt{NeurIPS\_2024\_1166}, \texttt{ICLR\_2022\_0132}} \\ Weighted-flow-diffusion uses high-dimensional concentration to prove the joint signal preserves cluster structure, Bayesian label mapping derives the aggregation from Bayes' theorem, RISP constructs the orthogonality loss from renderer Jacobians --- Oral meta-reviews specifically cite the principled derivation as what lifts the work above 'plausible heuristic'.
  \item \textbf{If the signal involves self-bootstrapping (model critiques / filters / improves itself), build in a stability mechanism --- an explicit filter criterion, a convergence argument, or a closed-form check --- so reviewers can rule out noise reinforcement.} \textcolor{gray}{\footnotesize evidence: \texttt{NeurIPS\_2023\_0744}, \texttt{ICLR\_2025\_1435}, \texttt{NeurIPS\_2025\_1925}} \\ Toolformer's perplexity-reduction threshold, EASYTOOL's grounded interaction-with-tool diagnosis, and the multimodal-boosting paper's convergence proof all explicitly address the 'circular self-evaluation' worry that recurs in reject meta-reviews.
  \item \textbf{Encode the auxiliary supervision into the model's native representation (vocabulary tokens, embedding head, output channels) when possible, so the auxiliary mechanism is end-to-end optimized rather than glued on as a separate module.} \textcolor{gray}{\footnotesize evidence: \texttt{ICLR\_2024\_0945}, \texttt{NeurIPS\_2023\_0745}, \texttt{NeurIPS\_2024\_1166}} \\ Self-RAG's reflection tokens, ToolkenGPT's toolken embeddings, and Bayesian label-mapping's probabilistic output aggregation all internalize the auxiliary decision into the model's generation framework --- meta-reviews praise the resulting modularity, efficiency, and absence of catastrophic forgetting.
  \item \textbf{Show the auxiliary signal works disproportionately to its size --- small fraction of supervision, light auxiliary head, or zero new annotations producing near-fully-supervised performance --- making it economically and conceptually distinctive.} \textcolor{gray}{\footnotesize evidence: \texttt{ICLR\_2025\_1431}, \texttt{ICLR\_2021\_0035}, \texttt{NeurIPS\_2024\_1171}} \\ FSBM gets near-fully-supervised performance from <8\% anchor pairs, inverse-graphics achieves results matching orders-of-magnitude-more-labeled baselines, and CAT3D replaces expensive multi-view capture entirely --- quantitatively surprising leverage is a recurring Oral signature.
\end{itemize}
\textbf{Failure modes} (Reject):
\begin{itemize}[leftmargin=12pt,topsep=2pt,itemsep=2pt]
  \item \textbf{Auxiliary signal is borrowed from a neighbouring paradigm (contrastive, mixup, denoising) and bolted onto a new setting without identifying a specific capability it uniquely repairs --- reviewers immediately read this as a domain-port rather than a methodological contribution.} \textcolor{gray}{\footnotesize evidence: \texttt{ICLR\_2023\_0188}, \texttt{ICLR\_2023\_0354}, \texttt{ICLR\_2024\_0901}} \\ ContraGen ports contrastive to causal LMs, SupCR ports SupCon to regression, SupReMix ports mixup to contrastive regression --- all three meta-reviews flag 'limited technical novelty' because the auxiliary objective is generic and the diagnosis of why this signal is needed is missing.
  \item \textbf{The auxiliary signal is plausibly motivated but never shown to be the load-bearing component --- ablations against strong baselines that solve the same need differently are absent, so the causal claim that the signal fixes the diagnosed problem is unverified.} \textcolor{gray}{\footnotesize evidence: \texttt{ICLR\_2024\_0824}, \texttt{ICLR\_2024\_0793}, \texttt{ICLR\_2023\_0328}} \\ These papers add an auxiliary alignment / shadow-set / cross-layer signal but reviewers explicitly call out that the structural-mismatch hypothesis is not isolated by ablation; the gain could come from the extra capacity or extra data instead.
  \item \textbf{Self-bootstrapping loop with no stability or quality check --- pseudo-labels / self-generated feedback are trusted blindly, so reviewers question whether the loop reinforces noise instead of correcting it.} \textcolor{gray}{\footnotesize evidence: \texttt{ICLR\_2024\_0850}, \texttt{ICLR\_2023\_0317}, \texttt{ICLR\_2024\_0937}} \\ Hexa's distant-supervision loop, ProsodyBERT's offline cluster pseudo-labels, and RepoFusion's structural-context training all leave open whether the auxiliary signal is reliable enough to substitute for direct supervision; meta-reviews ask for evidence the bootstrapping does not collapse.
  \item \textbf{The auxiliary signal lacks theoretical or mechanistic justification connecting it to the targeted capability, so reviewers cannot tell when or why it should work.} \textcolor{gray}{\footnotesize evidence: \texttt{ICLR\_2023\_0330}, \texttt{ICLR\_2025\_1394}, \texttt{NeurIPS\_2024\_1322}} \\ Safe RL's contrastive risk classifier has internal inconsistencies the AC explicitly flags ('it is unclear why and when the proposed approach works'); Thomson-energy and MMJ-distance papers borrow from physics/topology without grounding the borrowed criterion in the clustering objective.
  \item \textbf{Auxiliary signal targets a capability whose absence was never empirically demonstrated as the bottleneck --- the paper claims a need the benchmarks do not validate.} \textcolor{gray}{\footnotesize evidence: \texttt{ICLR\_2024\_0803}, \texttt{ICLR\_2025\_1455}, \texttt{NeurIPS\_2024\_1335}} \\ DS-CLIP's de-duplication, IDS's information-theoretic re-weighting, and hierarchical-prototypes-over-tree all assume their identified gap (redundancy, augmentation noise, hierarchical structure) is what limits performance, but reviewers note competitive baselines achieve similar gains without the auxiliary signal.
\end{itemize}
\textbf{Oral vs Reject gap.} Oral papers name a specific capability C that the primary loss provably underspecifies, then construct an auxiliary signal whose information content is causally tied to C and demonstrate (often via an ablation that turns the signal off) that C collapses without it --- Toolformer's perplexity filter, RG-SAN's position-only supervision, and DistiLLM-2's asymmetric contrastive loss all isolate exactly what the auxiliary signal teaches. Reject papers introduce an auxiliary signal but conflate two moves: they argue 'this signal could help' rather than 'this signal is the only thing that can teach C'. Concretely: Reject papers tend to (a) reuse a known auxiliary objective (contrastive, mixup, pseudo-label) ported to a new setting without diagnosing why the existing literature's version is insufficient, (b) skip ablations that would isolate the auxiliary signal's contribution from added capacity or extra data, and (c) leave the bootstrapping loop unaudited so reviewers cannot tell whether the self-generated supervision is stable. Oral papers also routinely show a counterintuitive sign --- adding less information, using model outputs as negatives, supervising with a 'lower-dimensional' source --- which forces them to provide a mechanistic explanation; Reject papers usually add information additively, which removes the need for explanation but also removes the conceptual hook.

\textbf{Oral vs HC gap.} HC papers execute auxiliary-signal patterns that already exist in the literature with high engineering quality on practical benchmarks (ANCE's dynamic hard-negative mining, DINO's contrastive denoising, Gorilla's retrieval-augmented training, Q-Align's discrete-then-aggregate scoring, Self-Debug's recursive self-critique). Reviewers consistently praise empirical results but flag the auxiliary signal itself as 'well-known in the literature' or 'incremental' (DINO meta-review explicitly, ANCE 'idea might not be so new'). Oral papers, by contrast, introduce an auxiliary signal that is structurally novel --- either an unexploited source (Toolformer's perplexity-as-filter, Self-RAG's reflection tokens, RG-SAN's rule-extracted spatial-only labels, FSBM's <8\% anchor constraints), or a non-obvious connection between two formalisms (LSEnet's hyperbolic + structural-entropy pairing, Lorentz multimodal-imbalance-as-classifier-gap reframing). The HC ceiling is set by reviewers reading the auxiliary signal as a 'combination of existing techniques' even when results are strong; Oral papers cross that ceiling by making the signal itself the conceptual contribution.

\textbf{Reviewer expectations:}
\begin{itemize}[leftmargin=12pt,topsep=2pt,itemsep=1pt]
  \item \textit{[both]} Demonstrate that the targeted capability cannot be recovered by the primary loss alone --- typically by showing the model fails on a probe of that capability without the auxiliary signal, then succeeds with it.
  \item \textit{[reject\_reviews]} Provide ablations that isolate the auxiliary signal from confounders (extra data, extra capacity, extra training stages); reject decisions cite missing or weak ablations as the dominant reason.
  \item \textit{[oral\_reviews]} Justify the auxiliary signal mechanistically --- why this signal carries information about the target capability, ideally with a theoretical guarantee, a concentration result, or a derivation (FSBM, LSEnet, weighted-flow-diffusion all earn praise for this).
  \item \textit{[reject\_reviews]} Show stability of any self-bootstrapping loop: filtering criteria, pseudo-label quality checks, or convergence proofs; reviewers reject when the loop could plausibly reinforce its own noise.
  \item \textit{[both]} Position cleanly against prior auxiliary-signal work in adjacent domains --- when contrastive / pseudo-label / self-critique techniques exist nearby, reviewers expect explicit articulation of what the new signal does that those prior signals cannot.
  \item \textit{[oral\_reviews]} When the auxiliary signal does something counterintuitive (uses less information, uses model outputs as negatives, uses a 'lower-dimensional' source for a higher-dimensional task), explain the mechanism explicitly --- Oral reviewers reward this; Reject reviewers punish hand-waving.
\end{itemize}
\textbf{Cognitive barriers:}
\begin{itemize}[leftmargin=12pt,topsep=2pt,itemsep=1pt]
  \item Researchers treat sources of supervision as fixed and externally-given (annotations, rewards, ground-truth labels), which makes it psychologically hard to reconceptualize a model's own loss curve, a renderer's gradients, or a search procedure's discarded candidates as legitimate training signal --- these were 'tools', 'analyses', or 'waste', not supervision (Toolformer NeurIPS\_2023\_0744, CPO NeurIPS\_2024\_1172, RISP ICLR\_2022\_0132).
  \item There is a default assumption that adding more information helps and removing it hurts; auxiliary-signal designs that deliberately strip information (RG-SAN's position-only supervision, Bayesian label mapping replacing hard assignments) feel like self-sabotage until reviewers see that the conflated signal was the actual blocker (RG-SAN NeurIPS\_2024\_1298, NeurIPS\_2024\_1166).
  \item Paradigm-boundary policing: 'unsupervised' must avoid all external knowledge, 'contrastive' is for representation, tools are 'external calls' --- these labels prevent practitioners from noticing that the underlying machinery is reusable across boundaries (ICML\_2024\_1046 image clustering with VLM guidance, ICLR\_2024\_0945 Self-RAG, NeurIPS\_2023\_0745 ToolkenGPT).
  \item Self-evaluation feels circular: using a model to critique itself or to filter its own training data triggers the intuition that errors will compound, masking that a sufficiently capable model already encodes the correct behaviour and only needs the right interface to surface it (NeurIPS\_2023\_0744 Toolformer, ICLR\_2024\_0945 Self-RAG, ICLR\_2025\_1435 documentation self-repair).
\end{itemize}
\textbf{Representative examples:}
\begin{itemize}[leftmargin=12pt,topsep=2pt,itemsep=1pt]
  \item \textbf{[Oral]} \texttt{NeurIPS\_2023\_0744} \textemdash\ Toolformer is the cleanest instance of using a model's own existing objective (perplexity) as an automatic filter for self-generated training data, closing the auxiliary-signal loop without any human labels --- it is the canonical 'capability emerges without new annotation' execution.
  \\ \textcolor{gray}{\footnotesize \textit{Lesson:} Use the model's own existing objective (perplexity) as the automatic filter for self-generated training data --- the canonical 'capability emerges without new annotation' execution.}
  \item \textbf{[Oral]} \texttt{NeurIPS\_2024\_1298} \textemdash\ RG-SAN is distinctive because it deliberately reduces the supervision (position-only, no semantic content) on the spatial-awareness component, isolating the failure mode and forbidding the semantic-matching shortcut --- a textbook capability-isolation move.
  \\ \textcolor{gray}{\footnotesize \textit{Lesson:} Deliberately reduce the supervision (position-only, no semantic content) on the spatial-awareness component --- isolating the failure mode by removing the shortcut-inviting signal is a textbook move.}
  \item \textbf{[Oral]} \texttt{ICLR\_2024\_0945} \textemdash\ Self-RAG converts pipeline-level decisions (when to retrieve, how to critique) into in-vocabulary reflection tokens, internalizing what would otherwise be external auxiliary supervision and making the entire control flow end-to-end trainable.
  \\ \textcolor{gray}{\footnotesize \textit{Lesson:} Internalize pipeline-level decisions as in-vocabulary tokens --- what would otherwise be external auxiliary supervision becomes end-to-end trainable.}
  \item \textbf{[Oral]} \texttt{ICLR\_2022\_0132} \textemdash\ RISP uses the differentiable renderer's gradients with respect to nuisance variables to define an orthogonality loss --- turning the source of unwanted variation into the supervision that removes it, a structurally non-obvious self-corrective signal.
  \\ \textcolor{gray}{\footnotesize \textit{Lesson:} Turn the source of unwanted variation into the supervision that removes it --- using the differentiable renderer's gradients w.r.t. nuisance variables as an orthogonality loss is structurally non-obvious.}
  \item \textbf{[Reject]} \texttt{ICLR\_2023\_0354} \textemdash\ SupCR adds a pairwise contrastive objective ordered by target distance for regression, but the meta-review reads it as a straightforward port of supervised-contrastive-for-classification --- the auxiliary signal lacks a uniquely diagnosed capability gap.
  \\ \textcolor{gray}{\footnotesize \textit{Lesson:} A straightforward port from classification to regression is read as a wrapper --- auxiliary signals must come with a uniquely diagnosed capability gap they close.}
  \item \textbf{[Reject]} \texttt{ICLR\_2024\_0824} \textemdash\ Grafted an auxiliary contrastive alignment objective onto T2I fine-tuning but, as the meta-review notes, the structural-mismatch hypothesis was never validated by controlled ablations --- the auxiliary signal's necessity was assumed, not demonstrated.
  \\ \textcolor{gray}{\footnotesize \textit{Lesson:} The structural-mismatch hypothesis must be validated by controlled ablations --- auxiliary-signal necessity cannot be assumed.}
  \item \textbf{[Reject]} \texttt{ICLR\_2023\_0330} \textemdash\ Safe RL with a contrastive risk classifier attempted to repurpose contrastive learning as a calibrated risk estimator, but the AC flagged theoretical inconsistencies that left it unclear when the auxiliary signal is reliable --- the mechanism was not grounded.
  \\ \textcolor{gray}{\footnotesize \textit{Lesson:} Repurposing contrastive learning as a calibrated risk estimator demands a theoretical grounding for when the signal is reliable --- without that, the mechanism is unjustified.}
\end{itemize}
\end{tcolorbox}

\begin{tcolorbox}[colback=bgskill, colframe=cblue!70, title={\textbf{Align Heterogeneous Sources in Shared Space} \textcolor{gray}{\small\texttt{shared\_latent\_alignment}} \hfill confidence: \textbf{high}}, breakable, before skip=4pt, after skip=4pt]
\textcolor{gray}{\small Oral $n=45$ (meta 34/45) | HC $n=18$ (meta 12/18) | Reject $n=52$ (meta 51/52)}

\textbf{Sample note.} \textit{Sample is large and balanced (n\_oral=45, n\_reject=52, n\_hc=18) with strong meta coverage on the reject side (51/52) and adequate coverage on the oral side (34/45, 76\%); HC sample of 18 supports the oral\_hc\_gap claim.}

\textbf{Step-by-Step:}
\begin{enumerate}[leftmargin=14pt,topsep=2pt,itemsep=1pt]
  \item Identify $\ge$ 2 heterogeneous sources / modalities / tasks the literature treats in isolation --- name what each contributes.
  \item Establish a structural identity (algebraic, geometric, semantic) that justifies merging them into one shared space --- not just empirical correlation between outputs.
  \item Choose the shared space Z carefully: continuous latent, tokenized vocabulary, or operator-level alignment --- pick the level that respects the structural identity.
  \item Learn encoders \{E\_i\} mapping each source into Z under a diagnosis-first alignment objective --- identify which component is the modality-bias bottleneck and individualize only that.
  \item Operate in Z with unified downstream operators --- alignment is a means; the contribution is what becomes derivable / transferable / interpretable in Z.
  \item Isolate which alignment component carries the gain via per-component ablation --- bundling without isolation draws 'where does the improvement come from' critique.
\end{enumerate}
\textbf{Success conditions} (Oral):
\begin{itemize}[leftmargin=12pt,topsep=2pt,itemsep=2pt]
  \item \textbf{Identify a specific structural property of the source/target modalities that justifies why a shared space is even reachable, and instantiate that property as a concrete architectural or objective choice (not a generic 'projection layer + alignment loss').} \textcolor{gray}{\footnotesize evidence: \texttt{ICLR\_2021\_0041}, \texttt{NeurIPS\_2025\_1906}, \texttt{ICML\_2025\_1825}, \texttt{ICLR\_2025\_1635}} \\ Bisimulation distance (0041) gives a fixed-point equation that defines what 'aligned' means; MokA isolates the A matrix as the modality-bias bottleneck; VideoRoPE reasons about frequency allocation rather than token-index allocation; the modality-gap paper traces two phenomena to one information-imbalance cause. In each case the shared-space construction follows from a diagnosed property, not from a generic recipe.
  \item \textbf{Introduce a domain-appropriate discretization or tokenization that lets a heterogeneous source enter the shared space using the dominant paradigm's machinery, rather than reusing an off-the-shelf tokenizer.} \textcolor{gray}{\footnotesize evidence: \texttt{ICLR\_2022\_0081}, \texttt{ICML\_2023\_0403}, \texttt{ICML\_2023\_0418}, \texttt{ICML\_2024\_1140}, \texttt{NeurIPS\_2025\_1894}} \\ BEiT introduces visual tokens to make masked prediction work for images; Mu$^2$SLAM quantizes speech into tokens that share the text vocabulary; BEATs makes the tokenizer itself learnable and co-evolving; Moirai designs patch+output-scaling to bridge time-series heterogeneity; InfinityStar designs a spacetime pyramid tokenizer. The tokenizer choice is treated as the load-bearing contribution, not as preprocessing.
  \item \textbf{Demonstrate that the shared space unlocks a qualitatively new capability (zero-shot transfer, cross-domain morphing, interpretability, any-to-any generation) rather than only incremental gains on existing benchmarks.} \textcolor{gray}{\footnotesize evidence: \texttt{ICLR\_2022\_0137}, \texttt{ICLR\_2024\_0863}, \texttt{ICML\_2024\_1074}, \texttt{ICML\_2024\_1143}, \texttt{NeurIPS\_2023\_0743}} \\ StyleAlign uses correspondence in the latent space to perform cross-domain morphing without retraining; CLIP-decomposition turns alignment into an interpretability tool; NExT-GPT enables any-to-any generation through frozen modules; VideoPoet shows emergent multi-task generation from unified tokens; LLaVA bootstraps multimodal instruction-following via text-only GPT-4. Reviewers respond strongly to capabilities that did not exist before alignment.
  \item \textbf{Bridge frozen pretrained components with a small learned interface trained in stages, rather than retraining either side end-to-end, when reusing strong existing models from each modality.} \textcolor{gray}{\footnotesize evidence: \texttt{ICML\_2023\_0433}, \texttt{ICML\_2024\_1074}, \texttt{NeurIPS\_2025\_1944}} \\ ORCA aligns distributions before fine-tuning a frozen backbone for cross-modal transfer; NExT-GPT trains only lightweight projections between frozen encoders/decoders; Vita-1.5 stages vision then speech curriculum to avoid mutual interference. The decomposition of 'where the alignment lives' into a separable bridging step is consistently rewarded.
  \item \textbf{Apply per-modality objectives or per-modality components within a shared backbone to prevent one modality from dominating, with explicit ablation showing the asymmetric treatment matters.} \textcolor{gray}{\footnotesize evidence: \texttt{ICLR\_2025\_1632}, \texttt{NeurIPS\_2025\_1906}, \texttt{ICLR\_2025\_1397}} \\ Transfusion mixes next-token loss for text and diffusion loss for images in one transformer; MokA splits A matrices per modality but shares B and adds cross-modal attention; DEPT decouples embeddings per source while sharing the transformer body. Each shows that uniform sharing creates interference and that selective splitting---at exactly the right component---resolves it.
\end{itemize}
\textbf{Failure modes} (Reject):
\begin{itemize}[leftmargin=12pt,topsep=2pt,itemsep=2pt]
  \item \textbf{Stack alignment as a generic recipe (separate encoders + projection + contrastive/reconstruction loss + downstream head) without identifying any structural property of the modalities that makes the alignment load-bearing.} \textcolor{gray}{\footnotesize evidence: \texttt{ICLR\_2024\_0782}, \texttt{ICLR\_2024\_0948}, \texttt{ICLR\_2025\_1515}, \texttt{ICLR\_2025\_1608}} \\ Reject meta-reviews repeatedly note 'combination of existing techniques' or 'limited methodological innovation' (MULAN, SMILE, ICLR\_2025\_1608). The latent-alignment scaffolding looks correct in isolation, but reviewers flag that the gains are not attributed to any specific property and could equally come from any of the components.
  \item \textbf{Inherit a tokenization or backbone from a source domain (ViT, LLM tokenizer, CLIP encoder) and apply it to the target domain without justifying that the target's structural primitives match the source's, or showing the choice is the right one.} \textcolor{gray}{\footnotesize evidence: \texttt{ICLR\_2023\_0391}, \texttt{ICLR\_2025\_1414}, \texttt{ICLR\_2024\_0966}} \\ ViTMTSC drops time series into ViT patches; MEIT encodes ECG via a generic 1D CNN for an MLLM; TEST aligns to LLM embedding anchors. Reviewers consistently ask why this specific bridge is appropriate and note absence of analysis showing what about the target's structure justifies the inherited tokenization.
  \item \textbf{Marginal/inconsistent improvements across datasets, where the aligned model wins on one benchmark but only ties or loses on others, exposing that the shared representation is not the binding constraint.} \textcolor{gray}{\footnotesize evidence: \texttt{ICLR\_2023\_0259}, \texttt{ICLR\_2025\_1588}, \texttt{NeurIPS\_2024\_1175}} \\ Cross-domain MUDA gains hold on DomainNet but not Office-31; DynaMer Adapter shows 0.5-1\% improvements; concept-text aggregation receives borderline accept-reject scores. ACs treat inconsistency across datasets as evidence that the alignment mechanism is not the actual driver of the wins.
  \item \textbf{Train the alignment with synthetically generated paired data without validating that the synthesis preserves the structural correspondence the alignment depends on.} \textcolor{gray}{\footnotesize evidence: \texttt{ICLR\_2025\_1387}, \texttt{ICLR\_2025\_1399}, \texttt{NeurIPS\_2023\_0610}, \texttt{NeurIPS\_2023\_0611}} \\ CtrlSynth, DiffTell, CoVR, and the captioning-curation paper all align via paired data fabricated by another generative model. Reviewers question whether gains transfer to real distributions and whether the synthetic pairings encode the relationship being learned, treating this as an unverified premise of the alignment.
  \item \textbf{Claim a unified shared space across many tasks but evaluate on a narrow set, leaving the universality claim unsupported.} \textcolor{gray}{\footnotesize evidence: \texttt{ICLR\_2024\_0782}, \texttt{ICLR\_2024\_0927}, \texttt{ICLR\_2024\_0986}} \\ Conditional Protein Generation claims joint-task versatility but is evaluated on inverse folding + docking only; PROSE claims bimodal-bimodal generality on 15 ODE systems; USLNet claims unsupervised sign-language translation/generation but reviewers note poor numerical results. ACs ask for either an extra task showing the unified space pays off, or a narrower claim.
  \item \textbf{Conflate distinct subproblems behind a unification banner without ablations isolating which alignment component carries the gain.} \textcolor{gray}{\footnotesize evidence: \texttt{NeurIPS\_2023\_0688}, \texttt{ICLR\_2024\_0801}, \texttt{ICML\_2025\_1737}} \\ Multitask face-forgery joint embedding bundles symbolic supervision + shared space + multitask loss without isolating contributions; LEAD bundles dataset curation + sample-level + subject-level contrastive without showing each is needed; multilingual-vision claims rest on aggregated metrics. Reviewers explicitly demand ablations attributing improvements.
\end{itemize}
\textbf{Oral vs Reject gap.} Oral papers diagnose a specific structural property (a fixed-point condition, an asymmetric matrix role, a frequency-allocation problem, an information-imbalance source) and let the shared-space construction fall out of that diagnosis---the alignment objective is not the contribution, the diagnosis is. Reject papers, by contrast, present the alignment scaffolding itself as the contribution: pick two encoders, project to a shared space, add contrastive or reconstruction loss, evaluate. Oral papers also typically introduce a domain-appropriate discretization or interface (BEiT visual tokens, Mu$^2$SLAM speech tokens, BEATs co-evolving tokenizer, Moirai patches) that becomes the load-bearing element; reject papers more often inherit a tokenizer/backbone from a source domain (ViT for time series, LLM embedding for ECG/protein) without showing that the target's primitives match. Finally, Oral papers demonstrate the alignment unlocks a new capability---cross-domain morphing, zero-shot generation across modalities, interpretability tooling---while reject papers report 0.5-2\% gains on existing benchmarks with inconsistent improvements across datasets, which ACs read as evidence the shared space is not the binding constraint.

\textbf{Oral vs HC gap.} HC papers (n=18) tend to be polished engineering instantiations of an established alignment paradigm: BLIP-2's Q-Former bridges frozen models, AIM adds spatiotemporal adapters to a frozen image model, Make-A-Video factorizes T2V into modular reuse, TimesNet reshapes 1D series into 2D for image-CNN reuse. They succeed by choosing the right bridge between known components and validating it across benchmarks. Oral papers go further by introducing a structural reframing that did not exist before: StyleAlign formalizes correspondence as a discovered invariant rather than a designed objective; CLIP-decomposition turns alignment into an interpretability protocol via additive residual decomposition; the modality-gap paper unifies two phenomena under one information-theoretic cause; MokA diagnoses A vs B asymmetry as the modality-imbalance bottleneck. The Oral move is to make the alignment's existence or geometry into the discovery; the HC move is to bridge well between pre-existing alignments. HC papers also more often defer the conceptual contribution to dataset/scale and validate empirically, while Oral papers spend more pages on why the construction had to take this shape.

\textbf{Reviewer expectations:}
\begin{itemize}[leftmargin=12pt,topsep=2pt,itemsep=1pt]
  \item \textit{[both]} Ablations isolating which alignment component (which projector, which loss, which matrix, which tokenizer choice) actually drives the gain, not just whole-system improvements.
  \item \textit{[oral\_reviews]} Demonstration that the shared space transfers or generalizes across tasks/datasets/modalities the model was not directly trained on---zero-shot or held-out behavior, not just in-distribution wins.
  \item \textit{[reject\_reviews]} Justification for why the chosen tokenization or bridging interface matches the target domain's structural primitives, especially when the bridge is borrowed from another domain.
  \item \textit{[reject\_reviews]} Comparisons against the strongest contemporary baselines in each modality, including modern foundation models and recent unification methods, to defend that the alignment is the source of improvement.
  \item \textit{[reject\_reviews]} Consistent improvements across multiple benchmarks/datasets; ACs treat dataset-by-dataset inconsistency as evidence that the alignment mechanism is not the actual driver.
  \item \textit{[reject\_reviews]} Clear articulation of what is conceptually novel beyond a 'separate encoders + shared projection + contrastive loss' template; reviewers explicitly flag and reject this template when the structural insight is missing.
\end{itemize}
\textbf{Cognitive barriers:}
\begin{itemize}[leftmargin=12pt,topsep=2pt,itemsep=1pt]
  \item Domain-separation prior: practitioners assume two modalities (speech vs. text, EEG vs. signal, time series vs. language) require different architectures, vocabularies, and objectives, blocking the move to recognize discretization as the bridge that resolves the modality gap (Mu$^2$SLAM, BEiT, time-series-as-pixels).
  \item Tokenizer-as-fixed-preprocessing: the prediction target / discretization is treated as an upstream design choice rather than a learnable parameter that should co-evolve with the representation (BEATs, BEiT).
  \item Direct-alignment default: the assumption that two heterogeneous sources must be aligned head-on, missing that an intermediate third modality (text, HTML, latent semantic space) can serve as a bridge that makes the problem tractable (Pix2Struct, UniMedI vs. its rejection, 3D-LLM).
  \item Architecture-over-distribution framing: cross-modal transfer is assumed to require modality-specific architectures or adapters, obscuring that pretrained backbones are modality-agnostic if input distributions are matched (ORCA, BLIP-2).
\end{itemize}
\textbf{Representative examples:}
\begin{itemize}[leftmargin=12pt,topsep=2pt,itemsep=1pt]
  \item \textbf{[Oral]} \texttt{NeurIPS\_2025\_1906} \textemdash\ MokA exemplifies the diagnosis-first pattern: it identifies the A matrix specifically (not B, not both) as the modality-bias bottleneck in LoRA-based MLLM tuning, and then individualizes only that component while keeping B shared---the alignment fix follows directly from the diagnosis.
  \\ \textcolor{gray}{\footnotesize \textit{Lesson:} The diagnosis-first pattern --- identify the A matrix specifically (not B, not both) as the modality-bias bottleneck, then individualize only that component. The fix follows directly from the diagnosis.}
  \item \textbf{[Oral]} \texttt{ICML\_2023\_0403} \textemdash\ Mu$^2$SLAM shows the canonical 'tokenize the continuous modality so it joins the discrete vocabulary' move, treating speech tokens as another sequential symbol system inside a multilingual text framework rather than building a separate speech encoder.
  \\ \textcolor{gray}{\footnotesize \textit{Lesson:} Tokenize the continuous modality so it joins the discrete vocabulary --- treating speech tokens as another sequential symbol system is a cleaner unification than separate-encoder + alignment-loss.}
  \item \textbf{[Oral]} \texttt{ICLR\_2024\_0863} \textemdash\ CLIP decomposition repurposes an existing shared latent space as an interpretability tool by additively decomposing the image representation across heads/layers/tokens and projecting summands into text---turning alignment into a downstream capability rather than a training objective.
  \\ \textcolor{gray}{\footnotesize \textit{Lesson:} Repurpose an existing shared latent as an interpretability tool --- additive decomposition across heads / layers / tokens, projecting summands into text. Alignment becomes a downstream capability, not the training objective.}
  \item \textbf{[Oral]} \texttt{ICLR\_2025\_1635} \textemdash\ This paper unifies two separate empirical phenomena (modality gap and object bias) under a single information-imbalance cause---showing the cognitive move that distinguishes Orals: collapsing apparently distinct symptoms into one root property of the shared space.
  \\ \textcolor{gray}{\footnotesize \textit{Lesson:} Collapse apparently-distinct symptoms into one root property of the shared space --- information-imbalance unifies modality gap and object bias under a single cause.}
  \item \textbf{[Reject]} \texttt{ICLR\_2025\_1515} \textemdash\ MULAN attaches a structure-aware adapter to a protein language model to inject a complementary modality, but reviewers found the alignment template indistinguishable from prior structure-aware PLMs and the gains incremental---a clean case of 'alignment scaffolding as the contribution' failing.
  \\ \textcolor{gray}{\footnotesize \textit{Lesson:} Alignment scaffolding indistinguishable from prior structure-aware work draws 'incremental' critique --- the structural identity must be specifically new.}
  \item \textbf{[Reject]} \texttt{ICLR\_2023\_0391} \textemdash\ ViTMTSC reshapes multivariate time series into 2D images and applies a pretrained ViT---exemplifying the inherited-tokenizer failure mode where the source-domain primitive is reused without justifying that target structure matches.
  \\ \textcolor{gray}{\footnotesize \textit{Lesson:} Inheriting a tokenizer from another domain (here: ViT for time series) without justifying that target structure matches the source primitive is the inherited-tokenizer failure.}
  \item \textbf{[Reject]} \texttt{NeurIPS\_2023\_0688} \textemdash\ Joint embedding for face forgery bundles multitask + symbolic supervision + shared latent space without isolating which alignment component carries the gain, drawing exactly the 'where does the improvement come from' critique reviewers raise repeatedly.
  \\ \textcolor{gray}{\footnotesize \textit{Lesson:} Bundling multitask + symbolic supervision + shared latent without per-component ablation draws 'where does the improvement come from' --- alignment claims must be ablated.}
\end{itemize}
\end{tcolorbox}

\begin{tcolorbox}[colback=bgskill, colframe=cblue!70, title={\textbf{Aggregate Over Diverse Candidates} \textcolor{gray}{\small\texttt{candidate\_aggregation\_over\_diversity}} \hfill confidence: \textbf{low}}, breakable, before skip=4pt, after skip=4pt]
\textcolor{gray}{\small Oral $n=6$ (meta 5/6) | HC $n=2$ (meta 1/2) | Reject $n=3$ (meta 3/3)}

\textbf{Sample note.} \textit{Sample sizes are small (n\_oral=6, n\_reject=3, n\_hc=2); meta-review coverage is adequate (Oral 83\%, Reject 100\%) but HC is too small to support an Oral-vs-HC contrast, so confidence is set low.}

\textbf{Step-by-Step:}
\begin{enumerate}[leftmargin=14pt,topsep=2pt,itemsep=1pt]
  \item Identify the regime where single-shot perfectionism dominates: stable correct answer across perturbations, individual passes noisy.
  \item Specify the source of diversity (random sampling, structural heterogeneity, learned policy diversity) and prove or empirically show it is genuinely independent --- not relabeling.
  \item Specify the aggregation rule: weighted vote, consensus, layered cross-conditioning, or closed-form weighted aggregation --- pick the rule the error model justifies.
  \item State the error model under which aggregation provably dominates single-best (independence, voting majority, support coverage) --- without an error model, aggregation is a heuristic.
  \item Run the aggregation; report the diversity statistics that justify the assumed error model --- independence claims need induced-diversity analysis.
  \item Compare against the strongest single-pass baseline and against simpler aggregation rules --- gains over both are the threshold for novelty.
\end{enumerate}
\textbf{Success conditions} (Oral):
\begin{itemize}[leftmargin=12pt,topsep=2pt,itemsep=2pt]
  \item \textbf{Accompany the aggregation with a formal guarantee or a measurable independence argument --- a bound, a theorem, or controlled ablations that show the candidates carry complementary (not redundant) errors.} \textcolor{gray}{\footnotesize evidence: \texttt{ICLR\_2023\_0165}, \texttt{ICML\_2025\_1710}} \\ ICLR\_2023\_0165 extends weighted least squares to vector-valued outputs so the aggregation parameter is computable from source data with a provable target-error bound; ICML\_2025\_1710 runs systematic FM-by-granularity ablations to identify exactly when complementary signals exist. Both replace the intuition 'diversity should help' with evidence it does.
  \item \textbf{Re-frame the problem so that aggregation becomes the natural primitive rather than a post-hoc combination trick.} \textcolor{gray}{\footnotesize evidence: \texttt{ICLR\_2025\_1509}, \texttt{NeurIPS\_2025\_1845}, \texttt{ICML\_2023\_0475}} \\ Mixture-of-Agents recasts inference as a collaboration medium where competing outputs are first-class inputs; COMPASS recasts RL performance ceilings as inference-sampling failures rather than training failures; Instant Soup recasts pruning-plus-ensembling as a single weight-averaging pass. Each reframing relocates the novelty from 'we combined k models' to 'the object we aggregate is different.'
  \item \textbf{Decompose the task before aggregating, assigning different generators to sub-problems matched to their competence.} \textcolor{gray}{\footnotesize evidence: \texttt{ICLR\_2025\_1360}, \texttt{ICLR\_2023\_0165}} \\ BIRD splits LLM work into factor enumeration vs. conditional probability estimation before assembling a Bayesian network; the UDA paper extends aggregation to vector-valued outputs so each dimension can be handled by a different system. Structured division of labor is what keeps aggregation from just averaging noise.
  \item \textbf{Empirically establish that diversity is genuine and load-bearing by varying candidate sources and showing the aggregate exceeds the strongest individual.} \textcolor{gray}{\footnotesize evidence: \texttt{ICLR\_2025\_1509}, \texttt{ICML\_2025\_1710}} \\ MoA shows performance improves even when peers are weaker than the synthesizer, proving the gain is not 'pick the best model'; the FM subset-selection paper shows the multi-FM ranker beats any single FM across fine-grained datasets. Without this evidence reviewers suspect the ensemble is just oracle selection in disguise.
  \item \textbf{Account explicitly for the k-fold candidate budget, either by bounding it, amortizing it, or treating compute as a first-class experimental axis.} \textcolor{gray}{\footnotesize evidence: \texttt{ICML\_2023\_0475}, \texttt{NeurIPS\_2025\_1845}} \\ Instant Soup collapses k training passes to 1 before claiming ensemble quality; COMPASS treats inference compute as an explicit axis against which ceiling-breaking is demonstrated. Orals preempt the 'you just used k$\times$ compute' critique that sinks many Reject papers.
\end{itemize}
\textbf{Failure modes} (Reject):
\begin{itemize}[leftmargin=12pt,topsep=2pt,itemsep=2pt]
  \item \textbf{Aggregation is asserted to induce diversity but the mechanism producing uncorrelated candidates is neither proven nor measured.} \textcolor{gray}{\footnotesize evidence: \texttt{ICML\_2025\_1734}, \texttt{ICLR\_2023\_0401}} \\ Jakiro's reviewers explicitly complained it 'lacks justification and convincing analysis why MoE would lead to more diversity benefits'; the retrieval hybrid paper likewise combined dense+sparse without quantifying independence between the two signals. Without a diversity measurement, the method looks like it could reduce to a single generator with extra parameters.
  \item \textbf{Additive stacking of off-the-shelf components where the aggregate gain cannot be distinguished from the strongest individual component.} \textcolor{gray}{\footnotesize evidence: \texttt{ICLR\_2023\_0401}, \texttt{ICML\_2025\_1734}} \\ Reviewers of the search-agent retrieval paper found the contribution over prior iterative retrieval small because dense+sparse+reranker are each pre-validated; similarly Jakiro stacks MoE onto EAGLE draft models. In both cases the aggregation is a composition, not a reframing, so reviewers credit the parts, not the ensemble.
  \item \textbf{No cost/latency accounting for the k-fold candidate budget that aggregation implies.} \textcolor{gray}{\footnotesize evidence: \texttt{ICLR\_2023\_0401}, \texttt{ICML\_2025\_1734}, \texttt{NeurIPS\_2025\_1832}} \\ Reject meta-reviews consistently demand latency discussion (retrieval paper), efficiency and memory ablations vs. non-ensemble baselines (Jakiro), or deeper insight beyond cost savings (AdaCoT). Diverse-candidate methods multiply compute, and reviewers refuse to grant acceptance when that multiplier is hidden.
  \item \textbf{Framing the aggregation as a principled optimization while the method collapses to hyperparameter tuning.} \textcolor{gray}{\footnotesize evidence: \texttt{NeurIPS\_2025\_1832}, \texttt{ICML\_2025\_1734}} \\ AdaCoT's Pareto framing 'basically reduces to tuning penalty coefficients in an RL reward function' per the AC; Jakiro's MoE diversity claim likewise reduces to weighting expert outputs without a formal independence argument. When the scaffolding can be unwrapped back to knob-turning, reviewers classify the work as engineering, not science.
\end{itemize}
\textbf{Oral vs Reject gap.} Oral papers prove or measure that their candidates carry uncorrelated errors before claiming aggregation helps --- via a target-error bound computable from source data (ICLR\_2023\_0165), controlled ablations across many foundation models and granularities (ICML\_2025\_1710), or a reframed problem statement that makes aggregation the primitive (NeurIPS\_2025\_1845 reframes performance ceilings as inference-sampling failures; ICLR\_2025\_1509 reframes inference itself as a collaboration medium). Reject papers introduce a diversity mechanism (MoE experts, dense+sparse retrievers, adaptive triggering) but skip the 'why are these candidates actually independent?' step, leaving reviewers with a composition of known parts whose gain could come from any single component. Orals also account explicitly for the multiplicative cost of producing k candidates (COMPASS treats compute budget as a first-class axis; Instant Soup collapses k passes to 1); Rejects either hide or under-analyze latency and memory overhead. The asymmetry is structural, not rhetorical: Orals change the object being aggregated, Rejects change only the aggregator on top of a conventional pipeline.

\textbf{Oral vs HC gap.} HC sample (n=2) too small for reliable contrast.

\textbf{Reviewer expectations:}
\begin{itemize}[leftmargin=12pt,topsep=2pt,itemsep=1pt]
  \item \textit{[both]} A theoretical or mechanistic explanation for why the chosen candidate-generation scheme produces uncorrelated (or at least complementary) errors, not just that it produces many candidates.
  \item \textit{[oral\_reviews]} Controlled experiments identifying the regime where aggregation wins and where it fails (e.g., granularity, task difficulty, budget), not only headline accuracy numbers.
  \item \textit{[reject\_reviews]} Explicit accounting for the compute/latency cost of generating and scoring k candidates, with comparison against matched-compute single-path baselines.
  \item \textit{[both]} A clear articulation of what conceptual shift the aggregation enables (e.g., inference as collaboration, ceiling as sampling failure) rather than presenting the method as a composition of existing tricks.
  \item \textit{[reject\_reviews]} Differentiation from close prior work that already combined similar ingredients --- reviewers penalize heavily when the delta over a cited predecessor is unclear.
  \item \textit{[oral\_reviews]} Evidence that the aggregation result exceeds the best single candidate, not merely the average candidate, ruling out the 'oracle selection would be sufficient' objection.
\end{itemize}
\textbf{Cognitive barriers:}
\begin{itemize}[leftmargin=12pt,topsep=2pt,itemsep=1pt]
  \item The prior belief that combining imperfect candidates amplifies rather than cancels their errors --- making researchers default to 'pick the single best' rather than 'weight them all' (ICLR\_2023\_0165, ICLR\_2023\_0339).
  \item Treating powerful models as monolithic black-box reasoners rather than as replaceable components in a structured division of labor (ICLR\_2025\_1360, ICML\_2024\_1047, ICLR\_2025\_1509).
  \item Framing performance limits as training-time failures that demand bigger models or more data, rather than inference-time sampling failures that admit cheap test-time fixes (NeurIPS\_2025\_1845, ICLR\_2023\_0339).
  \item Failing to connect separately-developed research threads --- the methodology often requires fusing two existing primitives that nobody had linked, e.g. model soups + lottery tickets, or CoT + majority vote (ICML\_2023\_0475, ICLR\_2023\_0339).
\end{itemize}
\textbf{Representative examples:}
\begin{itemize}[leftmargin=12pt,topsep=2pt,itemsep=1pt]
  \item \textbf{[Oral]} \texttt{ICLR\_2023\_0165} \textemdash\ Turns 'which model do I pick without a validation signal?' into a closed-form weighted aggregation with a provable target-error bound --- the only paper in the set that derives an aggregation guarantee rather than asserting one.
  \\ \textcolor{gray}{\footnotesize \textit{Lesson:} Convert 'which model do I pick' into a closed-form weighted aggregation with a provable target-error bound --- the only paper in the set that derives an aggregation guarantee rather than asserting one.}
  \item \textbf{[Oral]} \texttt{ICLR\_2025\_1509} \textemdash\ Reframes inference from a single-pass, single-model process into a layered cross-output conditioning pipeline, making 'feed competing outputs back as context' a first-class architectural primitive rather than a post-hoc combination trick.
  \\ \textcolor{gray}{\footnotesize \textit{Lesson:} Reframe inference as a layered conditioning pipeline --- feeding competing outputs back as context becomes a first-class architectural primitive, not a post-hoc combination.}
  \item \textbf{[Oral]} \texttt{NeurIPS\_2025\_1845} \textemdash\ Recasts RL performance ceilings as inference-time sampling failures (the right answer lives in the policy's support but not its mode), justifying diversity-driven test-time search without additional training.
  \\ \textcolor{gray}{\footnotesize \textit{Lesson:} Recast performance ceilings as inference-time sampling failures --- when the right answer is in the policy's support but not its mode, diversity-driven test-time search is justified without retraining.}
  \item \textbf{[Oral]} \texttt{ICML\_2023\_0475} \textemdash\ Fuses two previously disconnected threads (model soups and lottery tickets) to collapse an O(k) multi-retraining ensemble into an O(1) single-pass procedure that still preserves candidate diversity.
  \\ \textcolor{gray}{\footnotesize \textit{Lesson:} Fuse two previously-disconnected threads to collapse multi-retraining into single-pass --- the contribution is the structural fusion, not either component alone.}
  \item \textbf{[Reject]} \texttt{ICML\_2025\_1734} \textemdash\ Claims MoE experts generate statistically independent speculative-decoding candidates but provides no theoretical or empirical analysis of the induced diversity, exactly the gap reviewers flagged as fatal.
  \\ \textcolor{gray}{\footnotesize \textit{Lesson:} Independence claims need analysis --- 'MoE experts generate statistically independent candidates' without theoretical or empirical support is the gap reviewers flag as fatal.}
  \item \textbf{[Reject]} \texttt{ICLR\_2023\_0401} \textemdash\ Combines dense + sparse retrieval with a cross-encoder reranker --- three pre-validated components whose additive composition reviewers found indistinguishable from prior iterative-retrieval work.
  \\ \textcolor{gray}{\footnotesize \textit{Lesson:} Additive composition of pre-validated components reads as iterative-retrieval reuse --- the aggregation must produce a structural advantage the components don't have individually.}
\end{itemize}
\end{tcolorbox}

\begin{tcolorbox}[colback=bgskill, colframe=cblue!70, title={\textbf{Localize via Causal Probing} \textcolor{gray}{\small\texttt{mechanistic\_probing\_with\_intervention}} \hfill confidence: \textbf{low}}, breakable, before skip=4pt, after skip=4pt]
\textcolor{gray}{\small Oral $n=29$ (meta 22/29) | HC $n=4$ (meta 4/4) | Reject $n=20$ (meta 20/20)}

\textbf{Sample note.} \textit{n\_oral=29 and n\_reject=20 with 76\% and 100\% meta coverage are robust, but HC n=4 falls below the <5 threshold needed for a reliable Oral-vs-HC contrast, so oral\_hc\_gap is deliberately left blank.}

\textbf{Step-by-Step:}
\begin{enumerate}[leftmargin=14pt,topsep=2pt,itemsep=1pt]
  \item Hypothesize an internal mechanism M: which component (layer, head, neuron, activation direction, circuit) is responsible for the observed capability or failure?
  \item Pick the right unit (SAE feature, FFN value neuron, attention head, layer slice) --- the level at which M can be intervened on cleanly.
  \item Design a targeted intervention I (ablation, activation patching, weight surgery, controlled input perturbation, causal probing) --- I must single out M, not bundled co-located components.
  \item Predict the discriminating outcome BEFORE running I: what does the behavior look like if M is responsible vs. correlate? The prediction must be falsifiable.
  \item Run I; measure the behavioral change; rule out spurious co-located effects with control conditions.
  \item Exploit the attribution as an actionable artifact: targeted edit, capability transfer, bias removal, or interpretability tool --- descriptive findings without exploitation invite the 'descriptive rather than explanatory' critique.
\end{enumerate}
\textbf{Success conditions} (Oral):
\begin{itemize}[leftmargin=12pt,topsep=2pt,itemsep=2pt]
  \item \textbf{Pair the probe with a causal intervention (activation patching, ablation, editing, injection) that flips or reproduces the behavior, never stopping at probe accuracy alone.} \textcolor{gray}{\footnotesize evidence: \texttt{ICLR\_2023\_0227}, \texttt{ICLR\_2025\_1457}, \texttt{ICLR\_2025\_1489}, \texttt{ICLR\_2025\_1526}, \texttt{ICLR\_2024\_0836}, \texttt{ICLR\_2025\_1586}} \\ Oral papers treat probe-predicted structure as a hypothesis; the paper is built around the intervention experiment that converts 'encoded' into 'used'. Reviewers explicitly praise this causal step (e.g., ICLR\_2025\_1489's 'causal interventions' and ICLR\_2023\_0227's 'interventional experiments').
  \item \textbf{Convert the localized mechanism into an actionable downstream artifact --- a controllable knob, a targeted edit, a safety ablation, or a transferable vector --- rather than ending with a descriptive finding.} \textcolor{gray}{\footnotesize evidence: \texttt{ICML\_2024\_1051}, \texttt{ICML\_2023\_0430}, \texttt{ICLR\_2025\_1526}, \texttt{ICLR\_2025\_1405}, \texttt{ICLR\_2024\_0836}, \texttt{ICLR\_2025\_1586}} \\ Each Oral paper exits with a concrete capability: value-neuron editing fixes arithmetic, SAHARA zeroes <0.006\% of heads to break safety, SAE directions steer refusal, summed function vectors trigger behavior in unrelated contexts. The intervention is the product, not a diagnostic.
  \item \textbf{Ground the analysis in a controlled setting with known ground-truth structure (synthetic task, established continuous scale, simpler non-neural system) so the probe's claims are falsifiable.} \textcolor{gray}{\footnotesize evidence: \texttt{ICLR\_2023\_0227}, \texttt{ICLR\_2023\_0397}, \texttt{ICLR\_2024\_0969}, \texttt{ICLR\_2025\_1489}, \texttt{ICML\_2025\_1703}, \texttt{ICLR\_2024\_0980}} \\ Othello, linear-ICL, boolean ICL, modular arithmetic, and DW-NOMINATE ideology scales all provide external ground truth against which probe predictions and intervention effects can be cross-checked, eliminating the 'probe is just memorizing' escape hatch.
  \item \textbf{Triangulate with independent evidence streams --- constructive proof, behavioral discrimination across competing hypotheses, and mechanistic decoding --- rather than relying on a single probing technique.} \textcolor{gray}{\footnotesize evidence: \texttt{ICLR\_2023\_0397}, \texttt{ICLR\_2023\_0402}, \texttt{ICLR\_2024\_0969}, \texttt{ICLR\_2024\_0836}, \texttt{NeurIPS\_2024\_1229}} \\ Oral meta-reviews repeatedly praise combining 'theoretical foundation-laying and empirical experiments specifically designed to test the assumptions of those theories' (ICLR\_2024\_0969) and the design of discriminating experiments (ICLR\_2023\_0397/0402). Single-method probing is not enough at Oral bar.
  \item \textbf{Adopt an interpretable unit of analysis (SAE feature, concept-aligned neuron, routing weight, established scale) before running the circuit/attribution analysis, rather than analyzing raw polysemantic neurons or heads.} \textcolor{gray}{\footnotesize evidence: \texttt{ICLR\_2025\_1586}, \texttt{ICLR\_2025\_1405}, \texttt{ICML\_2025\_1711}, \texttt{ICLR\_2025\_1489}} \\ Reviewers for ICLR\_2025\_1586 explicitly reward the move from 'highly polysemantic' units to 'monosemantic features' as an atomic substrate. The same pattern drives the protein-SAE HC paper and the political-representation Oral: pick the right unit first.
  \item \textbf{Close the loop by showing the intervention transfers or generalizes --- across contexts, models, or tasks --- not just within the identification set.} \textcolor{gray}{\footnotesize evidence: \texttt{ICLR\_2024\_0836}, \texttt{ICLR\_2025\_1405}, \texttt{ICLR\_2025\_1489}, \texttt{ICML\_2025\_1703}} \\ Function vectors injected into unrelated contexts, SAE directions transferred from base to chat model, grokking reproduced in RFM --- these transfer demonstrations distinguish 'identified' from 'causally sufficient'. Meta-reviewers consistently cite these as strengths.
\end{itemize}
\textbf{Failure modes} (Reject):
\begin{itemize}[leftmargin=12pt,topsep=2pt,itemsep=2pt]
  \item \textbf{Descriptive probing without a causal account --- characterizing what is encoded internally and stopping there, leaving reviewers to call the work 'empirical characterization without a causal account of how the behavior emerges'.} \textcolor{gray}{\footnotesize evidence: \texttt{ICML\_2025\_1819}, \texttt{NeurIPS\_2024\_1149}, \texttt{ICML\_2025\_1721}} \\ Chess look-ahead probes quantified depth but had no mechanistic explanation; stochastic-parroting work used linear probes without mechanism; topological signatures were 'reporting of results ... not very well-motivated'. Meta-reviews diagnose this as descriptive, not explanatory.
  \item \textbf{Transplanting an existing attribution or probing method into a new architecture/domain with only engineering adaptation, producing 'limited novelty' and no domain-specific insight.} \textcolor{gray}{\footnotesize evidence: \texttt{ICLR\_2023\_0230}, \texttt{ICLR\_2023\_0161}, \texttt{ICML\_2025\_1725}, \texttt{NeurIPS\_2025\_1940}} \\ LRP$\to$GNN-RL, saliency$\to$medical imaging, Shapley$\to$proteins, and explainer$\to$heterogeneous graphs were all rejected with essentially the same complaint: the method was not rethought for the new setting and contributes no mechanistic claim reviewers could take home.
  \item \textbf{Relying on narrow or architecture-specific hooks (one normalization layer, one relation type, one dataset family) and failing to test generality, so reviewers question whether the finding is a property of the model or the probe.} \textcolor{gray}{\footnotesize evidence: \texttt{ICLR\_2023\_0255}, \texttt{NeurIPS\_2023\_0670}, \texttt{NeurIPS\_2023\_0644}} \\ Channel-awareness required CCBM, word2vec-style arithmetic only held on one-to-one relations, GOAt used too few datasets. Each meta-review flagged generality as the deciding weakness.
  \item \textbf{Formal or axiomatic critiques of an attribution method without a constructive replacement mechanism reviewers can use, leading to 'the title promises a bit more than is provided'.} \textcolor{gray}{\footnotesize evidence: \texttt{NeurIPS\_2023\_0577}, \texttt{NeurIPS\_2023\_0692}, \texttt{ICLR\_2023\_0187}} \\ Refutations of Shapley, formal-logic attribution, and Fourier-consistency attribution were rejected because the theory 'reiterates the definitions' or leaves practical guidance unclear. Critique alone, without a usable successor method, does not clear the bar.
  \item \textbf{Proposing a new probing/attribution recipe whose faithfulness is validated only by heuristic or subjective measures (heatmap quality, human preference) rather than against a downstream task with ground truth.} \textcolor{gray}{\footnotesize evidence: \texttt{ICLR\_2024\_0802}, \texttt{ICLR\_2024\_0928}, \texttt{NeurIPS\_2023\_0648}} \\ DAME's human study asked for subjective 'quality'; ProtoNMF lost discriminativeness; frequency-filtering attribution lacked formal grounding. Reviewers repeatedly demand an objective downstream evaluation.
  \item \textbf{Operationalizing a grand construct (self-awareness, parroting, general ideology) through a single, under-specified behavioral test without internal-mechanism evidence to back the claim.} \textcolor{gray}{\footnotesize evidence: \texttt{NeurIPS\_2024\_1150}, \texttt{NeurIPS\_2024\_1149}} \\ Both papers gesture at mechanistic/representational questions but deliver only a thin behavioral probe; meta-reviewers found the soundness and rigor of the central test insufficient.
\end{itemize}
\textbf{Oral vs Reject gap.} Oral papers always execute a probe--intervention--exploitation arc: the probe identifies candidate components, a causal intervention (activation patching, targeted ablation, editing, injection, or swap into a different context) either breaks or reproduces the behavior, and the work ends with an actionable artifact (an edit that fixes arithmetic, a head-set whose ablation jailbreaks safety, a transferable function vector, a controllable ideology axis). Reject papers stop at one of the earlier stages --- they characterize internal structure, propose a new attribution recipe, or critique an existing one, but either do not run the intervention that would convert correlation into causation (ICML\_2025\_1819, NeurIPS\_2024\_1149) or propose a technique whose only validation is heatmap-quality or a subjective human rating (ICLR\_2024\_0802, NeurIPS\_2023\_0648). Oral work also insists on triangulation --- constructive proof plus behavioral discrimination plus mechanistic decoding, or probe plus intervention plus transfer --- whereas Reject submissions typically have a single evidence stream and rely on the method's theoretical pedigree (Shapley, LRP, persistent homology, Fourier) to carry plausibility. Finally, Oral papers pick their unit of analysis deliberately (SAE feature, routing weight, concept neuron tied to an external scale) before probing; Reject papers analyze polysemantic raw units or architecture-specific hooks and get hit on generality.

\textbf{Oral vs HC gap.} HC sample (n=4) too small for reliable contrast.

\textbf{Reviewer expectations:}
\begin{itemize}[leftmargin=12pt,topsep=2pt,itemsep=1pt]
  \item \textit{[oral\_reviews]} Demand an interventional experiment that demonstrates causal --- not merely correlational --- use of the identified representation; probe accuracy alone does not satisfy them.
  \item \textit{[oral\_reviews]} Expect the paper to exit with a concrete downstream demonstration (edit, ablation effect, transfer, behavioral control) rather than a descriptive map.
  \item \textit{[reject\_reviews]} Require evidence that the finding generalizes beyond one architecture, one dataset, or one relation type; narrow scope is the single most common rejection reason.
  \item \textit{[reject\_reviews]} Reject papers that port a classic attribution/probing method to a new setting without rethinking what the new setting changes about the method's validity.
  \item \textit{[oral\_reviews]} Praise work that pairs theoretical construction (proof, formal analysis, controlled synthetic domain) with empirical measurement specifically designed to test the theory's assumptions.
  \item \textit{[reject\_reviews]} Refuse purely subjective evaluations of explanation quality; demand ground-truth or downstream-task metrics that falsify the attribution's faithfulness claims.
\end{itemize}
\textbf{Cognitive barriers:}
\begin{itemize}[leftmargin=12pt,topsep=2pt,itemsep=1pt]
  \item Novelty bias: researchers frame an impressive model behavior as a new phenomenon and ask 'what new thing is it doing?' instead of the harder, more deflationary 'is it re-executing something we already understand?' (ICLR\_2023\_0402).
  \item The 'distributed is inevitable' assumption: capabilities in large models are presumed to be smeared across all components, which blocks the hypothesis that a minimal parameter subset could mediate them and makes targeted editing seem implausible (ICML\_2024\_1051, ICLR\_2025\_1526).
  \item Cross-community isolation: causal-intervention tools, sparse decompositions, unlearning, influence functions, and counterfactual reasoning each live in separate sub-literatures, so the composition that Oral papers rely on (e.g., SAE $\times$ circuit discovery, unlearning $\times$ baseline selection) is non-obvious without bridging two methodological traditions (ICLR\_2025\_1586, ICLR\_2025\_1639).
  \item False-confidence barrier: a method with strong theoretical guarantees in its home domain (Shapley, influence functions, LRP) inherits unearned trust when moved to interpretability, obscuring the need to re-examine whether its axioms still apply (ICML\_2024\_1114, NeurIPS\_2023\_0577).
\end{itemize}
\textbf{Representative examples:}
\begin{itemize}[leftmargin=12pt,topsep=2pt,itemsep=1pt]
  \item \textbf{[Oral]} \texttt{ICLR\_2023\_0227} \textemdash\ Canonical demonstration of the probe$\to$intervention$\to$behavioral-change loop in a fully controlled synthetic task (Othello), giving reviewers ground-truth-backed causal evidence that a trained model uses --- not merely encodes --- its internal world state.
  \\ \textcolor{gray}{\footnotesize \textit{Lesson:} The probe $\to$ intervention $\to$ behavioral-change loop in a fully controlled synthetic task gives ground-truth-backed causal evidence --- a model uses (not merely encodes) its internal world state.}
  \item \textbf{[Oral]} \texttt{ICLR\_2025\_1586} \textemdash\ Exemplifies the 'pick the right unit first' move by composing SAE decomposition with circuit discovery, then using the interpretable circuits for SHIFT-style bias ablation --- turning interpretability into targeted control.
  \\ \textcolor{gray}{\footnotesize \textit{Lesson:} Pick the right unit first via SAE decomposition, then compose with circuit discovery --- turning interpretability into targeted control via SHIFT-style bias ablation.}
  \item \textbf{[Oral]} \texttt{ICML\_2024\_1051} \textemdash\ Shows that arithmetic, widely assumed to be distributed, is stored in identifiable FFN value neurons and can be fixed by direct inference-time activation edits with no gradient training --- the methodology's most striking 'actionable artifact' payoff.
  \\ \textcolor{gray}{\footnotesize \textit{Lesson:} The most striking actionable-artifact payoff --- what was widely assumed distributed is in fact stored in identifiable FFN value neurons and can be edited at inference time with no gradient training.}
  \item \textbf{[Oral]} \texttt{ICLR\_2025\_1489} \textemdash\ Anchors probing to an established external continuous scale (DW-NOMINATE) and closes the loop with head-level ablation that causally steers political bias, a textbook three-step execution of the methodology.
  \\ \textcolor{gray}{\footnotesize \textit{Lesson:} Anchor probing to an established external continuous scale --- closing the loop with head-level ablation that causally steers political bias is a textbook three-step execution.}
  \item \textbf{[Reject]} \texttt{ICML\_2025\_1819} \textemdash\ Extends prior probing infrastructure to quantify chess look-ahead depth but stops at characterization --- reviewers explicitly call it 'descriptive rather than explanatory' because no mechanism is causally identified.
  \\ \textcolor{gray}{\footnotesize \textit{Lesson:} Stopping at characterization is 'descriptive rather than explanatory' --- the methodology demands a causal intervention identifying a mechanism, not a quantification.}
  \item \textbf{[Reject]} \texttt{ICLR\_2023\_0230} \textemdash\ Transplants LRP onto a GNN-based RL policy with only engineering adaptation and no mechanistic claim, hitting the 'domain port without insight' failure mode head-on.
  \\ \textcolor{gray}{\footnotesize \textit{Lesson:} Domain port without insight --- transplanting probing infrastructure with only engineering adaptation hits the 'no mechanistic claim' failure head-on.}
  \item \textbf{[Reject]} \texttt{ICML\_2025\_1721} \textemdash\ Proposes topological signatures of adversarial latent-space compression but lacks causal intervention and downstream task validation, so reviewers treat it as unmotivated reporting rather than mechanism discovery.
  \\ \textcolor{gray}{\footnotesize \textit{Lesson:} Without causal intervention and downstream task validation, mechanism reporting is 'unmotivated' --- observation is not discovery.}
\end{itemize}
\end{tcolorbox}

\begin{tcolorbox}[colback=bgskill, colframe=cblue!70, title={\textbf{Audit Evaluation via Controlled Perturbation} \textcolor{gray}{\small\texttt{evaluation\_validity\_via\_controlled\_perturbation}} \hfill confidence: \textbf{high}}, breakable, before skip=4pt, after skip=4pt]
\textcolor{gray}{\small Oral $n=46$ (meta 36/46) | HC $n=25$ (meta 20/25) | Reject $n=22$ (meta 22/22)}

\textbf{Sample note.} \textit{Sample is well-powered for both directions (n\_oral=46, n\_reject=22, full reject meta coverage); Oral meta-coverage at 36/46 (78\%) is adequate but several Oral ICML 2023/2024 papers lack meta-reviews, which slightly weakens the reviewer-voice signal for the success side.}

\textbf{Step-by-Step:}
\begin{enumerate}[leftmargin=14pt,topsep=2pt,itemsep=1pt]
  \item Suspect a confound C in benchmark B --- shortcut, leakage, memorization, measurement artifact --- name C precisely, not vaguely.
  \item Construct controlled variants of B: vary the target capability while holding C fixed, OR vary C while holding the target capability fixed.
  \item Equalize confounders the variants don't address (capacity, hyperparameter search, dataset scale) --- without equalization, 'controlled' perturbations are not controlled.
  \item Re-evaluate methods on the variants; reattribute the apparent capability gap to either the target or the confound.
  \item Validate the controlled-perturbation move on a real-data bridge --- synthetic-only audits without showing the perturbation tracks natural confounds die in review.
  \item State what changes in the field's understanding: a published consensus reverses, a benchmark is invalidated, a measurement primitive is retired. The audit's payoff is the reattribution.
\end{enumerate}
\textbf{Success conditions} (Oral):
\begin{itemize}[leftmargin=12pt,topsep=2pt,itemsep=2pt]
  \item \textbf{Run a direct ablation that toggles the suspected confound on/off while holding everything else fixed, then report the effect size of the confound itself rather than just the relative ranking change.} \textcolor{gray}{\footnotesize evidence: \texttt{NeurIPS\_2024\_1168}, \texttt{NeurIPS\_2025\_1888}, \texttt{NeurIPS\_2024\_1248}, \texttt{ICLR\_2025\_1630}} \\ Cambrian-1 freezes the LLM and varies only vision encoders; ImageNet-trained CNNs paper independently suppresses each cue rather than letting them compete; LLM Evaluators watermark generations to manipulate self-recognition; Training on the Test Task fine-tunes base models on task format and shows it absorbs most of the apparent progress. In each case the methodology is a one-factor-at-a-time intervention with a measured magnitude, not just a rank flip.
  \item \textbf{Re-run a classic, named diagnostic test or community-accepted protocol on modern systems and use the replication itself as the contribution, rather than designing a new metric.} \textcolor{gray}{\footnotesize evidence: \texttt{ICLR\_2025\_1339}, \texttt{ICLR\_2024\_0985}, \texttt{ICLR\_2025\_1404}, \texttt{ICLR\_2024\_0915}} \\ A Decade's Battle directly replicates Torralba \& Efros (2011) at modern scale; Unprocessing re-evaluates 7 years of fairness papers under one harness; Conformity ports the Asch experiment; PPBench imports validated psychometric instruments. Reviewers explicitly praise the move of carrying an external/older protocol into a saturated subfield because the protocol's prior validation buys credibility for the negative findings.
  \item \textbf{Build a parameterized template or generator that produces semantically-matched variants and report how performance degrades along a single perturbation axis (irrelevant clause, value substitution, frame switch, syntactic transform).} \textcolor{gray}{\footnotesize evidence: \texttt{ICLR\_2025\_1443}, \texttt{ICML\_2025\_1747}, \texttt{ICLR\_2025\_1412}, \texttt{ICLR\_2025\_1407}} \\ GSM-Symbolic, LLM-SRBench, VLB, and COMFORT all share the structure: define a template, generate matched variants that vary one capability-relevant dimension while preserving difficulty/semantics, then plot degradation. The template move converts an ambiguous benchmark complaint into a continuous, replottable signal that reviewers can interrogate.
  \item \textbf{Construct an adversarial 'minimal counterexample' or null baseline whose success on the benchmark logically implies the benchmark is invalid, not merely weak.} \textcolor{gray}{\footnotesize evidence: \texttt{ICLR\_2025\_1371}, \texttt{NeurIPS\_2024\_1161}, \texttt{ICML\_2024\_1066}, \texttt{ICML\_2025\_1684}} \\ Cheating Automatic Benchmarks shows a constant null-model response wins; Are We on the Right Way runs vision benchmarks text-only and shows large fractions remain solvable; MLLM-as-Judge constructs cases where known biases would dominate; EMMA designs tasks that admit no single-modality shortcut by construction. The argument has the form of a proof-by-counterexample, which is qualitatively stronger than 'the benchmark is noisy'.
  \item \textbf{Pair the audit with a constructive remedy or design rule (filtered benchmark, calibrated harness, escalation protocol, intermediate criterion) so the paper is not purely a takedown.} \textcolor{gray}{\footnotesize evidence: \texttt{ICLR\_2025\_1633}, \texttt{ICML\_2025\_1740}, \texttt{ICLR\_2024\_0985}, \texttt{ICLR\_2025\_1551}} \\ Trust or Escalate provides a selective-prediction framework with provable agreement; eSEN proposes energy-conservation as an intermediate filter that predicts utility; Unprocessing offers a unified comparison protocol; Counterfactual Feedback proposes both new metrics and a training fix. Reviewers in this set repeatedly note that the constructive contribution prevents the work from reading as 'just a critique'.
  \item \textbf{Validate the audit across many models/configurations and report the consistency of the effect, not a single dramatic point estimate.} \textcolor{gray}{\footnotesize evidence: \texttt{NeurIPS\_2024\_1325}, \texttt{ICLR\_2025\_1339}, \texttt{ICLR\_2025\_1383}, \texttt{ICLR\_2024\_0985}} \\ Reasoning Boundary spans 25 models and 4 tasks; Decade's Battle sweeps architectures and dataset configurations; Context-Parametric Inversion tests multiple model families and IFT datasets; Unprocessing evaluates thousands of configurations. Meta-reviews credit acceptance specifically to this breadth, contrasted with rejected papers that show the effect on one model or one dataset.
\end{itemize}
\textbf{Failure modes} (Reject):
\begin{itemize}[leftmargin=12pt,topsep=2pt,itemsep=2pt]
  \item \textbf{The controlled environment is fully synthetic with no demonstrated transfer to realistic data, so reviewers cannot tell whether the diagnosed confound matters in deployment.} \textcolor{gray}{\footnotesize evidence: \texttt{ICLR\_2023\_0374}, \texttt{ICLR\_2024\_0958}} \\ Blindspot Discovery's SpotCheck is synthetic-only and reviewers explicitly flag uncertainty about whether artificially-induced blindspots track natural ones; SynBench builds a class-conditional Gaussian proxy and reviewers reject the claim that it predicts behavior on real downstream tasks. Both papers are technically clean controlled-perturbation work, but the absence of a real-data anchor sinks them.
  \item \textbf{The audited methodology overlaps too closely with existing critique papers, so the controlled-perturbation move reads as a re-skinning rather than a new audit.} \textcolor{gray}{\footnotesize evidence: \texttt{ICLR\_2024\_0873}, \texttt{ICLR\_2024\_0869}, \texttt{ICLR\_2025\_1402}} \\ Are LLMs Strong Abstract Reasoners is rejected for not diverging enough from existing reasoning critiques; Compositional Decision Making is treated as a variant of prior compositionality benchmarks; Direct Judgment Preference Optimization overlaps with concurrent judge-training work. The signal: reviewers ask 'what does this audit show that prior audits did not?' and accepted papers have a sharp answer.
  \item \textbf{The proposed alternative metric or benchmark relies on a subjective ground-truth pipeline (LLM-as-judge, in-house annotation) without showing that the new ground truth is itself less biased than the one being criticized.} \textcolor{gray}{\footnotesize evidence: \texttt{NeurIPS\_2024\_1276}, \texttt{ICLR\_2024\_0956}, \texttt{NeurIPS\_2025\_1851}} \\ PiCO uses peer LLMs as judges but reviewers question the consistency assumptions; Style Over Substance evaluates only 40 GPT-4 generations with the same model; the Visual Emotion paper bootstraps labels from MLLM majority-vote while simultaneously claiming MLLMs are unreliable on this task. The audit collapses because the new measurement instrument inherits or amplifies the very pathology being diagnosed.
  \item \textbf{The confound is named and a fix is proposed but neither is convincingly isolated --- the proposed correction is entangled with other moving parts, so reviewers cannot attribute improvements to the fix.} \textcolor{gray}{\footnotesize evidence: \texttt{NeurIPS\_2024\_1291}, \texttt{ICML\_2025\_1783}, \texttt{ICLR\_2024\_0914}} \\ Removing Length Bias proposes prompt-template bias calibration but cannot disentangle it from length effects; Reliable Image Quality proposes DQA guidance but reviewers find the disparity signal entangled with feature-extractor pretraining; Provable Benefits of Preferences argues bias is the cause but the ablation logic does not cleanly rule out alternative explanations. The methodology lives or dies on isolation, and partial isolation does not pass.
  \item \textbf{Audit is restricted to a narrow benchmark or single phenomenon whose flaws are already widely suspected, so the paper confirms folk knowledge rather than uncovering a hidden confound.} \textcolor{gray}{\footnotesize evidence: \texttt{ICLR\_2025\_1653}, \texttt{NeurIPS\_2025\_1841}, \texttt{ICLR\_2023\_0387}} \\ The LRA paper audits only Long Range Arena, whose local-bias problem was already well-known; Are We Using Motion shows benchmarks lack temporal demand but on a niche slice; Conceptual Consistency proposes a consistency metric whose insight reviewers find limited. When the audited target is too narrow or its problems pre-known, controlled perturbation does not produce a community-relevant update.
  \item \textbf{The benchmark or audit lacks scale (few items, few models, few seeds) so even a clean controlled-perturbation result cannot support the broad claim it makes.} \textcolor{gray}{\footnotesize evidence: \texttt{ICLR\_2024\_0956}, \texttt{ICLR\_2025\_1422}, \texttt{NeurIPS\_2024\_1169}} \\ Style Over Substance covers 40 questions and one generator; ERiC-UP\textasciicircum{}3 has expert annotation but limited scale and primarily text validation; the Hybrid Architectures ICL paper enumerates structural permutations but reviewers find the scope insufficient for the claimed conclusions. Reviewers explicitly tie rejection to 'cannot tell if the conclusion generalizes', which is a structural risk for any controlled-perturbation paper that has high methodological cost per data point.
\end{itemize}
\textbf{Oral vs Reject gap.} Oral papers couple the controlled perturbation to a falsifiable prior claim with a named owner --- Torralba \& Efros 2011 (ICLR\_2025\_1339), seven years of fairness papers (ICLR\_2024\_0985), CLIP's assumed superiority over captioners (NeurIPS\_2023\_0658), the texture-bias claim (NeurIPS\_2025\_1888) --- and report the confound's effect size as a single headline number that reverses or substantially shrinks the prior conclusion. They run the audit across many models or configurations (25 models in NeurIPS\_2024\_1325, thousands of configs in ICLR\_2024\_0985) so reviewers can see consistency, and they offer a constructive remedy (a corrected protocol, calibrated harness, escalation framework). Reject papers, by contrast, target a vaguely-defined community practice rather than a named claim, run the perturbation on one model or a small set (40 prompts in ICLR\_2024\_0956, single-model evaluation in ICLR\_2025\_1474), and stop at 'we built a cleaner benchmark' without either falsifying a specific result or providing a usable replacement. The other recurring reject pattern is full-synthetic evaluation with no real-data bridge (ICLR\_2023\_0374, ICLR\_2024\_0958), where reviewers cannot rule out that the artificially-induced confound has no real-world counterpart.

\textbf{Oral vs HC gap.} HC papers (e.g., MMLU ICLR\_2021\_0043, GSM-Symbolic ICLR\_2025\_1443, LiveCodeBench ICLR\_2025\_1490, Chatbot Arena ICML\_2024\_1012) deliver a clean controlled-evaluation instrument that the community quickly adopts, but tend to stop at constructing the harness and reporting initial findings. Oral papers go one step further: they use the harness to overturn or sharply qualify a specific named result and articulate a methodological lesson the field is expected to internalize --- Unprocessing reverses fairness-paper conclusions (ICLR\_2024\_0985), Decade's Battle falsifies the assumption that scale solved dataset bias (ICLR\_2025\_1339), Cheating Benchmarks shows null models top leaderboards (ICLR\_2025\_1371), the Texture-bias paper inverts a canonical finding (NeurIPS\_2025\_1888). The Oral graduation move is from 'here is a better measurement instrument' to 'here is what the field believed that turns out to be wrong, demonstrated with that instrument'.

\textbf{Reviewer expectations:}
\begin{itemize}[leftmargin=12pt,topsep=2pt,itemsep=1pt]
  \item \textit{[oral\_reviews]} Show that the proposed controlled protocol changes which method or claim wins, not just that scores shift in absolute terms.
  \item \textit{[both]} Run the audit across multiple model families, scales, and datasets so a single cherry-picked condition cannot explain the effect.
  \item \textit{[oral\_reviews]} Provide a constructive output (corrected protocol, replacement metric, fine-tuning fix, or filter) rather than only a critique, so the field has something to adopt.
  \item \textit{[reject\_reviews]} Anchor the audit to a real-world or ecologically valid distribution; if the diagnosis is run on synthetic data, demonstrate transfer to natural data.
  \item \textit{[reject\_reviews]} When the audit relies on an automated judge or model annotator, validate that the judge itself does not inherit the bias being investigated.
  \item \textit{[reject\_reviews]} Explicitly contrast the diagnosed confound with prior critiques in the same area and show what is uniquely revealed by this protocol.
\end{itemize}
\textbf{Cognitive barriers:}
\begin{itemize}[leftmargin=12pt,topsep=2pt,itemsep=1pt]
  \item Practitioners take rising aggregate scores as evidence the underlying capability is improving and never re-instantiate the original diagnostic test, so the field assumes the problem is being solved while the diagnostic has never actually been re-run (ICLR\_2025\_1339, ICLR\_2025\_1443).
  \item Implementation choices in evaluation protocols are inherited from early influential papers as defaults; questioning them feels like methodological pedantry rather than a research contribution, so the option of one-factor-at-a-time ablation across protocol parameters is invisible (ICLR\_2024\_0974, ICLR\_2025\_1630).
  \item Convenience metrics (single-turn, end-to-end, scalar preference, end-state correctness) are conflated with sufficient metrics; the gap between 'easier to measure' and 'enough to measure' is suppressed by the self-reinforcing incentive of clean leaderboards (ICLR\_2024\_0899, ICLR\_2023\_0264).
  \item Benchmarks claiming to test a capability are presumed to test it because they nominally include the relevant input modality or task structure, hiding the possibility that single-modality or single-frame shortcuts dominate the score (NeurIPS\_2024\_1161, ICML\_2025\_1684, NeurIPS\_2025\_1841).
\end{itemize}
\textbf{Representative examples:}
\begin{itemize}[leftmargin=12pt,topsep=2pt,itemsep=1pt]
  \item \textbf{[Oral]} \texttt{ICLR\_2024\_0985} \textemdash\ Equalizes capacity and hyperparameter search across thousands of configurations from seven years of fairness papers and shows the simple thresholding baseline they all beat is in fact optimal --- the textbook case of a unified harness reversing a field's published consensus.
  \\ \textcolor{gray}{\footnotesize \textit{Lesson:} Equalize capacity and hyperparameter search across thousands of configurations from years of papers --- show the simple thresholding baseline they all beat is in fact optimal. A unified harness can reverse a field's published consensus.}
  \item \textbf{[Oral]} \texttt{ICLR\_2025\_1630} \textemdash\ Identifies test-task fine-tuning as a previously-unacknowledged confound, fine-tunes older base models on task-format data to ablate it, and shows the practice eats most of the apparent emergent-capability gap --- a clean isolate-and-quantify-the-confound move.
  \\ \textcolor{gray}{\footnotesize \textit{Lesson:} Identify a previously-unacknowledged confound, fine-tune older base models on task-format data to ablate it --- most of the apparent emergent-capability gap dissolves.}
  \item \textbf{[Oral]} \texttt{ICLR\_2025\_1371} \textemdash\ Constructs the minimal adversarial example --- a constant null-model response --- and shows it tops AlpacaEval-style leaderboards, using a single counterexample to invalidate the evaluation paradigm rather than a metric debate.
  \\ \textcolor{gray}{\footnotesize \textit{Lesson:} A single counterexample (constant null response tops the leaderboard) invalidates an evaluation paradigm faster than a metric debate.}
  \item \textbf{[Oral]} \texttt{NeurIPS\_2025\_1888} \textemdash\ Replaces the cue-conflict paradigm with independent suppression of shape, texture, and color, separating cue salience from feature reliance and overturning the canonical 'CNNs are texture-biased' finding through a cleaner experimental design.
  \\ \textcolor{gray}{\footnotesize \textit{Lesson:} Replace cue-conflict with independent suppression --- separating cue salience from feature reliance overturns canonical findings via cleaner experimental design.}
  \item \textbf{[Reject]} \texttt{ICLR\_2023\_0374} \textemdash\ Builds a synthetic ground-truth blindspot benchmark with the right intent, but reviewers reject because the artificially-induced blindspots are never shown to track natural ones --- the controlled-perturbation move dies without a real-data bridge.
  \\ \textcolor{gray}{\footnotesize \textit{Lesson:} The controlled-perturbation move dies without a real-data bridge --- synthetic blindspots must be shown to track natural blindspots, not just exist.}
  \item \textbf{[Reject]} \texttt{ICLR\_2024\_0956} \textemdash\ Decomposes evaluator preferences along multiple style axes --- methodologically the right idea --- but runs the entire audit on 40 GPT-4 generations evaluated by GPT-4 and Claude-1, so reviewers cannot tell if the bias is real or a single-model artifact.
  \\ \textcolor{gray}{\footnotesize \textit{Lesson:} Insufficient evaluation scope when the methodology's central claim is bias detection --- single-model evaluators undermine the bias claim itself.}
  \item \textbf{[Reject]} \texttt{ICLR\_2025\_1653} \textemdash\ Audits Long Range Arena and shows simple training tricks match specialized architectures, but the diagnosed pathology (LRA's local biases) was already widely suspected and the audit does not generalize beyond a single benchmark.
  \\ \textcolor{gray}{\footnotesize \textit{Lesson:} If the diagnosed pathology was already widely suspected and the audit doesn't generalize beyond a single benchmark, the contribution reads as confirmatory rather than novel.}
\end{itemize}
\end{tcolorbox}

\begin{tcolorbox}[colback=bgskill, colframe=cblue!70, title={\textbf{Sharpen a Single Analytic Step} \textcolor{gray}{\small\texttt{analytic\_refinement\_of\_loose\_bounds}} \hfill confidence: \textbf{medium}}, breakable, before skip=4pt, after skip=4pt]
\textcolor{gray}{\small Oral $n=60$ (meta 41/60) | HC $n=0$ (meta 0/0) | Reject $n=33$ (meta 33/33)}

\textbf{Sample note.} \textit{Strong sample on the Oral/Reject contrast (n\_oral=60 with 68\% meta coverage, n\_reject=33 with 100\% meta coverage), but n\_hc=0 means no Oral-vs-HC graduation contrast can be drawn.}

\textbf{Step-by-Step:}
\begin{enumerate}[leftmargin=14pt,topsep=2pt,itemsep=1pt]
  \item Locate the loose step L in the existing analysis: worst-case argument, unused property of the function class, missing concentration tool, coarse parameter (diameter, global norm).
  \item Identify the structural source of the slack --- why is L loose? What property of the actual problem is L not exploiting?
  \item Construct a sharper substitute L*: tighter inequality, matched lower bound, refined concentration, contour integration, Poisson process where independence is exact, or a sequence-specific replacement.
  \item Propagate L* through the existing proof --- the rest of the analysis remains valid; the sharpening localizes to the substituted step.
  \item Match with a lower bound when possible --- sharpness is convincing only when paired with a tightness witness (matched complexity, lower bound, separation).
  \item Show the improved rate has practical consequence: a regime that becomes feasible, a separation previously unreachable, a parameter that can be replaced with a tighter functional one.
\end{enumerate}
\textbf{Success conditions} (Oral):
\begin{itemize}[leftmargin=12pt,topsep=2pt,itemsep=2pt]
  \item \textbf{Import the sharper substitute from an adjacent mathematical subfield, not from within the methodology's home literature, so the new bound carries genuinely new analytical machinery.} \textcolor{gray}{\footnotesize evidence: \texttt{NeurIPS\_2023\_0660}, \texttt{NeurIPS\_2023\_0715}, \texttt{NeurIPS\_2024\_1307}, \texttt{NeurIPS\_2025\_1932}, \texttt{ICML\_2024\_1045}, \texttt{ICML\_2025\_1765}} \\ Across orals the loose step is replaced by a tool the relevant community had not used: Ville-style martingale mixtures for linear bandits, Poisson processes for explainable k-medians, Hilbert-space Freedman for distributional TD, contour integration for spectral perturbation, communication-complexity lower bounds for FlashAttention, classical abstract convexity for clipping. The 'sharpening' is recognized as conceptually deep precisely because it bridges literatures.
  \item \textbf{Pair the tighter upper bound with a matching lower bound or adversarial hard instance that proves the new bound is tight, converting 'sharper' into 'optimal'.} \textcolor{gray}{\footnotesize evidence: \texttt{NeurIPS\_2024\_1306}, \texttt{NeurIPS\_2025\_1911}, \texttt{ICLR\_2025\_1614}, \texttt{ICML\_2023\_0555}, \texttt{NeurIPS\_2023\_0699}, \texttt{NeurIPS\_2024\_1317}} \\ Oral meta-reviews use 'tight', 'matching', 'optimal', 'closes a long-standing open problem' as the praise vocabulary; the lower bound is what licenses these phrases. ICML\_2023\_0555 and ICLR\_2025\_1614 in particular show that constructing regime-specific hard instances is the difference between a tighter inequality and a settled question.
  \item \textbf{Make the refinement cascade --- show the same single substitution propagates through the existing analysis to improve multiple downstream theorems with minimal further work.} \textcolor{gray}{\footnotesize evidence: \texttt{ICML\_2025\_1728}, \texttt{NeurIPS\_2025\_1890}, \texttt{NeurIPS\_2025\_1932}, \texttt{ICML\_2023\_0528}} \\ ICML\_2025\_1728's tighter posterior-variance bound 'unlocks simultaneous optimality across multiple regimes'; NeurIPS\_2025\_1890 substitutes algorithm-specific MIG into the existing proof and re-derives improved regret in three settings; NeurIPS\_2025\_1932's contour-integration bound feeds straight into private PCA. A one-shot constant improvement is rarely enough --- the AC needs to see leverage.
  \item \textbf{Identify that the loose step is a worst-case proxy the algorithm provably avoids, then bound the algorithm-specific quantity instead.} \textcolor{gray}{\footnotesize evidence: \texttt{NeurIPS\_2025\_1890}, \texttt{ICML\_2024\_1111}, \texttt{ICML\_2024\_1041}, \texttt{NeurIPS\_2024\_1306}} \\ Across these orals the existing proof bounded a worst-case object (worst-case MIG, worst-case variance, worst-case clipping behavior, diameter) when the algorithm actually behaves more favorably; the fix is to track what the algorithm does, not what it could do. This reframes the methodology from 'find a tighter inequality' to 'find a structural property of the algorithm prior analyses ignored'.
  \item \textbf{Use the sharper analysis to either (a) drop a restrictive assumption that prevented prior bounds from applying in practice, or (b) enable a new algorithmic variant whose existence is justified by the new proof.} \textcolor{gray}{\footnotesize evidence: \texttt{ICLR\_2023\_0334}, \texttt{ICLR\_2022\_0109}, \texttt{NeurIPS\_2023\_0731}, \texttt{ICLR\_2025\_1367}} \\ ICLR\_2023\_0334 drops the L$\infty$/functional-inequality assumptions to make diffusion-model theory cover the practical regime; ICLR\_2022\_0109 derives tight shuffling-SGD bounds and uses the proof to design a synchronized-shuffling variant; NeurIPS\_2023\_0731 closes a sample-complexity gap and shows smoothing as the algorithmic preprocessing that achieves it. The sharpening pays a non-analytical dividend.
  \item \textbf{State the assumption being relaxed (or the parameter being replaced) as the headline of the paper, not as a hidden technical condition.} \textcolor{gray}{\footnotesize evidence: \texttt{ICLR\_2023\_0334}, \texttt{NeurIPS\_2024\_1306}, \texttt{ICML\_2025\_1728}, \texttt{ICML\_2025\_1765}} \\ Oral titles and abstracts foreground the relaxation ('minimal data assumptions', 'span-based', 'tighter maximum posterior variance', '(L0,L1)-smoothness'). Reject meta-reviews repeatedly criticize the opposite: hidden conditions buried in Section 3 that make the claim narrower than the title implies.
\end{itemize}
\textbf{Failure modes} (Reject):
\begin{itemize}[leftmargin=12pt,topsep=2pt,itemsep=2pt]
  \item \textbf{Re-derives a result that follows as a corollary of (or is already standard in) prior work the authors did not survey, so the 'sharpening' is not new.} \textcolor{gray}{\footnotesize evidence: \texttt{ICLR\_2023\_0162}, \texttt{ICLR\_2023\_0245}, \texttt{ICLR\_2023\_0257}, \texttt{NeurIPS\_2025\_1877}} \\ Reject meta-reviews repeatedly point to a specific prior paper from which the main theorem is derivable, or note that the operator/inequality properties being 'derived' are textbook in the relevant subliterature; without locating exactly which step in prior analyses is loose, papers default to re-proving rather than tightening.
  \item \textbf{The 'sharper' bound carries a hidden assumption whose presence makes the improvement vacuous in the headline regime.} \textcolor{gray}{\footnotesize evidence: \texttt{ICLR\_2023\_0318}, \texttt{ICLR\_2024\_0953}, \texttt{ICML\_2025\_1674}} \\ Meta-reviews flag that the convergence is only to a small noise floor (not zero), or that constants hidden in big-O blow up at the parameter setting that makes the result interesting, or that an unverifiable local-convexity assumption silently does the work; the AC reads the proof and sees the sharpening did not actually cascade through.
  \item \textbf{Adapts an existing proof technique to a new algorithm/setting without controlling the new error term that the substitution introduces.} \textcolor{gray}{\footnotesize evidence: \texttt{ICLR\_2024\_0834}, \texttt{NeurIPS\_2024\_1182}, \texttt{NeurIPS\_2024\_1295}} \\ Reviewers explicitly flag that the regret/convergence analysis 'does not account for' the gap between the approximation actually tracked by the algorithm and the exact object the proof refers to (e.g., low-rank sketch vs. true SVD, factored second moment vs. full Adam state, retraction-free iterates vs. on-manifold iterates). The sharpening leaves a hole the authors do not close.
  \item \textbf{Treats the substitution as routine technique transfer rather than as a genuine analytical innovation, leaving reviewers unable to identify the new tool.} \textcolor{gray}{\footnotesize evidence: \texttt{ICLR\_2023\_0197}, \texttt{ICML\_2025\_1694}, \texttt{ICML\_2025\_1753}, \texttt{NeurIPS\_2023\_0586}} \\ ACs describe the work as 'incremental', 'fairly straightforward given the literature', or 'BCD applied in a standard way'; without a named new mathematical object (a new dimension, a new inequality, a new metric) reviewers see assembly rather than refinement.
  \item \textbf{Lower bounds, matching constructions, or supporting empirics are absent, so the claimed tightness is untestable.} \textcolor{gray}{\footnotesize evidence: \texttt{NeurIPS\_2025\_1897}, \texttt{ICLR\_2025\_1428}, \texttt{ICLR\_2025\_1636}, \texttt{ICML\_2025\_1752}} \\ Reject meta-reviews complain about no formal theorem statements, no comparison to existing rates that would expose whether the bound is actually tighter, or contrived/toy settings where the claimed gap appears. Sharpening without a matching lower bound or quantified prior-art comparison reads as a different proof, not a better one.
  \item \textbf{Overclaims scope --- the sharpening only fires in a narrow regime that is buried in the abstract.} \textcolor{gray}{\footnotesize evidence: \texttt{ICLR\_2023\_0163}, \texttt{ICLR\_2023\_0257}, \texttt{ICLR\_2023\_0318}, \texttt{ICLR\_2024\_0953}} \\ ACs explicitly call out 'overselling' --- titles and intros promise general improvements (over Adam, over SGD, over weight decay) but the proof requires sampling without replacement, finite-sum structure, or a coordinate-wise smoothness condition that is buried in Section 3.
\end{itemize}
\textbf{Oral vs Reject gap.} Oral papers name the loose step explicitly and import a sharper substitute from a *different* mathematical subfield (Poisson processes for union bounds in NeurIPS\_2023\_0715, contour integration for spectral perturbation in NeurIPS\_2025\_1932, Freedman's inequality in Hilbert space for distributional TD in NeurIPS\_2024\_1307, communication-complexity lower bounds for FlashAttention in ICML\_2024\_1045, martingale-mixture tail bounds for linear bandits in NeurIPS\_2023\_0660); rejects substitute one inequality for another from the *same* literature with no new tool. Orals also pair the sharper upper bound with a matching lower bound or hard-instance construction (NeurIPS\_2024\_1306, NeurIPS\_2025\_1911, ICLR\_2025\_1614, ICML\_2023\_0555), turning 'tighter' into 'optimal'; rejects rarely produce lower bounds, leaving reviewers unable to verify tightness. Orals show the refinement cascading to multiple downstream theorems with minimal modification (ICML\_2025\_1728 propagates one tighter posterior-variance bound through three regimes; NeurIPS\_2025\_1890 swaps worst-case MIG for an algorithm-specific bound and reuses the prior regret framework); rejects fix one step but introduce a new uncontrolled error term (ICLR\_2024\_0834's sketch tracking error, NeurIPS\_2024\_1295's missing self-adjointness step, ICLR\_2023\_0318's exploding constants when $\beta$$_2$$\to$1). Finally, orals state the assumption being relaxed as the headline (ICLR\_2023\_0334 advertises 'minimal data assumptions'); rejects bury the new restrictive assumption in Section 3 and oversell the headline.

\textbf{Oral vs HC gap.} HC sample (n=0) too small for reliable contrast.

\textbf{Reviewer expectations:}
\begin{itemize}[leftmargin=12pt,topsep=2pt,itemsep=1pt]
  \item \textit{[oral\_reviews]} Pair the sharper upper bound with a matching lower bound or adversarial-instance construction so the result is provably tight, not merely tighter.
  \item \textit{[oral\_reviews]} Show that the refinement cascades --- the same tightening should improve multiple downstream theorems or unlock simultaneous optimality across regimes (not just one corollary).
  \item \textit{[oral\_reviews]} Name the loose step explicitly and trace its provenance to a specific prior paper's lemma; reviewers credit the work for 'closing a long-standing gap' only when the gap and its source are made concrete.
  \item \textit{[reject\_reviews]} Survey the relevant prior literature exhaustively --- reject meta-reviews repeatedly identify a specific antecedent paper that already implies the result, and authors often admit they were unaware of it.
  \item \textit{[reject\_reviews]} Make any new restrictive assumption (without-replacement sampling, local convexity around the optimum, sketch-tracking accuracy, smoothness regime) visible in the abstract and theorem statement, not buried in Section 3, and exhibit a non-toy instance where it provably holds.
  \item \textit{[reject\_reviews]} When substituting a sharper object into an existing proof, explicitly bound the new error term the substitution introduces (sketch vs. true SVD, factored vs. full second moment, retraction-free iterate vs. on-manifold iterate) --- reviewers comb the proof for self-adjointness, tightness of constants, and missing approximation accounting.
\end{itemize}
\textbf{Cognitive barriers:}
\begin{itemize}[leftmargin=12pt,topsep=2pt,itemsep=1pt]
  \item Researchers treat the bottleneck step in a celebrated proof as a fundamental requirement of the result rather than as an artifact of one particular proof technique --- once a bound has been cited for years, the field stops asking whether it is tight, whether the assumption is load-bearing, or whether a worst-case proxy could be replaced by algorithm-specific behavior.
  \item Existing 'tagline' results (e.g., 'noise defeats all convex boosters', 'second-order methods need exact Hessians per step', 'shuffling needs i.i.d.-style analysis') are accepted as universally stated, even when the original proof implicitly relied on a specific subclass; spotting that the universality is conditional requires a re-reading of an old proof against its stated theorem.
  \item The right sharper substitute often lives in an adjacent mathematical subfield --- contour integration, Poisson processes, communication complexity, Hilbert-space Freedman, abstract convexity --- that the methodology's home community does not routinely consult, so the obstacle is not technical but knowing where to look.
  \item Convergence-style analysis and resource/efficiency analysis (sample, query, communication, memory, information) are conflated, so the community fails to notice that an algorithm whose convergence is well-understood is suboptimal in a different but observable resource.
\end{itemize}
\textbf{Representative examples:}
\begin{itemize}[leftmargin=12pt,topsep=2pt,itemsep=1pt]
  \item \textbf{[Oral]} \texttt{NeurIPS\_2023\_0715} \textemdash\ Replaces the discrete union bound that produced a log k slack in random-cuts analysis by recasting the randomness as a Poisson process where independence is exact, extracting the precise constant 2 ln k + 2.
  \\ \textcolor{gray}{\footnotesize \textit{Lesson:} Replace a discrete union bound with a setting where independence is exact (Poisson process) --- the sharpening extracts the precise constant, not just an order-of-magnitude improvement.}
  \item \textbf{[Oral]} \texttt{NeurIPS\_2025\_1890} \textemdash\ Identifies that GP-UCB regret was loose only because prior analysis bounded worst-case MIG, and substitutes a query-sequence-specific MIG bound directly into the existing analysis without changing the algorithm.
  \\ \textcolor{gray}{\footnotesize \textit{Lesson:} The algorithm doesn't change; what changes is the analysis --- substituting a query-sequence-specific MIG bound directly into the existing analysis is a clean sharpening move.}
  \item \textbf{[Oral]} \texttt{NeurIPS\_2025\_1932} \textemdash\ Bypasses the non-smoothness of the spectral norm by using contour integration around eigenvalues --- a tool from operator theory rarely seen in numerical-linear-algebra proofs --- to obtain perturbation bounds scaling with the eigengap rather than global matrix norm.
  \\ \textcolor{gray}{\footnotesize \textit{Lesson:} Borrow tools from operator theory rarely seen in numerical-linear-algebra proofs to bypass non-smoothness --- perturbation bounds scale with eigengap, not global matrix norm.}
  \item \textbf{[Oral]} \texttt{NeurIPS\_2024\_1306} \textemdash\ Closes the average-reward MDP rate gap by replacing the diameter (a worst-case geometric proxy) with the span (a tighter functional parameter) and reproves all complexity components against this new parameter, then matches with a lower bound.
  \\ \textcolor{gray}{\footnotesize \textit{Lesson:} Replace a worst-case geometric proxy (diameter) with a tighter functional parameter (span); reprove all complexity components and match with a lower bound.}
  \item \textbf{[Reject]} \texttt{ICLR\_2023\_0162} \textemdash\ Attempts adaptive-regret guarantees for adaptive optimizers, but the AC notes the main theorem follows as a corollary of an unsurveyed prior paper --- the 'sharpened' bound was already implicit in the literature.
  \\ \textcolor{gray}{\footnotesize \textit{Lesson:} If the main theorem is a corollary of an unsurveyed prior paper, the sharpening was already implicit --- literature surveys are mandatory for refinement papers.}
  \item \textbf{[Reject]} \texttt{ICLR\_2023\_0318} \textemdash\ Relaxes Lipschitz to coordinate-wise smoothness for Adam, but the proof's constants explode as $\beta$$_2$$\to$1 and the convergence is only to a noise floor, so the headline 'provable adaptivity' does not actually cascade.
  \\ \textcolor{gray}{\footnotesize \textit{Lesson:} A 'sharpened' bound whose constants explode at the regime boundary or that converges only to a noise floor doesn't actually deliver provable adaptivity.}
  \item \textbf{[Reject]} \texttt{ICLR\_2024\_0834} \textemdash\ Replaces full SVD with an incremental sketch in a second-order online kernel method, but the regret proof neglects the tracking error between the maintained sketch and the true low-rank approximation, leaving the sharper bound unjustified.
  \\ \textcolor{gray}{\footnotesize \textit{Lesson:} Neglecting the tracking error between the maintained sketch and the true low-rank approximation leaves the sharper bound unjustified --- refinement must propagate cleanly through ALL terms.}
\end{itemize}
\end{tcolorbox}
